% Lesson 60 — Writing: Writes compositions on difficulties with the use of ICTs at work or for learning — Anglais Tle A — Cours signé OnBuch+
% Programme officiel MINESEC. Compiler avec : tectonic EN60-writing-writes-compositions-on-difficulties-with-the-us.tex
\newcommand{\DOCMATIERE}{Anglais}
\newcommand{\DOCNIVEAU}{Tle A}
\newcommand{\DOCMODULE}{Chapter 5 : Media and Communication}
\newcommand{\DOCLECON}{EN60}
\newcommand{\DOCTITRE}{Lesson 60 — Writing: Writes compositions on difficulties with the use of ICTs at work or for learning}
% OnBuch+ — préambule « cours signés » ENRICHI (compilé avec tectonic / XeLaTeX)
% Système visuel « Pop » d'OnBuch+ : fond crème (jamais blanc), bordures encre
% épaisses, ombres « dures » décalées, titres Archivo Black en capitales.
% Version MATHÉMATIQUES : graphes (pgfplots/tikz), boîtes pédagogiques
% (définition, propriété, méthode, attention, le savais-tu, exercice résolu,
% exercice, TP, bilan), page de garde et signature OnBuch+ ancrée en bas.
%
% Macros attendues dans main.tex AVANT \input{preamble} :
%   \DOCMATIERE  \DOCNIVEAU  \DOCMODULE  \DOCLECON (numéro)  \DOCTITRE
\documentclass[11pt]{article}

\usepackage{fontspec}
\usepackage[english,french]{babel} % français = langue principale ; l'anglais via \eng{..} / textebox
\usepackage[a4paper,margin=2cm,top=3.4cm,bottom=2.8cm,headheight=1.4cm,footskip=1.3cm]{geometry}
\usepackage{amsmath,amssymb,amsfonts}
\usepackage{mathtools} % \xmapsto, \coloneqq... souvent utilisés par les modèles
\usepackage{cancel} % simplifications barrées (bilans de forces)
% \abs{z} (module, valeur absolue) et \norm{u} (norme), utilisés par les modèles
\providecommand{\abs}{}\let\abs\relax\DeclarePairedDelimiter{\abs}{\lvert}{\rvert}
\providecommand{\norm}{}\let\norm\relax\DeclarePairedDelimiter{\norm}{\lVert}{\rVert}
\usepackage[table]{xcolor} % [table] : fournit aussi \rowcolors (alternance de lignes)
\usepackage{pagecolor}
\usepackage{tikz}
\usetikzlibrary{arrows.meta,positioning,calc,shapes.geometric,shapes.misc,shapes.arrows,decorations.pathmorphing,decorations.pathreplacing,decorations.markings,patterns,shadows,mindmap,trees,fit,backgrounds,matrix,intersections,angles,quotes}
\usepackage{pgfplots}
\usepackage[version=4]{mhchem} % équations nucléaires \ce{^{235}_{92}U}
\usepackage[european,siunitx]{circuitikz} % schémas électriques
\pgfplotsset{compat=1.18}
\usepgfplotslibrary{fillbetween,groupplots} % aires entre courbes (calcul intégral)
\usepackage{enumitem}
\usepackage{fancyhdr}
% [most] charge toutes les bibliothèques tcolorbox (listings, minted, raster,
% documentation...) et gonfle inutilement la mémoire TeX sur un document long
% (~300 boîtes) : on ne charge que ce qui est réellement utilisé.
\usepackage[skins,breakable]{tcolorbox}
\usepackage{graphicx}
\usepackage{colortbl}
\usepackage{booktabs}
\usepackage{tabularx}
\usepackage{diagbox} % case d'en-tête diagonale des tableaux à double entrée
% Colonnes extensibles centrée / gauche / droite (C, L, R), souvent utilisées par les modèles
\newcolumntype{C}{>{\centering\arraybackslash}X}
\newcolumntype{L}{>{\raggedright\arraybackslash}X}
\newcolumntype{R}{>{\raggedleft\arraybackslash}X}
\usepackage{array}
\usepackage{multicol}
\usepackage{multirow}
\usepackage{siunitx}
% parse-numbers=false : accepte aussi \SI{2,5 \times 10^{-3}}{...}
\sisetup{locale=FR,output-decimal-marker={,},group-minimum-digits=4,parse-numbers=false}
\DeclareSIUnit\ohm{\ensuremath{\symup{\Omega}}} % U+2126 absent de la police
\DeclareSIUnit\division{div} % réglages d'oscilloscope (ms/div, V/div)
\DeclareSIUnit\tr{tr} % tours (vitesse de rotation, ex. \SI{1500}{\tr\per\minute})
\DeclareSIUnit\molar{M} % concentration molaire (ex. \SI{3}{\molar}, \SI{1}{\micro\molar})
\DeclareSIUnit\an{an} % année (le modèle confond parfois avec \year, un compteur TeX) — ex. \SI{5}{\kilo\meter\per\an}
\DeclareSIUnit\FCFA{FCFA} % franc CFA (ex. \SI{500}{\FCFA\per\kilo\gram})
\DeclareSIUnit\degreeBrix{\textdegree Brix} % teneur en sucre d'un sirop/jus (réfractométrie)
\DeclareSIUnit\month{mois} % le modèle utilise parfois \month, un compteur TeX, comme unité de durée
\usepackage{qrcode}
% \degree : commande fréquemment utilisée par les modèles pour un angle ou une
% température en dehors de \SI{}{} (non fournie par les paquets chargés).
\providecommand{\degree}{^{\circ}}
% \qed (fin de démonstration), utilisé par les modèles sans amsthm
\providecommand{\qed}{\hfill$\square$}
% \barre : double barre des valeurs interdites dans un tableau de variations (tkz-tab)
\providecommand{\barre}{\ensuremath{\|}}
% environnement proof (amsthm non chargé) utilisé par les modèles
\ifdefined\proof\else\newenvironment{proof}[1][Démonstration]{\par\noindent\textit{#1.}\ }{\qed\par}\fi
\usepackage{lastpage}
\usepackage{hyperref}
\hypersetup{hidelinks}

% ============================================================
% Palette « Pop » OnBuch+ — mêmes hex que la classe P (plus_ui.dart)
% ============================================================
\definecolor{poporange}{HTML}{F59321}
\definecolor{popdark}{HTML}{E07A0C}
\definecolor{popcream}{HTML}{FFF6E8}
\definecolor{popink}{HTML}{1C1714}
\definecolor{poppurple}{HTML}{7A5AE0}
\definecolor{popgreen}{HTML}{2E9E6A}
\definecolor{poppink}{HTML}{E0457B}
\definecolor{popblue}{HTML}{2F7DE1}
\definecolor{popgold}{HTML}{C98A12}
\definecolor{popmuted}{HTML}{8A7F76}
\definecolor{popline}{HTML}{EADFD2}
\definecolor{popgray}{HTML}{8A7F76}
\colorlet{popgrey}{popgray}
% Alias tolérants (le modèle nomme parfois les couleurs différemment)
\colorlet{popwhite}{white}
\colorlet{popblack}{popink}
\colorlet{popred}{poppink}
\colorlet{popyellow}{popgold}
% Teintes claires (fonds de tableaux/encadrés) — le modèle nomme parfois
% les couleurs avec un suffixe L (ex. popblueL) sans les avoir définies.
\colorlet{poporangeL}{poporange!20}
\colorlet{popdarkL}{popdark!20}
\colorlet{poppurpleL}{poppurple!20}
\colorlet{popgreenL}{popgreen!20}
\colorlet{poppinkL}{poppink!20}
\colorlet{popblueL}{popblue!20}
\colorlet{popgoldL}{popgold!20}
\colorlet{popredL}{popred!20}
\colorlet{popgrayL}{popgray!20}
% Teintes claires pour fonds de tableaux / zones
\colorlet{popblueL}{popblue!10!white}
\colorlet{poporangeL}{poporange!14!white}
\colorlet{popgreenL}{popgreen!12!white}
\colorlet{poppurpleL}{poppurple!12!white}
\colorlet{poppinkL}{poppink!12!white}
\colorlet{popgoldL}{popgold!14!white}

\pagecolor{popcream}

% ============================================================
% Typographie
% ============================================================
\defaultfontfeatures{Ligatures=TeX}
\setmainfont{PlusJakartaSans-Medium.ttf}[Path=../../../fonts/,BoldFont=PlusJakartaSans-Bold.ttf]
% Maths en sans-serif (Fira Math) avec chiffres et lettres droites de Plus
% Jakarta, pour que les formules se fondent dans le texte.
\usepackage[math-style=ISO,bold-style=ISO]{unicode-math}
\setmathfont{FiraMath-Regular.otf}[Path=../../../fonts/]
\setmathfont{PlusJakartaSans-Medium.ttf}[Path=../../../fonts/,range={up/num,up/{Latin,latin},\mathup}]
\usepackage{icomma} % virgule décimale sans espace en mode math
% Caractères mathématiques Unicode absents des polices de texte (Plus Jakarta,
% Archivo Black) : rendus par leur équivalent mathématique, en texte comme en
% formule — sinon ils s'affichent en carré vide (ex. « Arithmétique dans ℤ »).
\usepackage{newunicodechar}
\newunicodechar{ℝ}{\ensuremath{\mathbb{R}}}
\newunicodechar{ℤ}{\ensuremath{\mathbb{Z}}}
\newunicodechar{ℂ}{\ensuremath{\mathbb{C}}}
\newunicodechar{ℕ}{\ensuremath{\mathbb{N}}}
\newunicodechar{ℚ}{\ensuremath{\mathbb{Q}}}
\newunicodechar{𝔻}{\ensuremath{\mathbb{D}}}
\newunicodechar{α}{\ensuremath{\alpha}}
\newunicodechar{β}{\ensuremath{\beta}}
\newunicodechar{γ}{\ensuremath{\gamma}}
\newunicodechar{δ}{\ensuremath{\delta}}
\newunicodechar{ε}{\ensuremath{\varepsilon}}
\newunicodechar{θ}{\ensuremath{\theta}}
\newunicodechar{λ}{\ensuremath{\lambda}}
\newunicodechar{μ}{\ensuremath{\mu}}
\newunicodechar{σ}{\ensuremath{\sigma}}
\newunicodechar{τ}{\ensuremath{\tau}}
\newunicodechar{φ}{\ensuremath{\varphi}}
\newunicodechar{ψ}{\ensuremath{\psi}}
\newunicodechar{ω}{\ensuremath{\omega}}
\newunicodechar{Σ}{\ensuremath{\Sigma}}
% majuscules produites par \MakeUppercase dans les titres (α-aminés -> Α)
\newunicodechar{Α}{\ensuremath{\alpha}}
\newunicodechar{Β}{\ensuremath{\beta}}
\newunicodechar{Γ}{\ensuremath{\Gamma}}
\newunicodechar{Δ}{\ensuremath{\Delta}}
\newunicodechar{Ε}{\ensuremath{\varepsilon}}
\newunicodechar{Θ}{\ensuremath{\Theta}}
\newunicodechar{Λ}{\ensuremath{\Lambda}}
\newunicodechar{Μ}{\ensuremath{\mu}}
\newunicodechar{Τ}{\ensuremath{\tau}}
\newunicodechar{Φ}{\ensuremath{\Phi}}
\newunicodechar{Ψ}{\ensuremath{\Psi}}
\newunicodechar{Ω}{\ensuremath{\Omega}}
\newunicodechar{↦}{\ensuremath{\mapsto}}
\newunicodechar{∈}{\ensuremath{\in}}
\newunicodechar{∉}{\ensuremath{\notin}}
\newunicodechar{∧}{\ensuremath{\wedge}}
\newunicodechar{⇔}{\ensuremath{\Leftrightarrow}}
\newunicodechar{⟺}{\ensuremath{\Longleftrightarrow}}
\newunicodechar{⇒}{\ensuremath{\Rightarrow}}
\newunicodechar{⟹}{\ensuremath{\Longrightarrow}}
\newunicodechar{∘}{\ensuremath{\circ}}
\newunicodechar{∪}{\ensuremath{\cup}}
\newunicodechar{∩}{\ensuremath{\cap}}
\newunicodechar{≡}{\ensuremath{\equiv}}
\newunicodechar{⊂}{\ensuremath{\subset}}
\newunicodechar{⊥}{\ensuremath{\perp}}
\newunicodechar{∃}{\ensuremath{\exists}}
\newunicodechar{∀}{\ensuremath{\forall}}
\newunicodechar{∛}{\ensuremath{\sqrt[3]{\,}}}
\newunicodechar{∜}{\ensuremath{\sqrt[4]{\,}}}
\newunicodechar{⃗}{\ensuremath{{}^{\to}}}
\newunicodechar{ⁿ}{\ifmmode^{n}\else\textsuperscript{\ensuremath{n}}\fi}
\newunicodechar{ˣ}{\ifmmode^{x}\else\textsuperscript{\ensuremath{x}}\fi}
\newunicodechar{ᵇ}{\ifmmode^{b}\else\textsuperscript{\ensuremath{b}}\fi}
\newunicodechar{ʳ}{\ifmmode^{r}\else\textsuperscript{\ensuremath{r}}\fi}
\newunicodechar{ᵘ}{\ifmmode^{u}\else\textsuperscript{\ensuremath{u}}\fi}
\newunicodechar{ʸ}{\ifmmode^{y}\else\textsuperscript{\ensuremath{y}}\fi}
\newunicodechar{ᵃ}{\ifmmode^{a}\else\textsuperscript{\ensuremath{a}}\fi}
\newunicodechar{ᵉ}{\ifmmode^{e}\else\textsuperscript{\ensuremath{e}}\fi}
\newunicodechar{ᵏ}{\ifmmode^{k}\else\textsuperscript{\ensuremath{k}}\fi}
\newunicodechar{ᵖ}{\ifmmode^{p}\else\textsuperscript{\ensuremath{p}}\fi}
\newunicodechar{ᴺ}{\ifmmode^{N}\else\textsuperscript{\ensuremath{N}}\fi}
\newunicodechar{ᶜ}{\ifmmode^{c}\else\textsuperscript{\ensuremath{c}}\fi}
\newunicodechar{ʰ}{\ifmmode^{h}\else\textsuperscript{\ensuremath{h}}\fi}
\newunicodechar{ᵠ}{\ifmmode^{q}\else\textsuperscript{\ensuremath{q}}\fi}
\newunicodechar{ₙ}{\ifmmode_{n}\else\textsubscript{\ensuremath{n}}\fi}
\newunicodechar{ₐ}{\ifmmode_{a}\else\textsubscript{\ensuremath{a}}\fi}
\newunicodechar{ᵢ}{\ifmmode_{i}\else\textsubscript{\ensuremath{i}}\fi}
\newunicodechar{ᵣ}{\ifmmode_{r}\else\textsubscript{\ensuremath{r}}\fi}
\newunicodechar{ₖ}{\ifmmode_{k}\else\textsubscript{\ensuremath{k}}\fi}
\newunicodechar{ᵦ}{\ifmmode_{\beta}\else\textsubscript{\ensuremath{\beta}}\fi}
\newunicodechar{ₓ}{\ifmmode_{x}\else\textsubscript{\ensuremath{x}}\fi}
\newfontfamily\popdisplay{ArchivoBlack-Regular.ttf}[Path=../../../fonts/]
\newfontfamily\poplogo{SpaceGrotesk-Bold.ttf}[Path=../../../fonts/]
\newfontfamily\popsemi{PlusJakartaSans-SemiBold.ttf}[Path=../../../fonts/]
\newfontfamily\popxbold{PlusJakartaSans-ExtraBold.ttf}[Path=../../../fonts/]
\newfontfamily\popxboldls{PlusJakartaSans-ExtraBold.ttf}[Path=../../../fonts/,LetterSpace=3.0]

\setlist[enumerate]{leftmargin=1.7em,itemsep=5pt,topsep=4pt}
\setlist[itemize]{leftmargin=1.7em,itemsep=4pt,topsep=4pt,label={\color{poporange}\textbullet}}
\setlength{\parindent}{0pt}
\setlength{\parskip}{5pt}
\renewcommand{\arraystretch}{1.25}

% pgfplots : style Pop par défaut
\pgfplotsset{
  every axis/.append style={axis line style={popink,line width=1pt},tick style={popink},
    grid style={popline,line width=0.5pt},label style={font=\small},tick label style={font=\footnotesize}},
  popcurve/.style={poporange,line width=1.8pt,no markers,smooth},
}
% Même style disponible dans un \draw TikZ simple (hors pgfplots)
\tikzset{popcurve/.style={poporange,line width=1.8pt}}
\tikzset{popgrid/.style={popline,very thin}}
\colorlet{popcurve}{poporange} % le nom du style est parfois utilisé comme couleur

% ============================================================
% Pill / squiggle
% ============================================================
\newcommand{\poppill}[3]{%
  \tikz[baseline=-0.55ex]\node[draw=popink,line width=1.2pt,rounded corners=2.6pt,%
    fill=#2,inner xsep=6.5pt,inner ysep=3pt]{%
    \popxboldls\fontsize{7}{7}\selectfont\color{#3}\MakeUppercase{#1}};%
}
\newcommand{\popsquiggle}[1]{%
  \tikz[baseline=-0.4ex]{
    \draw[#1,line width=2.6pt,line cap=round]
      plot [smooth] coordinates {(0,0.09) (0.34,0.01) (0.68,0.21) (1.02,0.01) (1.36,0.10)};
  }%
}
% Mot-clé surligné (marqueur orange sous le texte)
\newcommand{\cle}[1]{{\popxbold\color{popdark}#1}}

% ============================================================
% En-tête / pied
% ============================================================
\pagestyle{fancy}
\fancyhf{}
\renewcommand{\headrulewidth}{0pt}
\renewcommand{\footrulewidth}{0pt}
\lhead{\raisebox{-0.15ex}{\poplogo\fontsize{13}{13}\selectfont\color{popink}OnBuch\color{poporange}+}}
\rhead{\poppill{\DOCMATIERE\ \textbullet\ \DOCNIVEAU\ \textbullet\ Leçon \DOCLECON}{white}{popink}}
\lfoot{\popxbold\fontsize{7.6}{7.6}\selectfont\color{popmuted}\MakeUppercase{Cours signé OnBuch+}\ \textbullet\ {\popsemi\DOCTITRE}}
\rfoot{\tikz[baseline=-0.6ex]\node[rounded corners=8.5pt,fill=popink,minimum height=17pt,minimum width=17pt,inner xsep=5pt,inner sep=0]{\popxbold\fontsize{8}{8}\selectfont\color{popcream}\thepage\,/\,\pageref*{LastPage}};}

% ============================================================
% Titres
% ============================================================
\newcommand{\chaptitre}[2]{%
  \begin{tcolorbox}[enhanced,colback=poporange,colframe=popink,boxrule=2.6pt,arc=11pt,
      drop shadow=popink,left=15pt,right=15pt,top=12pt,bottom=11pt,before skip=3pt,after skip=18pt]
    \poppill{#2}{popink}{popcream}\par\vspace{9pt}
    {\raggedright\hyphenpenalty=10000\exhyphenpenalty=10000\popdisplay\fontsize{22}{24}\selectfont\color{popcream}\MakeUppercase{#1}\par}
  \end{tcolorbox}
}
% Section numérotée : pastille orange avec le numéro + titre Archivo
\newcounter{popsec}
\newcommand{\coursec}[1]{%
  \stepcounter{popsec}\setcounter{popsub}{0}%
  \par\vspace{16pt}\noindent
  \begin{minipage}{\linewidth}
  \tikz[baseline=-0.7ex]\node[draw=popink,line width=1.6pt,fill=poporange,rounded corners=4pt,minimum size=22pt,inner sep=0]{\popdisplay\fontsize{12}{12}\selectfont\color{popcream}\thepopsec};%
  \hspace{8pt}{\popdisplay\fontsize{15}{17}\selectfont\color{popink}\MakeUppercase{#1}}\par
  \vspace{4pt}\noindent{\color{poporange}\rule{36pt}{3.6pt}}
  \end{minipage}\par\nopagebreak\vspace{8pt}
}
\newcounter{popsub}
\newcommand{\courssub}[1]{%
  \stepcounter{popsub}%
  \par\vspace{9pt}\noindent{\popxbold\fontsize{11.5}{13}\selectfont\color{popink}{\color{poporange}\thepopsec.\thepopsub}\hspace{6pt}#1}\par\nopagebreak\vspace{4pt}
}

% ============================================================
% Boîtes pédagogiques — même recette : fond blanc, bordure épaisse colorée,
% ombre dure encre, badge plein. Titre optionnel [#1].
% ============================================================
\tcbset{popbase/.style={breakable,enhanced,colback=white,boxrule=2.3pt,arc=9pt,
  shadow={3pt}{-3pt}{0pt}{popink},left=14pt,right=14pt,top=11pt,bottom=11pt,before skip=11pt,after skip=15pt,
  fontupper=\popsemi\fontsize{10.6}{14.5}\selectfont\color{popink},
  fontlower=\popsemi\fontsize{10.6}{14.5}\selectfont\color{popink}}}
% Coche dessinée (le glyphe ✓ n'existe pas dans les polices du document)
\newcommand{\popcheck}[1]{\tikz[baseline=-0.5ex]\draw[#1,line width=1.6pt,line cap=round,line join=round](-0.13,0.0)--(-0.04,-0.09)--(0.13,0.1);}
\providecommand{\coche}{\popcheck{popgreen}}
\providecommand{\resultat}[1]{\par\smallskip\noindent\fcolorbox{popgreen}{popgreenL}{\parbox{\dimexpr\linewidth-2\fboxsep-2\fboxrule\relax}{\textbf{Résultat.} #1}}\par\smallskip}
\newcommand{\popboxhead}[3]{\poppill{#1}{#2}{white}\ifx\relax#3\relax\else\hspace{6pt}{\popxbold\fontsize{10.6}{12}\selectfont\color{popink}#3}\fi\par\vspace{7pt}}

\newtcolorbox{aretenir}[1][]{popbase,colframe=popgreen,before upper={\popboxhead{À retenir}{popgreen}{#1}}}
% Texte en anglais (document de lecture, dialogue, extrait) : cadre
% d'étude. Un titre optionnel (source, auteur) s'affiche dans l'en-tête.
\newtcolorbox{textebox}[1][]{popbase,colframe=popblue,before upper={\popboxhead{Text}{popblue}{#1}\begin{otherlanguage}{english}},after upper={\end{otherlanguage}}}
% Marques ✓ / ✗ (forme correcte / fautive) souvent utilisées par les modèles
\providecommand{\cross}{\textcolor{poppink}{\ensuremath{\times}}}
\providecommand{\xmark}{\cross}\providecommand{\crossmark}{\cross}\providecommand{\cmark}{\ensuremath{\checkmark}}
% Ligne de réponse / justification à compléter (inventée par les modèles)
\providecommand{\justification}[1]{\par\noindent\textit{Justification~:} \dotfill\par}
% \textipa (package tipa non chargé) : la police ne couvre pas l'API -> simple passage
\providecommand{\textipa}[1]{#1}
% Soulignement (ulem non chargé) : \uline, \ul
\providecommand{\uline}[1]{\underline{#1}}\providecommand{\ul}[1]{\underline{#1}}
% \glossaire{\item ...} inventée par les modèles : liste de glossaire
\providecommand{\glossaire}[1]{\begin{itemize}#1\end{itemize}}
% \alert (beamer) inventée par les modèles : texte mis en évidence
\providecommand{\alert}[1]{\textbf{\textcolor{poppink}{#1}}}
% variantes ASCII d'ordinaux écrites en macro
\providecommand{\iere}{\textsuperscript{re}}\providecommand{\ieme}{\textsuperscript{e}}\providecommand{\ere}{\textsuperscript{re}}\providecommand{\eme}{\textsuperscript{e}}\providecommand{\er}{\textsuperscript{er}}\providecommand{\ie}{\textsuperscript{e}}
% Ordinaux français écrits en macro par les modèles : 1\ière, 3\ième
\newcommand{\ière}{\textsuperscript{re}}\newcommand{\ième}{\textsuperscript{e}}
% \exemple (inventée par les modèles) : étiquette « Exemple : »
\providecommand{\exemple}{\textbf{Exemple~:}\ }
% \square utilisable aussi hors mode math (cases à cocher)
\AtBeginDocument{\let\squareMath\square\renewcommand{\square}{\ensuremath{\squareMath}}}
% Mot ou phrase en anglais à l'intérieur d'un paragraphe français
\newcommand{\eng}[1]{\ifmmode\text{\textit{#1}}\else\begingroup\selectlanguage{english}\itshape #1\endgroup\fi}
\let\esp\eng % alias toléré
% flèches d'intonation inventées par les modèles
\providecommand{\textsearrow}{}\renewcommand{\textsearrow}{\ensuremath{\searrow}}\providecommand{\textnearrow}{}\renewcommand{\textnearrow}{\ensuremath{\nearrow}}
% \fr{..} : mot français (inventé par les modèles)
\providecommand{\fr}[1]{\textit{#1}}
\newtcolorbox{exemplebox}[1][]{popbase,colframe=poporange,before upper={\popboxhead{Exemple}{poporange}{#1}}}
\newtcolorbox{definition}[1][]{popbase,colframe=popblue,before upper={\popboxhead{Définition}{popblue}{#1}}}
\newtcolorbox{propriete}[1][]{popbase,colframe=poppurple,before upper={\popboxhead{Propriété}{poppurple}{#1}}}
\newtcolorbox{methode}[1][]{popbase,colframe=popgold,before upper={\popboxhead{Méthode}{popgold}{#1}}}
\newtcolorbox{attention}[1][]{popbase,colframe=poppink,before upper={\popboxhead{Attention}{poppink}{#1}}}
\newtcolorbox{experience}[1][]{popbase,colframe=popblue,colback=popblueL,before upper={\popboxhead{Expérience}{popblue}{#1}}}
% « Le savais-tu ? » : carte violette pleine (contexte, culture, Cameroun)
\newtcolorbox{savaistu}[1][]{popbase,colback=poppurple,colframe=popink,
  fontupper=\popsemi\fontsize{10.4}{14.2}\selectfont\color{white},
  before upper={\poppill{Le savais-tu\,?}{white}{poppurple}\ifx\relax#1\relax\else\hspace{6pt}{\popxbold\color{white}#1}\fi\par\vspace{7pt}}}
% Exercice résolu : énoncé en haut, \tcblower puis la solution sur fond crème
\newcounter{popexo}\newcounter{popexr}
\newtcolorbox{exoresolu}[1][]{popbase,colframe=poporange,colbacklower=popcream,
  segmentation style={popink,line width=1.2pt,dashed},
  before upper={\stepcounter{popexr}\popboxhead{Exercice résolu \thepopexr}{poporange}{#1}},
  before lower={\poppill{Solution}{popink}{popcream}\par\vspace{6pt}}}
% Exercice (non résolu en place) — la correction est dans la partie Corrigés
\newtcolorbox{exercice}[1][]{popbase,colframe=popink,
  before upper={\stepcounter{popexo}\popboxhead{Exercice \thepopexo}{popink}{#1}}}
% Corrigé
\newtcolorbox{corrige}[1][]{popbase,colframe=popgreen,colback=popgreenL,
  before upper={\popboxhead{Corrigé}{popgreen}{#1}}}
% Objectifs / prérequis / bilan
\newtcolorbox{objectifs}{popbase,colframe=popink,colback=poporangeL,
  before upper={\popboxhead{Objectifs de la leçon}{popink}{}}}
\newtcolorbox{prerequis}{popbase,colframe=popink,colback=popblueL,
  before upper={\popboxhead{Prérequis}{popblue}{}}}
\newtcolorbox{bilan}{popbase,colframe=popink,colback=white,boxrule=2.8pt,
  before upper={\popboxhead{Fiche bilan}{poporange}{L'essentiel à connaître pour l'examen}}}
% Figure encadrée avec légende
\newtcolorbox{popfigure}[1][]{enhanced,breakable=false,colback=white,colframe=popline,boxrule=1.4pt,arc=8pt,
  left=10pt,right=10pt,top=10pt,bottom=8pt,before skip=10pt,after skip=12pt,halign=center,
  lower separated=false,halign lower=center,
  fontlower=\popsemi\fontsize{9}{11.5}\selectfont\color{popmuted},
  #1}
\newcommand{\legende}[1]{\par\vspace{6pt}{\popsemi\fontsize{9}{11.5}\selectfont\color{popmuted}#1}}

% ============================================================
% Signature OnBuch+
% ============================================================
\newcommand{\poplogoinline}[1]{{\poplogo\fontsize{#1}{#1}\selectfont\color{popink}OnBuch\color{poporange}+}}
\newcommand{\signatureonbuch}{%
  \begin{tcolorbox}[enhanced,colback=popink,colframe=popink,boxrule=2.2pt,arc=10pt,
      drop shadow=poporange,left=14pt,right=14pt,top=10pt,bottom=10pt,before skip=0pt,after skip=0pt]
    \tikz[baseline=-0.7ex]\node[circle,fill=popgreen,draw=popcream,line width=1.3pt,minimum size=22pt,inner sep=0]{\popcheck{white}};%
    \hspace{9pt}\begin{minipage}[c]{0.62\linewidth}
      {\popxbold\fontsize{12}{13}\selectfont\color{popcream}Signé\ \textbullet\ {\poplogo OnBuch\color{poporange}+}}\par\vspace{2pt}
      {\raggedright\popsemi\fontsize{8}{10.5}\selectfont\color{popline}\MakeUppercase{Équipe pédagogique OnBuch+}\\\MakeUppercase{Conforme au programme officiel MINESEC}\par}
    \end{minipage}\hfill
    {\poplogo\fontsize{22}{22}\selectfont\color{popcream}OnBuch\color{poporange}+}
  \end{tcolorbox}%
}

% ============================================================
% Page de garde
% #1 titre · #2 matière · #3 niveau · #4 module · #5 n° leçon · #6 durée ·
% #7 description / objectifs (texte ou itemize)
% ============================================================
\newcommand{\pagegarde}[7]{%
  \thispagestyle{empty}
  \newgeometry{a4paper,margin=1.7cm,top=1.6cm,bottom=1.4cm}%
  \begin{tikzpicture}[remember picture,overlay]
    \fill[poporange!18] ([xshift=-3.2cm,yshift=-3.6cm]current page.north east) circle (4.6cm);
    \fill[poppurple!14] ([xshift=2.4cm,yshift=7.6cm]current page.south west) circle (3.8cm);
  \end{tikzpicture}%
  {\poplogo\fontsize{30}{30}\selectfont\color{popink}OnBuch\color{poporange}+}\hfill
  \raisebox{0.6ex}{\poppill{Cours signé \textbullet\ Premium}{popink}{popcream}}\par
  \vspace{1pt}\popsquiggle{poppurple}\par
  \vspace{8pt}
  \begin{tcolorbox}[enhanced,colback=poporange,colframe=popink,boxrule=2.8pt,arc=13pt,
      drop shadow=popink,left=18pt,right=18pt,top=12pt,bottom=13pt,after skip=10pt]
    \poppill{#2 \textbullet\ #3}{popink}{popcream}\hspace{5pt}\poppill{#4}{white}{popink}\par\vspace{8pt}
    {\popdisplay\fontsize{14}{14}\selectfont\color{popink}LEÇON #5}\par\vspace{4pt}
    {\raggedright\hyphenpenalty=10000\exhyphenpenalty=10000\popdisplay\fontsize{25}{27}\selectfont\color{popcream}\MakeUppercase{#1}\par}
    \vspace{8pt}
    {\popxbold\fontsize{9.5}{9.5}\selectfont\color{popink}\MakeUppercase{Durée conseillée : #6}}
  \end{tcolorbox}
  \begin{tcolorbox}[enhanced,colback=white,colframe=popink,boxrule=2.2pt,arc=10pt,
      drop shadow=popink,left=16pt,right=16pt,top=10pt,bottom=10pt,after skip=8pt]
    \poppill{Dans cette leçon}{poppurple}{white}\par\vspace{5pt}
    \popsemi\fontsize{9.3}{12.6}\selectfont\color{popink}#7
  \end{tcolorbox}
  \noindent\begin{minipage}[t]{0.487\linewidth}
    \begin{tcolorbox}[enhanced,colback=popgreen,colframe=popink,boxrule=2.2pt,arc=10pt,
        drop shadow=popink,left=11pt,right=11pt,top=10pt,bottom=10pt,height=3.1cm,valign=center]
      \tcbox[colback=white,colframe=popink,boxrule=1.3pt,arc=5pt,boxsep=0pt,nobeforeafter,
        left=3pt,right=3pt,top=3pt,bottom=3pt,valign=center]{\qrcode[height=1.9cm]{https://play.google.com/store/apps/details?id=com.luvvix.onbuch&hl=fr}}%
      \hspace{8pt}\begin{minipage}[c]{0.5\linewidth}
        {\popdisplay\fontsize{11}{12.5}\selectfont\color{white}\MakeUppercase{Télécharge OnBuch}}\par\vspace{4pt}
        {\raggedright\popsemi\fontsize{8.4}{11}\selectfont\color{white}Gratuit \textbullet\ Google Play\\Tuteur IA, annales, concours, résultats\par}
      \end{minipage}
    \end{tcolorbox}
  \end{minipage}\hfill
  \begin{minipage}[t]{0.487\linewidth}
    \begin{tcolorbox}[enhanced,colback=poppurple,colframe=popink,boxrule=2.2pt,arc=10pt,
        drop shadow=popink,left=11pt,right=11pt,top=10pt,bottom=10pt,height=3.1cm,valign=center]
      \tikz[baseline=-0.7ex]\node[circle,fill=white,draw=popink,line width=1.3pt,minimum size=1.6cm,inner sep=0]{\popdisplay\fontsize{24}{24}\selectfont\color{poppurple}+};%
      \hspace{8pt}\begin{minipage}[c]{0.6\linewidth}
        {\popdisplay\fontsize{11}{12.5}\selectfont\color{white}\MakeUppercase{Passe à OnBuch+}}\par\vspace{4pt}
        {\raggedright\popsemi\fontsize{8.4}{11}\selectfont\color{white}Tous les cours signés, Tuteur IA illimité, fiches PDF — onglet OnBuch+ de l'app\par}
      \end{minipage}
    \end{tcolorbox}
  \end{minipage}
  \vspace{8pt}
  \popcontenu
  \vfill
  \signatureonbuch
  \restoregeometry
  \newpage
}

% Rangée « Ce cours contient » (page de garde)
\newcommand{\poptile}[3]{% #1 couleur #2 gros texte #3 légende
  \begin{tcolorbox}[enhanced,colback=#1,colframe=popink,boxrule=2pt,arc=9pt,shadow={2.5pt}{-2.5pt}{0pt}{popink},
      left=6pt,right=6pt,top=9pt,bottom=9pt,halign=center,valign=center,height=1.75cm,width=\linewidth,nobeforeafter]
    {\popdisplay\fontsize{12.5}{13}\selectfont\color{white}\MakeUppercase{#2}}\par\vspace{3pt}
    {\centering\popsemi\fontsize{7.8}{9.5}\selectfont\color{white}#3\par}
  \end{tcolorbox}}
\newcommand{\popcontenu}{%
  \par\noindent{\popxboldls\fontsize{7.5}{7.5}\selectfont\color{popmuted}CE COURS SIGNÉ CONTIENT}\par\vspace{7pt}
  \noindent\begin{minipage}[t]{0.235\linewidth}\poptile{popblue}{Cours}{complet, illustré, pas à pas}\end{minipage}\hfill
  \begin{minipage}[t]{0.235\linewidth}\poptile{poporange}{Méthodes}{et exercices résolus}\end{minipage}\hfill
  \begin{minipage}[t]{0.235\linewidth}\poptile{poppink}{Exercices}{type examen + corrigés}\end{minipage}\hfill
  \begin{minipage}[t]{0.235\linewidth}\poptile{popgreen}{Fiche bilan}{l'essentiel à retenir}\end{minipage}\par}

% Sommaire
\newtcolorbox{sommaire}{enhanced,colback=white,colframe=popink,boxrule=2.2pt,arc=10pt,
  drop shadow=popink,left=18pt,right=18pt,top=13pt,bottom=13pt,before skip=6pt,after skip=20pt,
  before upper={\poppill{Sommaire}{poppurple}{white}\par\vspace{9pt}\popsemi\fontsize{10.6}{14.5}\selectfont\color{popink}}}

% Signature de fin de cours — poussée tout en bas de la dernière page
\newcommand{\findecours}{%
  \par\vfill\nopagebreak
  \begin{center}\popsquiggle{poporange}\end{center}\vspace{4pt}
  \signatureonbuch
}
\AtBeginDocument{\def\llbracket{\mathopen{[\mkern-3.5mu[}}\def\rrbracket{\mathclose{]\mkern-3.5mu]}}} % intervalles d'entiers (glyphe ⟦ absent de la police)
% Glyphes absents de Fira Math : redéfinis à partir de glyphes disponibles
\AtBeginDocument{%
  \def\setminus{\mathbin{\backslash}}\let\smallsetminus\setminus
  \def\vdots{\mathord{\vbox{\baselineskip3.2pt\lineskiplimit0pt\kern4pt\hbox{.}\hbox{.}\hbox{.}}}}%
  \def\ddots{\mathinner{\mkern1mu\raise6.5pt\hbox{.}\mkern2mu\raise3.6pt\hbox{.}\mkern2mu\raise0.7pt\hbox{.}\mkern1mu}}%
  \def\triangle{\mathord{\scalebox{0.9}{$\symup{\Delta}$}}}%
  \def\nabla{\mathord{\raisebox{\depth}{\scalebox{1}[-1]{$\symup{\Delta}$}}}}%
  \def\checkmark{\popcheck{popdark}}%
  \def\star{\mathbin{\ast}}\def\bigstar{\mathord{\ast}}%
  \def\diamond{\mathbin{\scalebox{0.6}{$\Diamond$}}}%
  \def\bigcup{\mathop{\mathchoice{\raisebox{-0.3em}{\scalebox{1.7}{$\cup$}}}{\scalebox{1.25}{$\cup$}}{\cup}{\cup}}\displaylimits}%
  \def\bigcap{\mathop{\mathchoice{\raisebox{-0.3em}{\scalebox{1.7}{$\cap$}}}{\scalebox{1.25}{$\cap$}}{\cap}{\cap}}\displaylimits}%
  \def\lesseqqgtr{\mathrel{\lessgtr}}\def\bigoplus{\mathop{\oplus}\displaylimits}%
}
\providecommand{\ssi}{si et seulement si} % abréviation fréquente
\AtBeginDocument{\providecommand{\wideparen}{\overparen}} % arc de cercle

\begin{document}

\pagegarde{Lesson 60 — Writing: Writes compositions on difficulties with the use of ICTs at work or for learning}{Anglais}{Tle A}{Chapter 5 : Media and Communication}{EN60}{2 h}{Cette leçon d'expression écrite apprend aux élèves de Terminale A à planifier, rédiger et corriger une composition sur les difficultés liées à l'usage des TIC au travail ou dans les apprentissages. Elle fournit la structure attendue au Bac, les connecteurs, le vocabulaire utile, un modèle commenté et une grille d'évaluation.
\vspace{4pt}
\begin{multicols}{2}\small
\begin{itemize}
\item À la fin de cette leçon, je sais identifier le plan en trois parties (introduction, développement, conclusion) d'une composition argumentative.
\item À la fin de cette leçon, je sais formuler une problématique (thesis statement) claire sur les difficultés des TIC.
\item À la fin de cette leçon, je sais utiliser le vocabulaire thématique des TIC et de leurs problèmes (breakdown, power cut, distraction, cybercrime...).
\item À la fin de cette leçon, je sais organiser mes idées avec des connecteurs d'addition, de cause, de conséquence et de concession.
\item À la fin de cette leçon, je sais employer correctement les temps verbaux adaptés (present simple, present perfect, modaux).
\item À la fin de cette leçon, je sais rédiger une composition de 250 à 300 mots sur les difficultés des TIC au travail ou à l'école.
\end{itemize}\end{multicols}}

\begin{sommaire}
\begin{multicols}{2}
\begin{enumerate}
\item Understanding the task and brainstorming
\item The structure of a composition
\item Useful vocabulary and connectors
\item Verb tenses and grammar for this essay
\item A commented model essay
\item Common mistakes and assessment grid
\item Activité d'intégration
\item Méthodes et exercices résolus
\item Exercices
\item Corrigés des exercices
\item Fiche bilan
\end{enumerate}
\end{multicols}
\end{sommaire}

\begin{objectifs}
\begin{itemize}
\item À la fin de cette leçon, je sais identifier le plan en trois parties (introduction, développement, conclusion) d'une composition argumentative.
\item À la fin de cette leçon, je sais formuler une problématique (thesis statement) claire sur les difficultés des TIC.
\item À la fin de cette leçon, je sais utiliser le vocabulaire thématique des TIC et de leurs problèmes (breakdown, power cut, distraction, cybercrime...).
\item À la fin de cette leçon, je sais organiser mes idées avec des connecteurs d'addition, de cause, de conséquence et de concession.
\item À la fin de cette leçon, je sais employer correctement les temps verbaux adaptés (present simple, present perfect, modaux).
\item À la fin de cette leçon, je sais rédiger une composition de 250 à 300 mots sur les difficultés des TIC au travail ou à l'école.
\item À la fin de cette leçon, je sais relire et corriger ma production à l'aide d'une grille d'évaluation.
\item À la fin de cette leçon, je sais éviter les erreurs fréquentes des francophones (faux amis, calques du français).
\end{itemize}
\end{objectifs}

\begin{prerequis}
\begin{itemize}
\item Le présent simple et le present perfect (leçons de Première).
\item Les connecteurs de base : and, but, because, so, however.
\item Le vocabulaire général des médias et des TIC vu dans le chapitre 5 (computer, Internet, network...).
\item La structure d'un paragraphe : phrase d'introduction, arguments, exemple.
\item Les modaux can, must, should pour exprimer capacité, obligation et conseil.
\end{itemize}
\end{prerequis}

\newpage

\coursec{Understanding the task and brainstorming}

\courssub{Hook: when technology lets you down}

During the last mock exam at GBHS Bamenda, the computer room froze in the middle of the e-learning session: no connection, no printer, and half the class had never saved their work. Sound familiar? From cybercafés in Yaoundé to universities in London or Lagos, ICTs are everywhere — but so are their problems. Today, you will learn to write a well-organised composition describing these difficulties, a classic topic at the Baccalauréat.

La problématique de notre leçon sera donc la suivante : \eng{Although ICTs make work and learning easier, they also create serious difficulties.} (« Bien que les TIC facilitent le travail et l'apprentissage, elles créent aussi de sérieuses difficultés. »). Avant d'écrire quoi que ce soit, il faut d'abord \cle{comprendre le sujet} et \cle{faire jaillir des idées} (brainstorming). C'est l'objet de cette première section.

\courssub{Reading a model text first}

Avant d'analyser le sujet, lisons un court texte qui illustre le thème. Observez les difficultés mentionnées : nous les réutiliserons dans notre brainstorming.

\begin{textebox}[A connected school, a disconnected day]
At 7.30 a.m., Mrs Nfor, a biology teacher in Buea, switches on the computer room. She has prepared a brilliant online quiz for her Form Five students. But at 8.15, the lights go off: another power cut. The generator is broken, and the school cannot afford a new one. When electricity returns two hours later, the Internet connection is so slow that the quiz page will not load. Three students try to download the document on their phones, but their data bundles are finished, and buying more data is expensive for their families. Meanwhile, two boys at the back are not working at all: they are watching football videos on social media. Mrs Nfor sighs. She has never received any real training in digital tools; she learned everything alone, watching tutorials at night. Last month, a colleague even lost his salary to a cybercriminal who sent him a fake bank message. At the end of the day, Mrs Nfor's eyes hurt after hours in front of the screen. She loves technology — she knows it can open the world to her students — but some days, she admits, a piece of chalk and a blackboard seem more reliable. Her story is not unusual. In many schools and offices, the promise of ICTs collides daily with technical, financial and human obstacles. Understanding these difficulties is the first step towards solving them — and, for you, the first step towards writing an excellent exam composition.
\end{textebox}

\textbf{Glossaire}\par
\begin{center}
\begin{tabularx}{\linewidth}{|X|X|X|}\hline
\rowcolor{popblueL} Mot anglais & Nature & Traduction \\ \hline
power cut & nom & coupure d'électricité, délestage \\ \hline
generator & nom & groupe électrogène \\ \hline
to afford & verbe & avoir les moyens d'acheter \\ \hline
data bundle & nom & forfait internet, forfait de données \\ \hline
training & nom & formation \\ \hline
fake & adjectif & faux, falsifié \\ \hline
reliable & adjectif & fiable \\ \hline
to collide with & verbe & se heurter à \\ \hline
\end{tabularx}
\end{center}

\textbf{Questions de compréhension}\par
\begin{enumerate}
\item What technical problems does Mrs Nfor face during the day?
\item Why can't the students download the document on their phones?
\item What human difficulties are mentioned in the text?
\end{enumerate}

\courssub{Analysing the exam topic}

Voici le sujet type que vous pourrez rencontrer au Baccalauréat :

\begin{exemplebox}[Sujet type Bac]
\eng{Write a composition of about 250--300 words on the difficulties people face when using ICTs at work or for learning.}

« Rédigez une composition d'environ 250 à 300 mots sur les difficultés que les gens rencontrent lorsqu'ils utilisent les TIC au travail ou pour apprendre. »
\end{exemplebox}

Un bon candidat ne commence jamais à écrire immédiatement. Il \cle{analyse le sujet} en soulignant les mots-clés. Dans ce sujet, quatre éléments sont essentiels :

\begin{itemize}
\item \eng{difficulties} : le sujet demande des \textbf{problèmes}, pas les avantages des TIC. Hors sujet si vous ne parlez que des bienfaits !
\item \eng{ICTs} : le champ lexical est celui des technologies de l'information et de la communication (ordinateurs, Internet, téléphones, logiciels).
\item \eng{at work} : un premier contexte obligatoire (bureaux, entreprises, administrations).
\item \eng{for learning} : un deuxième contexte obligatoire (école, université, formation).
\end{itemize}

\begin{methode}[Analyser un sujet de composition en trois étapes]
\begin{enumerate}
\item \textbf{Underline the key words.} Soulignez les mots porteurs de sens : ici \eng{difficulties}, \eng{ICTs}, \eng{at work}, \eng{for learning}. Chaque mot-clé doit apparaître dans votre copie.
\item \textbf{Identify the type of text.} Demande-t-on une lettre, un article, un discours, une composition ? Ici, c'est une \cle{composition} (essai argumenté descriptif) : pas d'adresse, pas de formule de politesse, mais une introduction, un développement et une conclusion.
\item \textbf{Respect the word limit.} \eng{About 250--300 words} signifie entre 240 et 320 mots environ. Comptez vos mots à la fin : une copie de 120 mots ou de 600 mots sera pénalisée.
\end{enumerate}
\end{methode}

\begin{attention}[Ne recopiez jamais le sujet comme introduction]
Beaucoup de candidats commencent leur copie par : \eng{I am going to write a composition of about 250--300 words on the difficulties people face when using ICTs at work or for learning.} C'est une erreur grave : le correcteur voit immédiatement que vous n'avez aucune idée. Il faut \textbf{reformuler le sujet avec vos propres mots}. Comparez :

\eng{ICTs have invaded our offices and classrooms, yet using them is not always easy.} (bonne reformulation)

\eng{The difficulties people face when using ICTs at work or for learning are many.} (simple paraphrase du sujet : faible)
\end{attention}

\courssub{Brainstorming: generating ideas}

Le \cle{brainstorming} (remue-méninges) consiste à noter \textbf{toutes} les idées qui viennent, sans les juger, puis à trier. Pour ce sujet, travaillez en deux colonnes correspondant aux deux contextes du sujet.

\begin{center}
\begin{tabularx}{\linewidth}{|X|X|}\hline
\rowcolor{popblueL} Difficulties at work & Difficulties for learning \\ \hline
power cuts stop the whole office & students cannot follow e-learning during blackouts \\ \hline
slow or poor Internet connection & online lessons and videos will not load \\ \hline
high cost of data and devices & many students cannot afford a laptop \\ \hline
lack of training for workers & teachers not trained in digital tools \\ \hline
computer breakdowns, lost files & unsaved work disappears before exams \\ \hline
distractions: social media at the office & students watch videos instead of studying \\ \hline
cybercrime: fake messages, stolen money & cheating and false information online \\ \hline
health problems: eye strain, back pain & eye strain after long hours on screens \\ \hline
\end{tabularx}
\end{center}

\begin{exemplebox}[Idées formulées en anglais correct]
\eng{Many students cannot afford a laptop.} « Beaucoup d'élèves n'ont pas les moyens d'acheter un ordinateur portable. »

\eng{A single power cut can stop a whole day's work.} « Une seule coupure d'électricité peut arrêter toute une journée de travail. »

\eng{Some workers have never received any training in digital tools.} « Certains employés n'ont jamais reçu aucune formation aux outils numériques. »

\eng{Social media distract students from their studies.} « Les réseaux sociaux distraient les élèves de leurs études. »

\eng{Staring at a screen for hours causes eye strain.} « Fixer un écran pendant des heures provoque une fatigue oculaire. »
\end{exemplebox}

Pour organiser ces idées, on peut les classer en quatre familles, comme le montre la carte mentale suivante.

\begin{popfigure}
\begin{tikzpicture}[mindmap, grow cyclic, every node/.style={concept, text=white}, root concept/.append style={concept color=popblue, font=\bfseries}, level 1/.append style={level distance=3.4cm, sibling angle=90}, level 2/.append style={level distance=2.2cm, sibling angle=45, concept color=popmuted, font=\small}]
\node[root concept]{ICT difficulties}
  child[concept color=poporange]{ node {technical} child { node {power cuts} } child { node {breakdowns} } }
  child[concept color=popgreen]{ node {financial} child { node {cost of data} } child { node {devices} } }
  child[concept color=poppurple]{ node {human} child { node {lack of training} } child { node {distraction} } }
  child[concept color=poppink]{ node {social} child { node {cybercrime} } child { node {isolation} } };
\end{tikzpicture}
\legende{Carte mentale : les quatre familles de difficultés liées aux TIC.}
\end{popfigure}

\begin{savaistu}[Le délestage, star des copies du Bac]
Au Cameroun, les coupures d'électricité — les fameux \eng{power cuts} ou « délestages » — sont l'une des difficultés les plus citées dans les copies du Baccalauréat sur ce thème. Les correcteurs les connaissent bien ! Pour vous distinguer, citez aussi des difficultés moins attendues : le coût des données (\eng{the high cost of data}), la fatigue visuelle (\eng{eye strain}) ou la cybercriminalité (\eng{cybercrime}). En anglais britannique, on dit aussi \eng{power failure} ou \eng{blackout}.
\end{savaistu}

\courssub{Selecting and ranking the best ideas}

Vous ne pouvez pas tout dire en 250--300 mots. Il faut \cle{sélectionner} trois ou quatre idées fortes et les \cle{classer par ordre d'importance}. Un bon plan pour ce sujet :

\begin{enumerate}
\item \textbf{Technical difficulties} (power cuts, poor connection, breakdowns) — les plus visibles.
\item \textbf{Financial difficulties} (cost of data and devices) — elles touchent surtout les élèves.
\item \textbf{Human difficulties} (lack of training, distractions, health problems) — les plus profondes.
\item \textbf{Social difficulties} (cybercrime, isolation) — pour conclure le développement en ouvrant vers la société.
\end{enumerate}

Ce classement du plus concret au plus large donne une progression naturelle à votre composition. La problématique qui guidera l'ensemble est : \eng{Although ICTs make work and learning easier, they also create serious difficulties.}

\begin{exoresolu}[Repérer les mots-clés et classer les idées]
Énoncé. Voici un sujet : \eng{Write a composition of about 250--300 words on the difficulties people face when using ICTs at work or for learning.} (a) Relevez les quatre mots-clés du sujet. (b) Parmi les idées suivantes, choisissez-en quatre et classez-les par ordre d'importance pour un développement : \eng{cybercrime, power cuts, the weather, lack of training, high cost of data, football results.}
\tcblower
Solution. (a) Les mots-clés sont : \eng{difficulties}, \eng{ICTs}, \eng{at work}, \eng{for learning}. (b) \eng{The weather} et \eng{football results} sont hors sujet : on les élimine. Classement possible : 1. \eng{power cuts} (difficulté technique majeure, vécue par tous) ; 2. \eng{high cost of data} (obstacle financier pour les élèves et les employés) ; 3. \eng{lack of training} (difficulté humaine : on ne peut pas utiliser ce qu'on ne comprend pas) ; 4. \eng{cybercrime} (danger social, idéal pour clore le développement). Ce plan couvre bien les deux contextes du sujet : le travail et l'apprentissage.
\end{exoresolu}

\begin{exercice}[À vous de jouer]
(a) Soulignez les mots-clés dans ce sujet voisin : \eng{Write a composition of about 250 words on the advantages and disadvantages of mobile phones in schools.} (b) Faites un brainstorming en deux colonnes (advantages / disadvantages) avec au moins quatre idées par colonne, puis sélectionnez vos trois meilleures idées et justifiez votre classement en une phrase.
\end{exercice}

\begin{corrige}[Exercice : À vous de jouer]
(a) Mots-clés : \eng{advantages and disadvantages}, \eng{mobile phones}, \eng{in schools}. Attention : ce sujet exige les DEUX faces, contrairement au sujet sur les difficultés. (b) Exemple de réponse : advantages — quick research, contacting parents, educational apps, calculator ; disadvantages — distraction in class, cheating in exams, cyberbullying, cost of data. Classement possible : distraction (la plus fréquente), cheating (la plus grave pour l'école), educational apps (le meilleur argument positif pour équilibrer la conclusion).
\end{corrige}

\begin{aretenir}[L'essentiel de la section]
\begin{itemize}
\item Analysez le sujet avant d'écrire : soulignez les mots-clés, identifiez le type de texte, respectez le nombre de mots.
\item Ne recopiez jamais le sujet en introduction : reformulez-le.
\item Faites un brainstorming en deux colonnes (\eng{at work} / \eng{for learning}), puis sélectionnez 3 ou 4 idées et classez-les du plus concret au plus large.
\item Votre composition entière doit répondre à la problématique : \eng{Although ICTs make work and learning easier, they also create serious difficulties.}
\end{itemize}
\end{aretenir}

\coursec{The structure of a composition}

Dans la section précédente, nous avons compris le sujet et récolté des idées grâce au brainstorming : coupures de courant, connexion lente, coût des équipements, manque de formation, distractions, conséquences sur les notes et le travail. Nous avons maintenant une \cle{liste d'idées}, mais une liste n'est pas une composition. Il faut organiser ces idées dans une \cle{structure claire} que le correcteur reconnaîtra immédiatement. C'est l'objet de cette section : apprendre à construire le squelette de votre composition \emph{avant} d'écrire la première phrase.

\courssub{Overview: the three parts of a composition}

Une composition au baccalauréat, comme un article de presse anglophone, se divise toujours en \cle{trois parties} :

\begin{itemize}
\item \textbf{\eng{The introduction}} (3--4 phrases) : elle présente le sujet et annonce votre point de vue.
\item \textbf{\eng{The body}} (2--3 paragraphes) : elle développe vos arguments, une idée par paragraphe.
\item \textbf{\eng{The conclusion}} (3--4 phrases) : elle résume et donne votre opinion ou une suggestion.
\end{itemize}

\begin{popfigure}
\begin{tikzpicture}[
  bloc/.style={draw=popdark, thick, rounded corners=4pt, align=center, text width=11.5cm, inner sep=7pt},
  fleche/.style={-{Stealth[length=3mm]}, very thick, poporange}
]
\node[bloc, fill=popblueL, align=center] (intro) at (0,0){\textbf{\textcolor{popblue}{INTRODUCTION}}\\ Hook (accroche) + general presentation + thesis statement};
\node[bloc, fill=popgreenL, below=0.55cm of intro, align=center] (body){\textbf{\textcolor{popgreen}{BODY}}\\ Paragraph 1: technical difficulties (power cuts, poor connection)\\ Paragraph 2: human and financial difficulties (cost, training, distraction)\\ Paragraph 3 (optional): consequences on work and school results};
\node[bloc, fill=poppinkL, below=0.55cm of body, align=center] (concl){\textbf{\textcolor{poppink}{CONCLUSION}}\\ Summary of the main ideas + personal opinion or suggested solutions};
\draw[fleche] (intro) -- (body);
\draw[fleche] (body) -- (concl);
\end{tikzpicture}
\legende{La structure en trois blocs d'une composition : chaque flèche représente la progression logique de votre texte.}
\end{popfigure}

\begin{aretenir}[Target length]
Au baccalauréat, visez \cle{250 à 300 mots}, soit environ \cle{25 à 30 lignes}. Une composition de 150 mots est trop pauvre ; une composition de 500 mots contient presque toujours des répétitions et des hors-sujets. Répartition conseillée : introduction $\approx$ 40 mots, chaque paragraphe du corps $\approx$ 60--70 mots, conclusion $\approx$ 40 mots.
\end{aretenir}

\courssub{The introduction: from general to particular}

Observez d'abord cette introduction réussie, écrite par un élève de Terminale :

\begin{exemplebox}[A model introduction]
\eng{Nowadays, information and communication technologies are present in every office and every classroom in Cameroon. Many people believe that computers and the internet make work and learning easier. However, using ICTs also creates serious difficulties. This essay will examine the main technical, human and financial problems that workers and students face with ICTs.}

« Aujourd'hui, les technologies de l'information et de la communication sont présentes dans chaque bureau et chaque salle de classe au Cameroun. Beaucoup de gens pensent que les ordinateurs et Internet facilitent le travail et l'apprentissage. Cependant, l'utilisation des TIC crée aussi de sérieuses difficultés. Cette composition examinera les principaux problèmes techniques, humains et financiers auxquels les travailleurs et les élèves font face avec les TIC. »
\end{exemplebox}

Remarquez le mouvement : la première phrase est très \cle{générale} (les TIC partout), puis le texte se \cle{rétrécit} progressivement jusqu'à la dernière phrase, qui annonce exactement le plan. C'est le principe de l'\emph{entonnoir} (\eng{funnel}) : large au début, précis à la fin.

\begin{definition}[Thesis statement]
La \eng{thesis statement} (problématique ou phrase de thèse) est la phrase, généralement placée \emph{à la fin de l'introduction}, qui annonce le point de vue défendu ou le plan de la composition. Elle répond à la question : « What is this essay going to prove or discuss? ». Sans \eng{thesis statement}, le lecteur ne sait pas où vous allez.
\end{definition}

Une introduction complète contient donc trois éléments :

\begin{enumerate}
\item \textbf{\eng{The hook}} (l'accroche) : une première phrase qui attire l'attention — un constat général, une question, un contraste. \eng{Nowadays, almost every student in Yaoundé owns a smartphone.} « De nos jours, presque chaque élève de Yaoundé possède un smartphone. »
\item \textbf{\eng{The general presentation}} : une ou deux phrases qui situent le thème et montrent les deux faces du sujet. \eng{Many people believe that ICTs make learning easier. However, they also create serious difficulties.} « Beaucoup pensent que les TIC facilitent l'apprentissage. Cependant, elles créent aussi de sérieuses difficultés. »
\item \textbf{\eng{The thesis statement}} : la phrase qui annonce le plan. \eng{This essay will examine the main technical, human and financial problems faced by ICT users.} « Cette composition examinera les principaux problèmes techniques, humains et financiers rencontrés par les utilisateurs des TIC. »
\end{enumerate}

\begin{methode}[How to write your introduction in four steps]
\begin{enumerate}
\item Commencez \textbf{large} : une phrase générale sur les TIC dans la société ou à l'école (\eng{Nowadays, ICTs are everywhere...}).
\item Ajoutez une phrase sur l'\textbf{opinion courante} (\eng{Many people think that...}).
\item Introduisez le \textbf{contraste} avec \eng{however} ou \eng{nevertheless} : le sujet comporte des difficultés.
\item Terminez par la \eng{thesis statement} qui annonce vos deux ou trois paragraphes (\eng{This essay will discuss...}).
\end{enumerate}
Vérifiez : si l'on lit uniquement votre dernière phrase, on doit pouvoir deviner le plan de votre développement.
\end{methode}

\courssub{The body: one paragraph, one idea}

Le développement (\eng{body}) est le cœur de la composition. Sa loi fondamentale est simple :

\begin{propriete}[The paragraph rule]
Un paragraphe = \cle{une seule idée principale}. Cette idée est annoncée dans la \cle{première phrase} du paragraphe, appelée \eng{topic sentence}. Le reste du paragraphe ne fait que trois choses : \textbf{expliquer} l'idée, la \textbf{illustrer} par un exemple concret, éventuellement la \textbf{nuancer}. Toute phrase qui ne sert pas l'idée principale doit être supprimée ou déplacée dans un autre paragraphe.
\end{propriete}

Structure interne d'un bon paragraphe :

\begin{center}
\begin{tabularx}{\linewidth}{|l|X|X|}
\hline
\rowcolor{popblueL} \textbf{Élément} & \textbf{\eng{English example}} & \textbf{Traduction} \\
\hline
Topic sentence & \eng{First of all, power cuts are a major problem.} & Tout d'abord, les coupures de courant sont un problème majeur. \\
\hline
Explanation & \eng{When electricity goes off, computers stop working and unsaved documents are lost.} & Quand l'électricité se coupe, les ordinateurs s'arrêtent et les documents non sauvegardés sont perdus. \\
\hline
Concrete example & \eng{For instance, during examinations in Douala, many students cannot print their documents because of blackouts.} & Par exemple, pendant les examens à Douala, beaucoup d'élèves ne peuvent pas imprimer leurs documents à cause des coupures. \\
\hline
Link (optionnel) & \eng{Moreover, a poor internet connection makes the situation worse.} & De plus, une mauvaise connexion internet aggrave la situation. \\
\hline
\end{tabularx}
\end{center}

\begin{exemplebox}[More topic sentences for this essay]
\eng{In addition, the high cost of ICT equipment prevents many families from buying computers.} « De plus, le coût élevé du matériel informatique empêche de nombreuses familles d'acheter des ordinateurs. »

\eng{Another serious difficulty is the lack of training.} « Une autre difficulté sérieuse est le manque de formation. »

\eng{As a result, these difficulties have negative consequences on school results.} « Par conséquent, ces difficultés ont des conséquences négatives sur les résultats scolaires. »
\end{exemplebox}

Pour notre sujet, voici le plan de développement recommandé :

\begin{itemize}
\item \textbf{Paragraph 1 — Technical difficulties} : \eng{power cuts} (coupures de courant), \eng{poor internet connection} (connexion lente), \eng{breakdowns} (pannes), \eng{outdated computers} (ordinateurs vétustes).
\item \textbf{Paragraph 2 — Human and financial difficulties} : \eng{the high cost of equipment and data} (coût du matériel et des données), \eng{lack of training} (manque de formation), \eng{distractions} comme les réseaux sociaux et les jeux.
\item \textbf{Paragraph 3 (optional) — Consequences} : \eng{wasted time, poor marks, stress, inequality between rich and poor students} (temps perdu, mauvaises notes, stress, inégalités entre élèves riches et pauvres).
\end{itemize}

\begin{attention}[One idea, one paragraph]
Erreur classique des candidats francophones : mélanger dans le même paragraphe les coupures de courant, le prix des smartphones et les mauvaises notes. Le correcteur anglophone sanctionne immédiatement ce manque de \eng{paragraphing}. Dès que vous changez d'idée, \textbf{allez à la ligne} et commencez par une nouvelle \eng{topic sentence} introduite par un connecteur (\eng{First of all}, \eng{Moreover}, \eng{In addition}, \eng{Another difficulty is}, \eng{As a result}).
\end{attention}

\courssub{The conclusion: summary and opinion}

La conclusion est le miroir de l'introduction : elle repart du particulier vers le général. Elle contient deux éléments :

\begin{enumerate}
\item \textbf{\eng{A summary}} des idées principales, en d'autres mots (\eng{In conclusion}, \eng{To sum up}, \eng{All in all}).
\item \textbf{\eng{A personal opinion or a suggestion}} : votre avis ou une solution proposée (\eng{In my opinion}, \eng{I believe that}, \eng{The government should...}).
\end{enumerate}

\begin{exemplebox}[A model conclusion]
\eng{To sum up, ICT users in Cameroon face technical problems such as power cuts, as well as human and financial difficulties like the high cost of equipment and the lack of training. In my opinion, the government and schools should invest in solar energy and offer free computer courses, so that technology can truly help workers and learners.}

« Pour résumer, les utilisateurs des TIC au Cameroun font face à des problèmes techniques comme les coupures de courant, ainsi qu'à des difficultés humaines et financières comme le coût élevé du matériel et le manque de formation. À mon avis, le gouvernement et les écoles devraient investir dans l'énergie solaire et offrir des cours d'informatique gratuits, afin que la technologie aide réellement les travailleurs et les apprenants. »
\end{exemplebox}

\begin{attention}[Never introduce a new idea in the conclusion]
Ne jamais introduire une \cle{idée nouvelle} dans la conclusion. Si vous écrivez \eng{Moreover, ICTs can also spread fake news} dans votre dernier paragraphe, vous ouvrez un sujet que vous ne développerez jamais : le correcteur le sanctionne. La conclusion \textbf{résume} et \textbf{conclut}, elle n'\textbf{ajoute} rien. De même, évitez de recopier mot pour mot votre introduction : reformulez (\eng{paraphrase}).
\end{attention}

\courssub{French and English essays: a useful comparison}

La tradition de la dissertation française (thèse, antithèse, synthèse) diffère de la composition anglaise, plus directe. Ce tableau vous évite les transferts négatifs :

\begin{center}
\begin{tabularx}{\linewidth}{|X|X|}
\hline
\rowcolor{popblueL} \textbf{Dissertation française} & \textbf{\eng{English composition}} \\
\hline
Plan dialectique long (thèse / antithèse / synthèse) & Plan direct : thesis statement claire dès l'introduction \\
\hline
Introduction en trois parties très codifiées (amorce, présentation, problématique, annonce du plan) & Hook + general presentation + thesis statement (3--4 phrases seulement) \\
\hline
Conclusion avec ouverture vers un autre sujet & Conclusion fermée : résumé + opinion, aucune idée nouvelle \\
\hline
Connecteurs implicites tolérés & Connecteurs explicites obligatoires (\eng{first, moreover, as a result}) \\
\hline
Le point de vue personnel souvent caché & L'opinion personnelle attendue (\eng{In my opinion...}) \\
\hline
\end{tabularx}
\end{center}

\begin{exoresolu}[Rebuild the structure]
Voici les phrases d'une composition mélangées. Remettez-les dans l'ordre et indiquez à quelle partie (\eng{introduction}, \eng{body} 1, \eng{body} 2, \eng{conclusion}) chacune appartient :

a) \eng{As a result, students who depend on computers often get poor marks.}

b) \eng{Nowadays, ICTs are used in almost every school and office in Cameroon.}

c) \eng{In my opinion, schools should provide free computer training to all learners.}

d) \eng{First of all, frequent power cuts make it difficult to work with computers.}

e) \eng{This essay will discuss the main difficulties of using ICTs at work and for learning.}

f) \eng{Moreover, many people cannot afford expensive laptops and internet data.}

\tcblower
\textbf{Solution détaillée.}

Ordre correct : \textbf{b -- e -- d -- f -- a -- c}.

\begin{itemize}
\item \textbf{\eng{Introduction}} : b (hook général sur les TIC) puis e (thesis statement qui annonce le plan).
\item \textbf{\eng{Body}, paragraph 1} : d (topic sentence sur les difficultés techniques, signalée par \eng{First of all}).
\item \textbf{\eng{Body}, paragraph 2} : f (difficultés financières, signalées par \eng{Moreover}).
\item \textbf{\eng{Body}, paragraph 3} : a (conséquences sur les résultats, signalées par \eng{As a result}).
\item \textbf{\eng{Conclusion}} : c (opinion personnelle et suggestion, signalées par \eng{In my opinion}).
\end{itemize}

On reconnaît chaque partie grâce aux \cle{connecteurs} et au contenu : le \eng{hook} est général, les \eng{topic sentences} introduisent une idée unique, la \eng{conclusion} donne un avis sans idée nouvelle.
\end{exoresolu}

\begin{exercice}[Your turn]
Sans rédiger encore la composition complète, écrivez uniquement :
\begin{enumerate}
\item une phrase d'accroche (\eng{hook}) sur les TIC dans votre ville ou votre lycée ;
\item une \eng{thesis statement} annonçant deux ou trois types de difficultés ;
\item trois \eng{topic sentences} pour vos trois paragraphes de développement ;
\item une phrase de conclusion avec votre opinion personnelle.
\end{enumerate}
Vous obtiendrez ainsi le squelette complet de votre composition : il ne restera plus, dans les sections suivantes, qu'à développer chaque phrase.
\end{exercice}

\begin{corrige}[Exercice : éléments de réponse]
Réponses possibles (plusieurs variantes sont acceptables) :

1. Hook : \eng{In Bamenda, almost every student now uses a smartphone to do homework.}

2. Thesis statement : \eng{However, using ICTs for learning creates technical, financial and human difficulties that this essay will examine.}

3. Topic sentences : \eng{First of all, a poor internet connection is a daily problem for learners.} / \eng{In addition, the cost of data and devices is too high for many families.} / \eng{Finally, these difficulties seriously affect students' marks and motivation.}

4. Conclusion : \eng{To sum up, although ICTs are useful, they remain difficult to use; I believe that cheaper internet and better training would help every learner.}
\end{corrige}

\begin{aretenir}[The golden rules of structure]
\begin{itemize}
\item Trois parties : \eng{introduction} (3--4 phrases), \eng{body} (2--3 paragraphes), \eng{conclusion} (3--4 phrases).
\item L'\eng{introduction} va du général au particulier et se termine par la \eng{thesis statement}.
\item Un paragraphe = une idée, annoncée par une \eng{topic sentence}, puis expliquée et illustrée.
\item La \eng{conclusion} résume et donne une opinion, sans aucune idée nouvelle.
\item Longueur cible : 250--300 mots.
\end{itemize}
\end{aretenir}

Maintenant que le squelette est en place, la section suivante vous fournira la chair : le vocabulaire précis et les connecteurs logiques pour remplir chaque bloc de cette structure.

\coursec{Useful vocabulary and connectors}

Dans la section précédente, nous avons appris qu'une composition se construit en trois temps : une \cle{introduction} qui présente le sujet, un \cle{development} qui organise les arguments en paragraphes, et une \cle{conclusion} qui synthétise. Mais une bonne architecture ne suffit pas : il faut encore disposer des \emph{mots justes} pour parler des difficultés liées aux TIC, et des \emph{charnières} qui relient les idées entre elles. C'est exactement l'objet de cette section : enrichir votre vocabulaire thématique, puis maîtriser les \cle{connectors} (\eng{link words}) qui donnent du rythme et de la logique à votre texte.

\courssub{Observing the language in context}

Avant toute liste de vocabulaire, lisons un court texte authentique. Soulignez mentalement les mots qui désignent des problèmes et les mots qui relient les idées.

\begin{textebox}[Reading: ICTs at school — a blessing or a burden?]
In many Cameroonian schools, computers and smartphones have changed the way students learn. On the one hand, ICTs save time: a student in Buea can download past examination papers in a few seconds. On the other hand, they can be a source of distraction and frustration.

First of all, technical problems are frequent. Power cuts often interrupt online lessons, and a slow connection makes video conferences almost impossible. Moreover, many computers are out of order because schools cannot afford maintenance costs. As a result, students sometimes lose their documents when a machine crashes or when a virus attacks the system.

Furthermore, human difficulties should not be forgotten. Because of a lack of training, some teachers feel uncomfortable with digital tools, and computer illiteracy remains common among parents. In addition, many teenagers suffer from eye strain and addiction to social media, which reduces their concentration in class.

Finally, there are serious risks such as hacking, scams and fake news. Therefore, although ICTs are useful, schools must train users and protect their privacy. All in all, technology helps learning only when it is well managed.
\end{textebox}

\textbf{Glossaire}

\begin{center}
\begin{tabularx}{\linewidth}{|l|X|}\hline
\rowcolor{popblueL} \textbf{English} & \textbf{Français} \\ \hline
power cut & coupure de courant \\ \hline
out of order & hors service, en panne \\ \hline
to afford & avoir les moyens (financiers) de \\ \hline
to crash & planter (ordinateur, système) \\ \hline
lack of training & manque de formation \\ \hline
eye strain & fatigue visuelle \\ \hline
scam & arnaque \\ \hline
privacy & vie privée, confidentialité \\ \hline
\end{tabularx}
\end{center}

\textbf{Questions de compréhension}
\begin{enumerate}
\item What technical problems are mentioned in the text?
\item Why do students sometimes lose their documents?
\item What human difficulties are described?
\item What solution does the writer suggest at the end?
\end{enumerate}

\courssub{The vocabulary of ICT difficulties}

\begin{definition}[Champ lexical]
Un \cle{champ lexical} est l'ensemble des mots qui se rapportent à un même thème. Pour rédiger sur les difficultés des TIC, il faut mobiliser quatre champs lexicaux : les problèmes techniques, les coûts, les difficultés humaines et les risques.
\end{definition}

\textbf{1. Technical problems (les problèmes techniques)}

\begin{center}
\begin{tabularx}{\linewidth}{|X|X|X|}\hline
\rowcolor{popblueL} \textbf{English} & \textbf{Français} & \textbf{Exemple} \\ \hline
power cut & coupure de courant & \eng{A power cut stopped the online exam.} \\ \hline
breakdown & panne & \eng{The network breakdown lasted two days.} \\ \hline
poor / slow connection & connexion mauvaise / lente & \eng{Because of a slow connection, the video froze.} \\ \hline
virus & virus & \eng{A virus destroyed all my files.} \\ \hline
to crash & planter & \eng{My computer crashed during the test.} \\ \hline
out of order & hors service & \eng{The printer is out of order.} \\ \hline
\end{tabularx}
\end{center}

\textbf{2. Costs (les coûts)}

\begin{center}
\begin{tabularx}{\linewidth}{|X|X|X|}\hline
\rowcolor{popblueL} \textbf{English} & \textbf{Français} & \textbf{Exemple} \\ \hline
expensive devices & appareils coûteux & \eng{Smartphones are expensive devices for many families.} \\ \hline
data bundle & forfait internet & \eng{She buys a data bundle every week.} \\ \hline
to afford & avoir les moyens de & \eng{Many students cannot afford a laptop.} \\ \hline
maintenance costs & frais de maintenance & \eng{The school cannot pay maintenance costs.} \\ \hline
\end{tabularx}
\end{center}

\textbf{3. Human difficulties (les difficultés humaines)}

\begin{center}
\begin{tabularx}{\linewidth}{|X|X|X|}\hline
\rowcolor{popblueL} \textbf{English} & \textbf{Français} & \textbf{Exemple} \\ \hline
lack of training & manque de formation & \eng{There is a lack of training for teachers.} \\ \hline
computer illiteracy & illectronisme & \eng{Computer illiteracy affects many parents.} \\ \hline
distraction & distraction & \eng{Phones are a source of distraction.} \\ \hline
addiction to social media & addiction aux réseaux sociaux & \eng{Addiction to social media wastes time.} \\ \hline
eye strain & fatigue visuelle & \eng{Long screen time causes eye strain.} \\ \hline
\end{tabularx}
\end{center}

\textbf{4. Risks (les risques)}

\begin{center}
\begin{tabularx}{\linewidth}{|X|X|X|}\hline
\rowcolor{popblueL} \textbf{English} & \textbf{Français} & \textbf{Exemple} \\ \hline
cybercrime & cybercriminalité & \eng{Cybercrime is increasing in big cities.} \\ \hline
hacking & piratage & \eng{Hacking can steal your passwords.} \\ \hline
scam & arnaque & \eng{He lost money in an online scam.} \\ \hline
fake news & fausses informations & \eng{Fake news spreads quickly on WhatsApp.} \\ \hline
privacy & vie privée & \eng{We must protect our privacy online.} \\ \hline
\end{tabularx}
\end{center}

\begin{attention}[Prononciation et pièges lexicaux]
\begin{itemize}
\item \eng{virus} se prononce « vaï-reuss » ; \eng{device} se prononce « di-vaïce » (un appareil), à ne pas confondre avec \eng{devise} (une devise, une devise monétaire).
\item \eng{to afford} se construit souvent avec \eng{can} : \eng{I cannot afford a laptop.} — On ne dit pas \eng{*I cannot afford to buy 500,000 francs} (le verbe \eng{buy} attend un objet, pas un prix). On dit \eng{I cannot afford a laptop} ou \eng{I cannot afford to buy a laptop}.
\item \eng{fake news} est invariable (singulier) : on ne dit pas \emph{fakes news}. Le verbe est au singulier : \eng{Fake news spreads fast.}
\end{itemize}
\end{attention}

\begin{popfigure}
\begin{center}
\begin{tikzpicture}[mindmap, grow cyclic, every node/.style={concept, font=\small}, root concept/.append style={concept color=popblue, font=\bfseries}, level 1/.append style={level distance=3.2cm, sibling angle=90}, level 2/.append style={level distance=2.4cm, sibling angle=45, font=\scriptsize}]
\node[root concept]{ICT difficulties}
  child[concept color=poporange]{ node {Technical} 
      child { node {power cut} }
      child { node {to crash} }
      child { node {virus} } }
  child[concept color=popgreen]{ node {Costs}
      child { node {data bundle} }
      child { node {to afford} } }
  child[concept color=poppurple]{ node {Human}
      child { node {lack of training} }
      child { node {eye strain} } }
  child[concept color=poppink]{ node {Risks}
      child { node {hacking} }
      child { node {fake news} } };
\end{tikzpicture}
\end{center}
\legende{Carte mentale des quatre champs lexicaux des difficultés liées aux TIC.}
\end{popfigure}

\courssub{Connectors: the skeleton of your essay}

Observez ces deux versions du même paragraphe :

\begin{exemplebox}[Avec ou sans connecteurs ?]
\textbf{Sans connecteurs :} \eng{ICTs are useful. They cause problems. Computers are expensive. Many students cannot buy them.}

\textbf{Avec connecteurs :} \eng{Although ICTs are useful, they also cause problems. For instance, computers are expensive; as a result, many students cannot afford them.} — « Bien que les TIC soient utiles, elles posent aussi des problèmes. Par exemple, les ordinateurs coûtent cher ; par conséquent, beaucoup d'élèves ne peuvent pas s'en offrir. »
\end{exemplebox}

La seconde version est plus riche, plus logique, plus mature. Les connecteurs indiquent au lecteur la \emph{relation} entre les idées : ajout, cause, opposition, exemple, conclusion.

\begin{propriete}[Les grandes familles de connecteurs]
\begin{itemize}
\item \textbf{Addition} : \eng{first of all} (tout d'abord), \eng{moreover} (de plus), \eng{furthermore} (en outre), \eng{in addition} (en plus), \eng{besides} (d'ailleurs).
\item \textbf{Cause et conséquence} : \eng{because} (parce que), \eng{since} (puisque), \eng{as a result} (par conséquent), \eng{consequently} (en conséquence), \eng{therefore} (donc), \eng{that is why} (c'est pourquoi).
\item \textbf{Concession et contraste} : \eng{although} (bien que), \eng{however} (cependant), \eng{whereas} (tandis que), \eng{on the one hand... on the other hand} (d'une part... d'autre part), \eng{nevertheless} (néanmoins).
\item \textbf{Exemple et conclusion} : \eng{for instance} (par exemple), \eng{such as} (tel que), \eng{to sum up} (pour résumer), \eng{in conclusion} (en conclusion), \eng{all in all} (somme toute).
\end{itemize}
\end{propriete}

\begin{popfigure}
\begin{center}
\begin{tikzpicture}
\node[draw=popblue, fill=popblueL, rounded corners, minimum width=4.2cm, minimum height=2.2cm, align=center, font=\scriptsize] (add) at (0,2.6) {\textbf{ADDITION}\\ first of all, moreover,\\ furthermore, in addition,\\ besides};
\node[draw=poporange, fill=poporangeL, rounded corners, minimum width=4.2cm, minimum height=2.2cm, align=center, font=\scriptsize] (cause) at (6.2,2.6) {\textbf{CAUSE / RESULT}\\ because, since, as a result,\\ consequently, therefore,\\ that is why};
\node[draw=poppurple, fill=poppurpleL, rounded corners, minimum width=4.2cm, minimum height=2.2cm, align=center, font=\scriptsize] (contrast) at (0,0) {\textbf{CONTRAST}\\ although, however, whereas,\\ on the one hand...\\ on the other hand, nevertheless};
\node[draw=popgreen, fill=popgreenL, rounded corners, minimum width=4.2cm, minimum height=2.2cm, align=center, font=\scriptsize] (ex) at (6.2,0) {\textbf{EXAMPLE / CONCLUSION}\\ for instance, such as,\\ to sum up, in conclusion,\\ all in all};
\draw[-{Stealth}, thick, popmuted] (add) -- (cause);
\draw[-{Stealth}, thick, popmuted] (add) -- (contrast);
\draw[-{Stealth}, thick, popmuted] (cause) -- (ex);
\draw[-{Stealth}, thick, popmuted] (contrast) -- (ex);
\end{tikzpicture}
\end{center}
\legende{Les connecteurs classés par fonction logique : à utiliser pour enchaîner les paragraphes.}
\end{popfigure}

\begin{propriete}[Règles de construction à respecter]
\begin{itemize}
\item \eng{because}, \eng{since}, \eng{although}, \eng{whereas} introduisent une \textbf{subordonnée} (sujet + verbe conjugué) : \eng{Although the connection is slow, students keep trying.} — « Bien que la connexion soit lente, les élèves continuent d'essayer. »
\item \eng{however}, \eng{therefore}, \eng{moreover}, \eng{as a result}, \eng{consequently}, \eng{furthermore} introduisent une \textbf{nouvelle phrase} (ou suivent un point-virgule) et sont suivis d'une virgule : \eng{The connection is slow. However, students keep trying.}
\item \eng{because of} et \eng{due to} sont suivis d'un \textbf{nom} (ou groupe nominal), pas d'un verbe : \eng{Because of a power cut, the lesson stopped.} — « À cause d'une coupure de courant, le cours s'est arrêté. »
\end{itemize}
\end{propriete}

\begin{exemplebox}[Exemples traduits]
\begin{itemize}
\item \eng{As a result, many workers lose their documents.} — « Par conséquent, beaucoup d'employés perdent leurs documents. »
\item \eng{Moreover, data bundles are expensive for most families.} — « De plus, les forfaits internet coûtent cher à la plupart des familles. »
\item \eng{Since many teachers lack training, they avoid using computers.} — « Puisque beaucoup d'enseignants manquent de formation, ils évitent d'utiliser les ordinateurs. »
\item \eng{On the one hand, ICTs save time; on the other hand, they can be a source of distraction.} — « D'une part, les TIC font gagner du temps ; d'autre part, elles peuvent être une source de distraction. »
\item \eng{All in all, technology is helpful only if users are well trained.} — « Somme toute, la technologie n'est utile que si les utilisateurs sont bien formés. »
\end{itemize}
\end{exemplebox}

\begin{attention}[Faux amis et erreurs fréquentes des francophones]
\begin{itemize}
\item \eng{actually} signifie « en fait, en réalité », et non « actuellement ». Pour « actuellement », dites \eng{currently} ou \eng{at present}. \eng{Actually, the computer is out of order.} — « En fait, l'ordinateur est hors service. »
\item « Une déception » ne se dit pas \emph{a deception} : dites \eng{a disappointment}. (\eng{deception} signifie « tromperie ».)
\item Ne traduisez pas « d'abord » par \emph{at first} dans une énumération : utilisez \eng{first of all} ou \eng{firstly}. (\eng{at first} signifie « au début », avec une nuance de changement par la suite.)
\item Attention à la double négation française : « personne ne peut » se dit \eng{nobody can} (et non \emph{nobody cannot}).
\end{itemize}
\end{attention}

\begin{methode}[Varier les connecteurs pour enrichir le style]
\begin{enumerate}
\item Rédigez votre premier jet sans vous soucier des connecteurs.
\item À la relecture, soulignez tous les connecteurs utilisés.
\item Règle d'or : \textbf{ne jamais répéter le même connecteur dans deux paragraphes consécutifs}. Si vous avez écrit \eng{moreover} au paragraphe 2, utilisez \eng{furthermore} ou \eng{in addition} au paragraphe 3.
\item Vérifiez la construction : subordonnée après \eng{although/because}, nouvelle phrase après \eng{however/therefore}.
\item Terminez toujours par un connecteur de conclusion (\eng{to sum up}, \eng{all in all}) pour annoncer clairement votre dernière partie.
\end{enumerate}
\end{methode}

\begin{exoresolu}[Compléter avec le bon connecteur]
Énoncé — Complete the sentences with a suitable connector from the box: \eng{although, as a result, moreover, for instance, to sum up}.

1. \eng{........... the computers are new, the connection is very slow.}

2. \eng{A virus attacked the school server. ..........., all the students' files were lost.}

3. \eng{Many risks exist online, ........... hacking and scams.}

4. \eng{The devices are expensive. ..........., the school cannot afford maintenance costs.}

5. \eng{..........., ICTs are useful but they require training and discipline.}
\tcblower
Solution détaillée :

1. \eng{Although} — il y a opposition entre « ordinateurs neufs » et « connexion lente » ; \eng{although} introduit la subordonnée en tête de phrase.

2. \eng{As a result} — la perte des fichiers est la conséquence de l'attaque du virus ; le connecteur ouvre une nouvelle phrase, suivi d'une virgule.

3. \eng{for instance} — « hacking and scams » sont des exemples de risques. (\eng{such as} serait aussi possible, sans virgule avant.)

4. \eng{Moreover} — on ajoute un argument de coût supplémentaire : les frais de maintenance.

5. \eng{To sum up} — la phrase résume l'ensemble du raisonnement : c'est une conclusion.
\end{exoresolu}

\begin{exercice}[À vous de jouer]
Rewrite the following sentences using the connector in brackets.

1. \eng{The connection is slow. Students continue their research online.} (\eng{however})

2. \eng{Many parents cannot use a computer. They cannot help their children with online homework.} (\eng{because})

3. \eng{ICTs help students learn faster. They also expose them to fake news.} (\eng{on the one hand... on the other hand})
\end{exercice}

\begin{corrige}[Exercice « À vous de jouer »]
1. \eng{The connection is slow. However, students continue their research online.} — « La connexion est lente. Cependant, les élèves poursuivent leurs recherches en ligne. »

2. \eng{Because many parents cannot use a computer, they cannot help their children with online homework.} — « Parce que beaucoup de parents ne savent pas utiliser un ordinateur, ils ne peuvent pas aider leurs enfants à faire leurs devoirs en ligne. »

3. \eng{On the one hand, ICTs help students learn faster; on the other hand, they also expose them to fake news.} — « D'une part, les TIC aident les élèves à apprendre plus vite ; d'autre part, elles les exposent aussi aux fausses informations. »
\end{corrige}

\begin{savaistu}[Les connecteurs au baccalauréat]
Au Bac A, l'usage \textbf{varié et correct} des connecteurs est un critère explicite de la grille de notation de la composition : les correcteurs évaluent la \emph{cohérence} et la \emph{cohésion} du texte. Un essai qui enchaîne mécaniquement dix fois \eng{and} ou \eng{but} perd des points, même si les idées sont bonnes. À l'inverse, un texte qui utilise à bon escient \eng{moreover}, \eng{however}, \eng{as a result} et \eng{all in all} donne immédiatement une impression de maîtrise. Entraînez-vous à utiliser au moins cinq connecteurs différents par composition.
\end{savaistu}

\begin{aretenir}[L'essentiel de la section]
\begin{itemize}
\item Quatre champs lexicaux pour parler des difficultés des TIC : problèmes techniques (\eng{power cut, breakdown, to crash, out of order}), coûts (\eng{data bundle, to afford, maintenance costs}), difficultés humaines (\eng{lack of training, computer illiteracy, eye strain}) et risques (\eng{hacking, scam, fake news, privacy}).
\item Cinq familles de connecteurs : addition, cause/conséquence, contraste, exemple, conclusion.
\item \eng{although/because} + subordonnée ; \eng{however/therefore} + nouvelle phrase avec virgule.
\item Varier les connecteurs : jamais le même dans deux paragraphes consécutifs.
\item Méfiez-vous des faux amis : \eng{actually} = « en fait » ; « actuellement » = \eng{currently}.
\end{itemize}
\end{aretenir}

\coursec{Verb tenses and grammar for this essay}

Dans la section précédente, nous avons rassemblé le vocabulaire et les connecteurs indispensables pour parler des difficultés liées aux TIC (\eng{ICTs}) au travail ou dans l'apprentissage. Mais de beaux mots ne suffisent pas : pour obtenir une excellente note à l'épreuve d'expression écrite du Baccalauréat, il faut les enfermer dans des phrases grammaticalement irréprochables. Cette section vous donne la « boîte à outils grammaticale » de la composition : les trois temps essentiels, les modaux des solutions, la voix passive, le gérondif après préposition et les accords piégeux.

\courssub{Observing the grammar in action}

Commençons par observer un court texte original, du type de celui que vous pourriez produire à l'examen. Lisez-le attentivement et soulignez mentalement les formes verbales.

\begin{textebox}[Reading text — “ICTs at school: progress and problems”]
Nowadays, ICTs play an important role in Cameroonian schools. Computers are used in many offices, and more and more students are using smartphones to do their homework. Technology has changed the way we learn: instead of copying long notes, pupils can download lessons and watch videos at home.

However, difficulties remain. The news is worrying: in many schools, the Internet connection breaks down easily, and much information is lost when there is a power cut. Some teachers have never been trained, so they avoid using computers in class. Worse, data can be stolen when students connect to unsafe networks.

Solutions exist. The government should train teachers regularly, and schools must protect their networks with strong passwords. By working together, teachers, parents and students could make technology a real tool for success, instead of a source of problems.
\end{textebox}

\textbf{Glossaire}

\begin{center}
\begin{tabularx}{\linewidth}{|l|l|X|}
\hline
\rowcolor{popblueL} Mot anglais & Nature & Traduction française \\
\hline
to break down & verbe & tomber en panne \\
\hline
power cut & nom & coupure de courant \\
\hline
to be trained & verbe (passif) & être formé \\
\hline
unsafe network & nom & réseau non sécurisé \\
\hline
password & nom & mot de passe \\
\hline
source of problems & expression & source de problèmes \\
\hline
\end{tabularx}
\end{center}

\textbf{Questions d'observation}

\begin{enumerate}
\item Relevez trois verbes au \cle{present simple}. Que décrivent-ils ?
\item Trouvez une phrase au \cle{present continuous}. Exprime-t-elle un fait permanent ou une tendance actuelle ?
\item Trouvez une phrase au \cle{present perfect}. Pourquoi ce temps est-il choisi ?
\item Relevez deux verbes à la \cle{voix passive} et deux \cle{modaux}. Quelle fonction ont les modaux dans le texte ?
\end{enumerate}

\courssub{The present simple: general facts}

\begin{propriete}[Le present simple : vérités générales et faits permanents]
\textbf{Forme} : sujet + base verbale (+ \textbf{-s} à la 3\up{e} personne du singulier : \eng{he plays}). Négation : \eng{do/does not + base verbale}. Question : \eng{Do/Does + sujet + base verbale ?}\par
\textbf{Emploi} : on l'utilise pour les vérités générales, les faits permanents, les habitudes et les opinions présentées comme des évidences. C'est le temps de base de votre composition.\par
\textbf{Exemple} : \eng{Computers break down easily.} « Les ordinateurs tombent facilement en panne. »
\end{propriete}

\begin{exemplebox}[Exemples traduits]
\begin{itemize}
\item \eng{ICTs play an important role in education.} « Les TIC jouent un rôle important dans l'éducation. »
\item \eng{Many students do not have access to the Internet at home.} « Beaucoup d'élèves n'ont pas accès à Internet à la maison. »
\item \eng{A slow connection wastes a lot of time.} « Une connexion lente fait perdre beaucoup de temps. » (attention au \textbf{-s} de \eng{wastes} : sujet singulier)
\end{itemize}
\end{exemplebox}

\courssub{The present continuous: current trends}

\begin{propriete}[Le present continuous : tendances en cours]
\textbf{Forme} : sujet + \eng{be} (am/is/are) + verbe en \textbf{-ing}.\par
\textbf{Emploi dans une composition} : il décrit une \cle{tendance actuelle}, un changement en train de se produire, souvent avec \eng{nowadays}, \eng{more and more}, \eng{increasingly}, \eng{at present}. C'est le temps idéal pour ouvrir un développement sur l'évolution de la société.\par
\textbf{Comparaison avec le français} : le français utilise le présent simple pour les deux cas (« les élèves utilisent des smartphones ») ; l'anglais, lui, distingue le fait général (\eng{students use smartphones}) de la tendance en cours (\eng{students are using smartphones more and more}).
\end{propriete}

\begin{exemplebox}[Exemples traduits]
\begin{itemize}
\item \eng{More and more students are using smartphones in class.} « De plus en plus d'élèves utilisent des smartphones en classe. »
\item \eng{Companies are storing their data online nowadays.} « Les entreprises stockent leurs données en ligne de nos jours. »
\item \eng{Cybercrime is increasing in many countries.} « La cybercriminalité augmente dans de nombreux pays. »
\end{itemize}
\end{exemplebox}

\courssub{The present perfect: assessing changes}

\begin{propriete}[Le present perfect : le bilan passé-présent]
\textbf{Forme} : sujet + \eng{have/has} + participe passé (\eng{changed}, \eng{become}, \eng{transformed}).\par
\textbf{Emploi} : il fait le bilan d'un changement qui a commencé dans le passé et dont le résultat est visible aujourd'hui. Il est parfait pour les phrases d'introduction et de conclusion.\par
\textbf{Piège} : jamais de date passée précise avec le present perfect ! On dit \eng{Technology has changed our lives} mais \eng{The Internet arrived in Cameroon in the 1990s} (preterite, car la date est donnée).
\end{propriete}

\begin{exemplebox}[Exemples traduits]
\begin{itemize}
\item \eng{Technology has changed the way we learn.} « La technologie a changé notre façon d'apprendre. »
\item \eng{ICTs have become essential in the workplace.} « Les TIC sont devenues essentielles au travail. »
\item \eng{Many schools have not received enough computers yet.} « Beaucoup d'écoles n'ont pas encore reçu assez d'ordinateurs. »
\end{itemize}
\end{exemplebox}

\begin{popfigure}
\begin{center}
\begin{tikzpicture}[font=\small]
\draw[-{Stealth}, thick, popdark] (0,0) -- (12.5,0) node[below left=0.1cm and 0cm, popdark] {temps};
\fill[popblue] (2,0) circle (0.12cm);
\node[above=0.25cm, align=center, text=popblue] at (2,0) {\textbf{Present simple}\\ faits permanents\\ \emph{ICTs play a role.}};
\fill[poporange] (6.2,0) circle (0.12cm);
\node[below=0.25cm, align=center, text=poporange] at (6.2,0) {\textbf{Present continuous}\\ tendances en cours\\ \emph{Students are using phones.}};
\fill[popgreen] (10.2,0) circle (0.12cm);
\node[above=0.25cm, align=center, text=popgreen] at (10.2,0) {\textbf{Present perfect}\\ bilan passé-présent\\ \emph{Technology has changed learning.}};
\draw[popline, dashed] (2,-0.15) -- (2,-1.6);
\draw[popline, dashed] (10.2,-0.15) -- (10.2,-1.6);
\node[align=center, text=popmuted] at (6.2,-1.9) {passé \hspace{2.2cm} aujourd'hui};
\end{tikzpicture}
\end{center}
\legende{Frise des trois temps clés de la composition sur les TIC}
\end{popfigure}

\courssub{Modal verbs for solutions: should, must, could}

Une composition sur les \emph{difficultés} doit obligatoirement proposer des \emph{solutions} : c'est le rôle des modaux.

\begin{propriete}[Les modaux du conseil et de l'obligation]
\textbf{Forme} : modal + base verbale (jamais de \eng{to}, jamais de \textbf{-s}, jamais de \eng{do} à la négation : \eng{should not}).\par
\begin{itemize}
\item \cle{should} = conseil, recommandation (« devrait ») : \eng{The government should train teachers.}
\item \cle{must} = obligation forte, nécessité (« doit ») : \eng{Schools must protect their networks.}
\item \cle{could} = possibilité, suggestion prudente (« pourrait ») : \eng{Teachers could use free online resources.}
\end{itemize}
\end{propriete}

\begin{center}
\begin{tabularx}{\linewidth}{|X|X|X|}
\hline
\rowcolor{popblueL} Français & English & Nuance \\
\hline
Les enseignants devraient être formés. & Teachers should be trained. & conseil \\
\hline
Les élèves doivent sauvegarder leurs données. & Students must save their data. & obligation \\
\hline
On pourrait installer plus d'ordinateurs. & More computers could be installed. & suggestion \\
\hline
\end{tabularx}
\end{center}

\courssub{The passive voice: a formal and useful tool}

\begin{propriete}[La voix passive]
\textbf{Forme} : sujet + \eng{be} (au temps voulu) + participe passé (+ \eng{by} + agent, souvent omis).\par
\textbf{Emploi} : très fréquente dans les écrits formels anglais, elle met l'accent sur le résultat plutôt que sur l'auteur de l'action. Elle se combine avec les modaux : modal + \eng{be} + participe passé.\par
\textbf{Exemples} : \eng{Computers are used in every office.} « On utilise des ordinateurs dans chaque bureau. » / \eng{Data can be stolen.} « Les données peuvent être volées. »
\end{propriete}

\begin{exemplebox}[Passif + modal : la combinaison gagnante]
\eng{Teachers should be trained regularly.} « Les enseignants devraient être formés régulièrement. »\par
\eng{Passwords must be changed often.} « Les mots de passe doivent être changés souvent. »\par
\eng{Broken computers could be repaired locally.} « Les ordinateurs en panne pourraient être réparés localement. »
\end{exemplebox}

\courssub{The gerund after prepositions}

\begin{propriete}[Verbe en -ing après toute préposition]
En anglais, après \textbf{toute} préposition (\eng{without}, \eng{by}, \eng{instead of}, \eng{before}, \eng{after}, \eng{in spite of}...), le verbe se met au \cle{gérondif} en \textbf{-ing}. Le français, lui, utilise l'infinitif : c'est une différence majeure.\par
\begin{itemize}
\item \eng{Students leave the cybercafé without saving their work.} « Les élèves quittent le cybercafé sans sauvegarder leur travail. »
\item \eng{We save paper by using e-mails.} « Nous économisons le papier en utilisant les e-mails. »
\item \eng{Many pupils play games instead of printing their lessons.} « Beaucoup d'élèves jouent au lieu d'imprimer leurs leçons. »
\end{itemize}
\end{propriete}

\courssub{Subject-verb agreement: tricky nouns ending in -s}

\begin{propriete}[Les noms en -s qui sont singuliers]
Certains noms anglais se terminent en \textbf{-s} mais sont \cle{singuliers} : \eng{news}, \eng{mathematics}, \eng{physics}, \eng{economics}, \eng{politics}.\par
\eng{The news is worrying.} « Les nouvelles sont inquiétantes. » (et non \emph{are})\par
\eng{Mathematics is difficult for some learners.} « Les mathématiques sont difficiles pour certains apprenants. »
\end{propriete}

\begin{attention}[« Information » est indénombrable]
Ne dites jamais \emph{informations} avec un \textbf{s} en anglais : \eng{information} est indénombrable. On écrit \eng{much information}, \eng{some information}, \eng{a piece of information}, mais jamais \emph{informations} ni \emph{an information}. Même piège avec \eng{equipment}, \eng{homework}, \eng{advice} et \eng{progress} : \eng{Technology has made great progress.} (sans \textbf{s}).
\end{attention}

\begin{attention}[« To learn » n'est pas « to teach »]
Faux ami classique du francophone : \eng{to learn} = \textbf{apprendre}, \eng{to teach} = \textbf{enseigner}. On dit \eng{ICTs help students to learn} « les TIC aident les élèves à apprendre » et \eng{Teachers teach ICT skills} « les enseignants enseignent les compétences numériques ». Ne dites jamais \emph{The teacher learns us English} !
\end{attention}

\begin{savaistu}[Pourquoi « news » prend-il un verbe au singulier ?]
Le mot \eng{news} ressemble à un pluriel, mais son \textbf{s} final ne vient pas du pluriel : il s'agit à l'origine de l'expression \eng{new things} (« choses nouvelles »), devenue un nom indénombrable unique. C'est pourquoi les journalistes anglophones disent toujours \eng{The news is at eight o'clock} et jamais \emph{the news are}. Retenez aussi que l'on dit \eng{a piece of news} pour « une nouvelle ».
\end{savaistu}

\begin{methode}[Choisir le bon temps ou la bonne forme en trois questions]
\begin{enumerate}
\item \textbf{Est-ce un fait général, une vérité permanente ?} $\rightarrow$ \cle{present simple} : \eng{Computers break down easily.}
\item \textbf{Est-ce une tendance actuelle ou un changement dont on voit le résultat ?} $\rightarrow$ tendance en cours : \cle{present continuous} (\eng{More students are using phones.}) ; bilan avec résultat présent : \cle{present perfect} (\eng{Technology has changed learning.}).
\item \textbf{Est-ce un conseil, une solution, une recommandation ?} $\rightarrow$ \cle{modal} (\eng{should/must/could}), souvent à la voix passive : \eng{Teachers should be trained.}
\end{enumerate}
\end{methode}

\begin{center}
\begin{tabularx}{\linewidth}{|l|X|X|}
\hline
\rowcolor{popblueL} Outil grammatical & Emploi dans la composition & Exemple \\
\hline
Present simple & faits généraux & \eng{ICTs play an important role.} \\
\hline
Present continuous & tendances actuelles & \eng{More students are using smartphones.} \\
\hline
Present perfect & bilan des changements & \eng{Technology has changed the way we learn.} \\
\hline
Modaux & solutions et conseils & \eng{The government should train teachers.} \\
\hline
Voix passive & style formel, résultat & \eng{Data can be stolen.} \\
\hline
Gérondif & après préposition & \eng{instead of printing} \\
\hline
\end{tabularx}
\end{center}

\begin{exoresolu}[Put the verbs in the correct form]
Énoncé — Mettez les verbes entre parenthèses à la forme correcte :
\begin{enumerate}
\item Nowadays, more and more workers (use) e-mails to communicate.
\item Technology (transform) the way companies work.
\item A slow Internet connection (waste) a lot of time.
\item The government should (train) all teachers in ICT skills.
\item Many students chat online instead of (revise) their lessons.
\item The news about cybercrime (be) worrying.
\end{enumerate}
\tcblower
Solution détaillée :
\begin{enumerate}
\item \eng{are using} — tendance actuelle signalée par \eng{nowadays} et \eng{more and more} : present continuous.
\item \eng{has transformed} — bilan d'un changement dont le résultat est visible aujourd'hui : present perfect (\eng{technology} est singulier : \eng{has}).
\item \eng{wastes} — fait général : present simple, avec le \textbf{-s} de la 3\up{e} personne du singulier.
\item \eng{train} — après un modal, toujours la base verbale, sans \eng{to} et sans \textbf{-s}.
\item \eng{revising} — après la préposition \eng{instead of}, le verbe prend la forme en \textbf{-ing}.
\item \eng{is} — \eng{news} est singulier malgré son \textbf{s} final.
\end{enumerate}
\end{exoresolu}

\begin{exercice}[Your turn]
Réécrivez les phrases suivantes en corrigeant l'erreur grammaticale :
\begin{enumerate}
\item \emph{Teachers learns students how to use computers.}
\item \emph{I found many useful informations on the Internet.}
\item \emph{The government should to build more computer rooms.}
\item \emph{Students save time by use online dictionaries.}
\item \emph{The news are good: our school has received ten new computers.}
\end{enumerate}
\end{exercice}

\begin{corrige}[Exercice « Your turn »]
\begin{enumerate}
\item \eng{Teachers teach students how to use computers.} (\eng{teach} = enseigner ; \eng{learn} = apprendre.)
\item \eng{I found much useful information on the Internet.} (\eng{information} est indénombrable : pas de \textbf{s}, et \eng{much} à la place de \eng{many}.)
\item \eng{The government should build more computer rooms.} (jamais de \eng{to} après un modal.)
\item \eng{Students save time by using online dictionaries.} (gérondif après la préposition \eng{by}.)
\item \eng{The news is good: our school has received ten new computers.} (\eng{news} est singulier.)
\end{enumerate}
\end{corrige}

\begin{aretenir}[L'essentiel grammatical de votre composition]
\begin{itemize}
\item Fait général $\rightarrow$ present simple ; tendance actuelle $\rightarrow$ present continuous ; bilan $\rightarrow$ present perfect.
\item Solutions $\rightarrow$ \eng{should / must / could} + base verbale, souvent au passif : \eng{Teachers should be trained.}
\item Après une préposition, verbe en \textbf{-ing} : \eng{without saving}, \eng{by using}, \eng{instead of printing}.
\item \eng{Information}, \eng{news}, \eng{equipment}, \eng{homework} : indénombrables ou singuliers, jamais de \textbf{s} ni de pluriel.
\item \eng{Learn} = apprendre ; \eng{teach} = enseigner.
\end{itemize}
\end{aretenir}

Vous disposez maintenant de tous les outils grammaticaux nécessaires. Dans la section suivante, nous les verrons réunis dans un modèle de composition commenté, avant de passer en revue les erreurs à éviter et la grille d'évaluation.

\coursec{A commented model essay}

After mastering the verb tenses and grammar tools in the previous section, we now turn to a full‑length model composition. Studying a well‑structured example is the most efficient way to internalise the organisation, the connectors and the tone expected in the Baccalauréat writing paper. Read the essay once for general meaning, then a second time to dissect its architecture — this is the method we will apply together.

\courssub{Reading the model essay}

\begin{textebox}{Model composition — Difficulties with the Use of ICTs at Work and for Learning}
Nowadays, computers, smartphones and the Internet are everywhere. From the office in Douala to the classroom in Buea, Information and Communication Technologies (ICTs) have become indispensable tools for working, studying and communicating. However, despite their obvious advantages, they also bring serious difficulties that cannot be ignored.

First of all, technical problems are a daily reality for many users. In Cameroon, frequent power cuts — locally known as “délestage” — interrupt online lessons and stop office work for hours. Moreover, Internet connection is often slow or unstable, especially in rural areas. For instance, a student in a village near Bamenda may wait ten minutes just to load a single page of an educational platform. Consequently, precious time is wasted and deadlines are missed.

Furthermore, the financial cost of ICTs excludes many families and small businesses. A decent laptop costs at least 300,000 FCFA, and a monthly data bundle can reach 10,000 FCFA. Many parents cannot afford these expenses, so their children fall behind at school because they cannot do online research or submit assignments on time. Similarly, small enterprises struggle to equip their staff with modern devices, which reduces their competitiveness.

In addition, human factors create further obstacles. Numerous teachers and employees lack proper training to use digital tools efficiently. As a result, they feel anxious and make avoidable mistakes. Besides, the Internet is a source of distraction: social media and video games tempt students and workers away from their tasks. Last but not least, cybercrime — such as phishing and identity theft — threatens both personal data and company secrets.

All in all, ICTs are useful but demanding tools. To benefit fully from them, schools and companies should invest in continuous training and in alternative energy sources like solar panels. Only then will the digital divide be bridged and productivity improved.
\end{textebox}

\noindent
\textbf{Glossaire des mots difficiles}
\begin{center}
\begin{tabularx}{\linewidth}{|X|X|X|}
\hline
\rowcolor{popblueL} \textbf{English} & \textbf{Nature} & \textbf{Français} \\
\hline
indispensable & adj. & indispensable \\
délestage & n. (masc.) & coupure programmée d'électricité \\
unstable & adj. & instable \\
rural & adj. & rural \\
decent & adj. & correct, convenable \\
data bundle & n. & forfait données (internet) \\
afford & v. & avoir les moyens de s'offrir \\
fall behind & v. (phrasal) & prendre du retard \\
competitiveness & n. (fém.) & compétitivité \\
efficiently & adv. & efficacement \\
anxious & adj. & anxieux, inquiet \\
distraction & n. (fém.) & distraction \\
tempt & v. & tenter \\
phishing & n. (masc.) & hameçonnage (arnaque par mail) \\
identity theft & n. (fém.) & usurpation d'identité \\
bridge the gap & loc. & combler le fossé \\
\hline
\end{tabularx}
\end{center}

\courssub{Analysing the introduction: hook and thesis}

The opening paragraph performs two essential functions: it catches the reader’s attention (\textbf{hook}) and states the writer’s position (\textbf{thesis statement}).

\begin{exemplebox}{Hook}
\eng{Nowadays, computers, smartphones and the Internet are everywhere.} « De nos jours, les ordinateurs, les smartphones et Internet sont partout. » \\
\eng{From the office in Douala to the classroom in Buea, Information and Communication Technologies (ICTs) have become indispensable tools for working, studying and communicating.} « Du bureau à Douala à la salle de classe à Buea, les TIC sont devenus des outils indispensables pour travailler, étudier et communiquer. »
\end{exemplebox}

\begin{propriete}[The hook]
\begin{itemize}
\item Presents the topic in a general, factual way.
\item Uses a \textbf{spatial reference} (Douala, Buea) to anchor the essay in the Cameroonian context — examiners appreciate this local colour.
\item Employs the present simple (\eng{are}, \eng{have become}) to describe a current reality.
\end{itemize}
\end{propriete}

\begin{exemplebox}{Thesis statement}
\eng{However, despite their obvious advantages, they also bring serious difficulties that cannot be ignored.} « Cependant, malgré leurs avantages évidents, elles apportent aussi de sérieuses difficultés qu'on ne peut ignorer. »
\end{exemplebox}

\begin{propriete}[The thesis statement]
\begin{itemize}
\item Signals a \textbf{contrast} with the hook through the connector \eng{However}.
\item Concedes a point (\eng{despite their obvious advantages}) to show balance — this is a mark of maturity.
\item Announces the \textbf{controlling idea}: the essay will discuss \textbf{difficulties} (not advantages).
\item Uses the modal \eng{cannot be ignored} (passive voice) for an objective, authoritative tone.
\end{itemize}
\end{propriete}

\begin{attention}[Piège fréquent]
Ne confondez pas \textbf{thesis statement} et \textbf{topic sentence}. La thèse guide \textbf{tout l'essai} ; la topic sentence guide \textbf{un seul paragraphe}. À l'examen, soulignez votre thèse dans l'introduction pour la repérer vite.
\end{attention}

\courssub{Analysing body paragraph 1: Technical difficulties}

\begin{exemplebox}{Topic sentence + development}
\eng{First of all, technical problems are a daily reality for many users.} \\
\eng{In Cameroon, frequent power cuts — locally known as “délestage” — interrupt online lessons and stop office work for hours.} \\
\eng{Moreover, Internet connection is often slow or unstable, especially in rural areas.} \\
\eng{For instance, a student in a village near Bamenda may wait ten minutes just to load a single page of an educational platform.} \\
\eng{Consequently, precious time is wasted and deadlines are missed.}
\end{exemplebox}

\begin{propriete}[Structure of a body paragraph]
\begin{enumerate}
\item \textbf{Topic sentence} (\eng{First of all, technical problems…}) : annonce l'idée directrice du paragraphe.
\item \textbf{Explanation + local example} (\eng{In Cameroon… délestage}) : ancre l'argument dans le réel.
\item \textbf{Additional argument + connector} (\eng{Moreover, Internet connection…}) : développe une seconde facette.
\item \textbf{Concrete illustration} (\eng{For instance, a student…}) : rend l'abstrait tangible.
\item \textbf{Consequence + connector} (\eng{Consequently, precious time…}) : montre l'impact réel.
\end{enumerate}
\end{propriete}

\begin{center}
\begin{tabularx}{\linewidth}{|X|X|X|}
\hline
\rowcolor{popblueL} \textbf{Connector} & \textbf{Function} & \textbf{French equivalent} \\
\hline
First of all & Opening / enumerating & Tout d'abord \\
Moreover & Adding a related point & De plus / En outre \\
For instance & Giving an example & Par exemple \\
Consequently & Showing result / consequence & Par conséquent / Du coup \\
\hline
\end{tabularx}
\end{center}

\begin{savaistu}[Le contexte camerounais paye]
Les correcteurs du Baccalauréat valorisent les \textbf{exemples locaux pertinents} : mentionner le « délestage », les cybercafés de quartier, les campus numériques de l'Université de Yaoundé I ou de Buea, les bendskins qui transportent des routeurs 4G… montre que vous maîtrisez le sujet \emph{et} votre environnement. Évitez les généralités vagues (« In Africa, electricity is bad ») — soyez précis.
\end{savaistu}

\courssub{Analysing body paragraph 2: Financial difficulties}

\begin{exemplebox}{Topic sentence + development}
\eng{Furthermore, the financial cost of ICTs excludes many families and small businesses.} \\
\eng{A decent laptop costs at least 300,000 FCFA, and a monthly data bundle can reach 10,000 FCFA.} \\
\eng{Many parents cannot afford these expenses, so their children fall behind at school because they cannot do online research or submit assignments on time.} \\
\eng{Similarly, small enterprises struggle to equip their staff with modern devices, which reduces their competitiveness.}
\end{exemplebox}

\begin{propriete}[Connectors of addition and similarity]
\begin{itemize}
\item \eng{Furthermore} = \eng{Moreover} = \eng{In addition} : ajoutent un argument \textbf{nouveau} (ici, le coût financier après le technique).
\item \eng{Similarly} : établit une \textbf{analogie} entre deux situations (familles / petites entreprises).
\item \eng{so} (coordinating conjunction) : lie cause (\eng{cannot afford}) et conséquence (\eng{fall behind}).
\item \eng{which} (relative pronoun) : renvoie à toute la proposition précédente (\eng{small enterprises struggle…}) pour en tirer une conséquence.
\end{itemize}
\end{propriete}

\begin{exoresolu}{Repérer les connecteurs et leur fonction}
\textbf{Énoncé :} Dans le paragraphe 2, soulignez chaque connecteur et indiquez sa fonction (addition, illustration, conséquence, analogie).

\textbf{Solution détaillée :}
\begin{enumerate}
\item \eng{Furthermore} — \textbf{addition} : introduit un nouvel argument (financier) après le technique.
\item \eng{so} — \textbf{conséquence} : relie \eng{cannot afford} (cause) à \eng{fall behind} (effet).
\item \eng{Similarly} — \textbf{analogie} : compare la situation des familles à celle des petites entreprises.
\item \eng{which} — \textbf{conséquence (relative)} : renvoie à la difficulté des entreprises pour en déduire la perte de compétitivité.
\end{enumerate}
\end{exoresolu}

\courssub{Analysing body paragraph 3: Human difficulties}

\begin{exemplebox}{Topic sentence + development}
\eng{In addition, human factors create further obstacles.} \\
\eng{Numerous teachers and employees lack proper training to use digital tools efficiently.} \\
\eng{As a result, they feel anxious and make avoidable mistakes.} \\
\eng{Besides, the Internet is a source of distraction: social media and video games tempt students and workers away from their tasks.} \\
\eng{Last but not least, cybercrime — such as phishing and identity theft — threatens both personal data and company secrets.}
\end{exemplebox}

\begin{propriete}[Connectors for the third argument and internal logic]
\begin{itemize}
\item \eng{In addition} : signale le \textbf{troisième grand axe} (humain) après technique et financier.
\item \eng{As a result} : conséquence directe du manque de formation (\eng{anxious}, \eng{mistakes}).
\item \eng{Besides} : ajoute un argument \textbf{secondaire mais important} (distraction) au sein du même paragraphe.
\item \eng{Last but not least} : marque le \textbf{point final et souvent le plus grave} (cybercrime) — expression idiomatique à connaître.
\item \eng{such as} : introduit des \textbf{exemples précis} d'un terme général (\eng{cybercrime}).
\end{itemize}
\end{propriete}

\begin{center}
\begin{tabularx}{\linewidth}{|X|X|X|}
\hline
\rowcolor{popblueL} \textbf{Connector / Structure} & \textbf{Role in paragraph 3} & \textbf{Français} \\
\hline
In addition & Third main argument & De plus / En outre \\
As a result & Consequence of lack of training & Par conséquent \\
Besides & Additional point within paragraph & De plus / Par ailleurs \\
Last but not least & Final, weighty point & Last but not least / Pour finir \\
such as & Exemplification & comme / tels que \\
\hline
\end{tabularx}
\end{center}

\courssub{Analysing the conclusion: summary and recommendation}

\begin{exemplebox}{Conclusion}
\eng{All in all, ICTs are useful but demanding tools.} \\
\eng{To benefit fully from them, schools and companies should invest in continuous training and in alternative energy sources like solar panels.} \\
\eng{Only then will the digital divide be bridged and productivity improved.}
\end{exemplebox}

\begin{propriete}[Anatomy of a strong conclusion]
\begin{enumerate}
\item \textbf{Summary connector} : \eng{All in all} (synonymes : \eng{In conclusion}, \eng{To sum up}, \eng{Overall}) — résume la thèse sans la répéter mot pour mot.
\item \textbf{Balanced assessment} : \eng{useful but demanding} — nuance, évite le manichéisme.
\item \textbf{Recommendation / Call to action} : \eng{should invest in…} — modal \eng{should} + acteurs concrets (\eng{schools and companies}) + solutions précises (\eng{training}, \eng{solar panels}).
\item \textbf{Inversion for emphasis} : \eng{Only then will the digital divide be bridged…} — structure \eng{Only + adverbial + auxiliary + subject + verb} pour insister sur la condition nécessaire. Cette inversion est un marqueur de niveau C1/C2 très apprécié au Bac.
\end{enumerate}
\end{propriete}

\begin{attention}[Erreurs à éviter en conclusion]
\begin{itemize}
\item \textbf{Nouvel argument} : n'introduisez jamais une idée nouvelle (ex. « Also, 5G will solve everything »).
\item \textbf{Répétition servile} : ne recopiez pas la thèse mot pour mot. Reformulez (\eng{difficulties} $\rightarrow$ \eng{demanding tools}).
\item \textbf{Conclusion trop courte} : une seule phrase (\eng{In conclusion, ICTs have problems.}) est insuffisante. Visez 3--4 phrases.
\end{itemize}
\end{attention}

\courssub{Annotated structure of the model}

The following diagram maps each paragraph to its rhetorical function and highlights the logical flow. Use it as a checklist when you plan your own essay.

\begin{popfigure}
\begin{tikzpicture}[
  node distance=1.2cm and 2.5cm,
  box/.style={draw=popblue, thick, rounded corners, fill=popblueL, text width=3.2cm, align=center, minimum height=1.8cm},
  arrow/.style={-Stealth, thick, popdark},
  label/.style={font=\small, text=popmuted, align=center}
]
% Nodes
\node[box, align=center] (intro){Introduction\\ \textbf{Hook} + \textbf{Thesis}};
\node[box, below=of intro, align=center] (p1){Paragraph 1\\ \textbf{Technical difficulties}\\ (power cuts, slow net)};
\node[box, below=of p1, align=center] (p2){Paragraph 2\\ \textbf{Financial difficulties}\\ (cost, exclusion)};
\node[box, below=of p2, align=center] (p3){Paragraph 3\\ \textbf{Human difficulties}\\ (training, distraction, cybercrime)};
\node[box, below=of p3, align=center] (conc){Conclusion\\ \textbf{Summary} + \textbf{Recommendation}};

% Arrows
\draw[arrow] (intro.south) -- (p1.north) node[midway, right, label] {First of all};
\draw[arrow] (p1.south) -- (p2.north) node[midway, right, label] {Furthermore};
\draw[arrow] (p2.south) -- (p3.north) node[midway, right, label] {In addition};
\draw[arrow] (p3.south) -- (conc.north) node[midway, right, label] {All in all};

% Side annotations
\node[label, left=0.3cm of intro.west] {environ 40 mots};
\node[label, left=0.3cm of p1.west] {environ 60 mots};
\node[label, left=0.3cm of p2.west] {environ 60 mots};
\node[label, left=0.3cm of p3.west] {environ 60 mots};
\node[label, left=0.3cm of conc.west] {environ 40 mots};

% Connectors highlighted
\node[label, right=0.3cm of intro.east] {\eng{However, despite…}};
\node[label, right=0.3cm of p1.east] {\eng{Moreover, For instance, Consequently}};
\node[label, right=0.3cm of p2.east] {\eng{Similarly, so, which}};
\node[label, right=0.3cm of p3.east] {\eng{As a result, Besides, Last but not least}};
\node[label, right=0.3cm of conc.east] {\eng{should invest, Only then will…}};
\end{tikzpicture}
\legende{Schéma annoté du modèle : structure paragraphique, connecteurs principaux et volume indicatif par section.}
\end{popfigure}

\courssub{How to exploit a model essay — step‑by‑step method}

\begin{methode}{Exploiter un modèle de composition}
\begin{enumerate}
\item \textbf{Lecture globale} (2 min) : lisez une première fois \textbf{sans stylo} pour saisir le sens général, le ton et la position de l'auteur.
\item \textbf{Repérage structurel} (3 min) : soulignez la \textbf{thèse} dans l'introduction ; encadrez la \textbf{topic sentence} de chaque paragraphe ; cerclez les \textbf{connecteurs} logiques.
\item \textbf{Analyse lexicale} (3 min) : relevez 5--7 expressions utiles (ex. \eng{fall behind}, \eng{bridge the gap}, \eng{Last but not least}) et notez leur traduction.
\item \textbf{Analyse grammaticale} (2 min) : identifiez 2--3 structures ciblées (passif, inversion, relativisation, modalités) que vous pourrez réutiliser.
\item \textbf{Transfert} (10 min) : sur un \textbf{nouveau sujet} (ex. « Advantages and disadvantages of social media for students »), rédigez votre propre plan en imitant la \textbf{structure} (pas le contenu) du modèle.
\item \textbf{Auto-évaluation} : comparez votre brouillon à la grille d'évaluation (cohérence, connecteurs, vocabulaire, grammaire, longueur).
\end{enumerate}
\end{methode}

\begin{attention}[Ne pas apprendre par cœur]
\textbf{Interdiction formelle :} ne recopiez jamais ce modèle (ou un autre) par cœur le jour de l'examen. Les correcteurs détectent immédiatement les essais « photocopiés » (mêmes exemples, mêmes chiffres, même formulation). \textbf{S'inspirer} = reprendre la \textbf{structure}, les \textbf{connecteurs}, le \textbf{ton} et quelques \textbf{tournures} pour produire un texte \textbf{personnel} adapté au sujet posé.
\end{attention}

\begin{savaistu}[Ce que les correcteurs du Bac recherchent]
\begin{itemize}
\item \textbf{Pertinence locale} : exemples camerounais réels (délestage, cybercafés, campus numériques, forfaits MTN/Orange, bendskins, cacao, football).
\item \textbf{Cohérence logique} : chaque paragraphe = une idée + explication + exemple + conséquence.
\item \textbf{Richesse connecteur} : au moins 8--10 connecteurs variés (pas seulement \eng{and}, \eng{but}, \eng{because}).
\item \textbf{Précision lexicale} : vocabulaire technique (\eng{bandwidth}, \eng{phishing}, \eng{digital divide}, \eng{solar panel}) + verbes d'action (\eng{afford}, \eng{struggle}, \eng{bridge}, \eng{enhance}).
\item \textbf{Maîtrise grammaticale} : passif, relativisation, inversion (\eng{Only then will…}), conditionnel (\eng{should invest}), modaux.
\end{itemize}
\end{savaistu}

\courssub{Consolidation exercise: reconstruct the logical flow}

\begin{exercice}{Reconstituer la progression logique}
Les phrases suivantes sont extraites du modèle mais mélangées. Remettez-les dans l'ordre chronologique de l'essai (1 à 10) et indiquez pour chacune : \textbf{paragraphe} (Intro, P1, P2, P3, Concl.) et \textbf{fonction} (Hook, Thesis, Topic sentence, Example, Consequence, Recommendation, etc.).
\begin{enumerate}
\item A decent laptop costs at least 300,000 FCFA.
\item Consequently, precious time is wasted and deadlines are missed.
\item Only then will the digital divide be bridged and productivity improved.
\item In Cameroon, frequent power cuts — locally known as “délestage” — interrupt online lessons.
\item All in all, ICTs are useful but demanding tools.
\item Nowadays, computers, smartphones and the Internet are everywhere.
\item Furthermore, the financial cost of ICTs excludes many families and small businesses.
\item Last but not least, cybercrime — such as phishing and identity theft — threatens both personal data and company secrets.
\item However, despite their obvious advantages, they also bring serious difficulties that cannot be ignored.
\item In addition, human factors create further obstacles.
\end{enumerate}
\end{exercice}

\begin{corrige}{Exercice n°1}
\begin{enumerate}
\item \textbf{6 — Intro — Hook} : ouverture générale.
\item \textbf{9 — Intro — Thesis} : position de l'auteur (contraste \eng{However}).
\item \textbf{4 — P1 — Explanation + local example} : illustration technique (délestage).
\item \textbf{2 — P1 — Consequence} (\eng{Consequently}) : résultat des problèmes techniques.
\item \textbf{7 — P2 — Topic sentence} : nouvel argument (financier) introduit par \eng{Furthermore}.
\item \textbf{1 — P2 — Concrete detail} : chiffre précis (300 000 FCFA).
\item \textbf{10 — P3 — Topic sentence} : troisième axe (humain) introduit par \eng{In addition}.
\item \textbf{8 — P3 — Final weighty example} (\eng{Last but not least}) : cybercriminalité.
\item \textbf{5 — Concl. — Summary} (\eng{All in all}) : synthèse nuancée.
\item \textbf{3 — Concl. — Emphatic recommendation} (\eng{Only then will…}) : inversion pour insister sur la condition.
\end{enumerate}
\end{corrige}

\coursec{Common mistakes and assessment grid}

Dans la section précédente, nous avons lu et commenté un modèle de composition réussie sur les difficultés liées à l'usage des TIC au travail ou dans les études. Mais entre le modèle idéal et la copie réelle, il y a souvent un fossé. Cette dernière section est peut-être la plus rentable de toute la leçon : connaître les erreurs typiques — celles que les correcteurs du Baccalauréat voient revenir chaque année — et savoir comment une copie est notée, c'est déjà gagner des points. Nous allons d'abord observer une copie faible, puis classer les erreurs en trois catégories, et enfin découvrir la grille d'évaluation et la méthode de relecture C-O-L.

\courssub{Observing a weak essay: spot the mistakes}

Lisons d'abord cet extrait de copie d'élève. Il contient volontairement les erreurs les plus fréquentes des candidats francophones. Essayez de les repérer avant de lire la correction.

\begin{textebox}[A weak student's essay — can you find the mistakes?]
The difficulties of ICTs in learning

I am agree that ICTs are very important today. In my school in Yaoundé, we have many informatics tools but there is many problems.

Because the connection is very slow, so the students cannot download their lessons. The teachers give us many informations and advices but we cannot use them well. Last week, I assisted at a computer lesson and the electricity went off after ten minutes. But the teacher continued. But we could not do the exercises. But it was very frustrating.

ICTs have many advantages. They help us to communicate, to find documents, to learn English, to watch videos, to play games and to listen to music. Technology is the future of the world and every student must have a computer and a smartphone to succeed in life and to be modern like the students in Europe and America.
\end{textebox}

\textbf{Glossaire :} \emph{slow} (adj., lent) ; \emph{to download} (v., télécharger) ; \emph{to go off} (v., ici : se couper, à propos du courant électrique) ; \emph{frustrating} (adj., frustrant).

Combien d'erreurs avez-vous trouvées ? Cette copie en contient au moins \cle{dix}, réparties en trois catégories : erreurs de langue, erreurs d'organisation et hors sujet. Analysons-les une par une.

\courssub{Language mistakes: uncountable nouns and false friends}

\begin{propriete}[Les noms indénombrables ne prennent jamais de -s]
En anglais, certains noms sont \cle{indénombrables} (\emph{uncountable nouns}) : ils désignent une masse ou une notion globale et n'ont \textbf{jamais} de pluriel en -s, même si leur équivalent français est pluriel. Ils s'emploient avec un verbe au \textbf{singulier}, sans article \emph{a/an}, et avec \emph{much} ou \emph{a little} plutôt que \emph{many} ou \emph{a few}. Pour les compter, on utilise \emph{a piece of} : \eng{a piece of information} « une information », \eng{two pieces of advice} « deux conseils ».
\end{propriete}

\begin{center}
\begin{tabularx}{\linewidth}{|X|X|X|}
\hline
\rowcolor{popblueL} Erreur fréquente & Forme correcte & Traduction \\
\hline
many informations & much information / a lot of information & beaucoup d'informations \\
\hline
many advices & some advice / pieces of advice & des conseils \\
\hline
a good new & a piece of good news & une bonne nouvelle \\
\hline
homeworks & homework & les devoirs \\
\hline
equipments & equipment & le matériel, l'équipement \\
\hline
knowledges & knowledge & les connaissances \\
\hline
\end{tabularx}
\end{center}

\begin{exemplebox}[Exemples traduits]
\eng{The teacher gave us useful information about the exam.} « Le professeur nous a donné des informations utiles sur l'examen. »

\eng{My father gave me a piece of advice: save your work every five minutes.} « Mon père m'a donné un conseil : sauvegarde ton travail toutes les cinq minutes. »

\eng{The news is good: our school has received new computers.} « Les nouvelles sont bonnes : notre école a reçu de nouveaux ordinateurs. » (Attention : \emph{news} se termine en -s mais reste \textbf{singulier}.)
\end{exemplebox}

\begin{attention}[Faux amis et calques du français]
Ne traduisez jamais mot à mot depuis le français :
\begin{itemize}
\item \eng{I am agree} est impossible : \emph{agree} est un \textbf{verbe}, pas un adjectif. Écrivez \eng{I agree with this idea.} « Je suis d'accord avec cette idée. » De même, \eng{The students are not agree with the new rule.} est faux ; dites \eng{The students do not agree with the new rule.} « Les élèves ne sont pas d'accord avec la nouvelle règle. »
\item \eng{to assist at a lesson} est un calque d'« assister à ». Dites \eng{to attend a lesson} « suivre un cours ». Le verbe \emph{to assist} existe, mais il signifie « aider » : \eng{She assisted the teacher.} « Elle a aidé le professeur. »
\item \eng{informatics tools} : le mot \emph{informatics} est rare en anglais courant ; dites \eng{computer tools} ou \eng{digital tools}.
\item \eng{there is many problems} : avec un nom pluriel, utilisez \eng{there are many problems}.
\end{itemize}
\end{attention}

\begin{exemplebox}[Les formes correctes à retenir]
\eng{I attended a computer lesson yesterday.} « J'ai assisté à un cours d'informatique hier. »

\eng{There are many problems with the school network.} « Il y a de nombreux problèmes avec le réseau de l'école. »

\eng{I completely agree with you.} « Je suis tout à fait d'accord avec toi. »
\end{exemplebox}

\courssub{Connector mistakes: "because... so" and repeated "but"}

\begin{attention}[\eng{Because} et \eng{so} ne s'emploient jamais ensemble]
Sous l'influence du français « parce que... donc... », de nombreux élèves écrivent \eng{Because the connection was slow, so I left.} C'est \textbf{incorrect} en anglais : les deux connecteurs expriment la même relation de cause à conséquence, il faut choisir \textbf{l'un ou l'autre} :
\begin{itemize}
\item \eng{Because the connection was slow, I left.} « Comme la connexion était lente, je suis parti. »
\item \eng{The connection was slow, so I left.} « La connexion était lente, donc je suis parti. »
\end{itemize}
Même règle pour \emph{although} et \emph{but} : \eng{Although it was raining, we went out.} ou \eng{It was raining, but we went out.} « Bien qu'il pleuve, nous sommes sortis. » — jamais les deux ensemble.
\end{attention}

\begin{propriete}[Varier les connecteurs]
Commencer trois phrases de suite par \emph{But} (comme dans la copie faible) appauvrit le texte et coûte des points d'organisation. Remplacez par : \emph{however, nevertheless, yet, on the contrary, despite this}. Et rappelez-vous : \emph{however} en début de phrase est suivi d'une virgule : \eng{However, the teacher continued the lesson.} « Cependant, le professeur a continué le cours. »
\end{propriete}

\courssub{Structural mistakes and going off-topic}

Trois erreurs de structure reviennent constamment dans les copies :

\begin{itemize}
\item \textbf{Absence de problématique} : l'introduction doit annoncer clairement le problème traité, par exemple \eng{Although ICTs are useful, they create serious difficulties for students and workers.} « Bien que les TIC soient utiles, elles créent de sérieuses difficultés pour les élèves et les travailleurs. » Sans cette phrase d'accroche, le lecteur ne sait pas où vous allez.
\item \textbf{Paragraphes non délimités} : chaque idée principale = un paragraphe, visible par un alinéa ou une ligne de saut. Un bloc unique de 250 mots est sanctionné.
\item \textbf{Conclusion manquante} : beaucoup de candidats s'arrêtent au milieu du développement, faute de temps.
\end{itemize}

\begin{propriete}[La conclusion est obligatoire]
Au Baccalauréat, une composition \textbf{sans conclusion perd automatiquement des points d'organisation}, même si le reste est excellent. Prévoyez toujours deux ou trois phrases de conclusion : un résumé de votre position (\eng{To conclude, ICTs are powerful tools, but...} « Pour conclure, les TIC sont des outils puissants, mais... ») et éventuellement une ouverture ou une recommandation (\eng{Schools should therefore train both teachers and students to use them wisely.} « Les écoles devraient donc former les enseignants comme les élèves à les utiliser judicieusement. »).
\end{propriete}

\begin{attention}[Le hors sujet : le piège des avantages]
Le sujet porte sur les \textbf{difficultés} (\emph{difficulties}) des TIC au travail ou pour apprendre — pas sur leurs avantages. Dans la copie faible, tout le dernier paragraphe vante les mérites de la technologie : c'est un \cle{hors sujet} partiel, lourdement sanctionné. Vous pouvez \emph{concéder} brièvement un avantage (\eng{Admittedly, ICTs make research faster. However, ...} « Certes, les TIC accélèrent la recherche. Cependant... ») à condition de revenir immédiatement aux difficultés : coupures d'électricité, connexion lente, coût des données, distraction, manque de formation.
\end{attention}

\begin{exoresolu}[Corriger une phrase de copie]
Énoncé : la phrase suivante contient quatre erreurs. Identifiez-les et réécrivez la phrase correctement.

\eng{Because I am agree that ICTs give many informations, so I assisted at the computer lesson with pleasure.}
\tcblower
Solution détaillée :
\begin{enumerate}
\item \eng{I am agree} $\rightarrow$ \eng{I agree} (\emph{agree} est un verbe, pas un adjectif).
\item \eng{informations} $\rightarrow$ \eng{information} (indénombrable, jamais de -s ; on dira aussi \eng{useful information} plutôt que \eng{many information}).
\item \eng{Because ... so} ensemble $\rightarrow$ supprimer l'un des deux connecteurs.
\item \eng{assisted at} $\rightarrow$ \eng{attended} (faux ami d'« assister à »).
\end{enumerate}
Correction : \eng{Because I agree that ICTs provide useful information, I attended the computer lesson with pleasure.} « Comme je reconnais que les TIC fournissent des informations utiles, j'ai suivi le cours d'informatique avec plaisir. »
\end{exoresolu}

\courssub{The assessment grid: how your essay is marked}

Voici la grille d'évaluation type, inspirée des pratiques du Baccalauréat camerounais : cinq critères notés chacun sur 4 points, pour un total de 20.

\begin{popfigure}
\begin{tikzpicture}[font=\small]
\fill[popblue] (0,0) rectangle (12.6,0.8);
\node[text=white, anchor=west, font=\bfseries] at (0.15,0.4) {Criterion};
\node[text=white, anchor=west, font=\bfseries] at (5.6,0.4) {What the examiner expects};
\node[text=white, anchor=east, font=\bfseries] at (12.45,0.4) {Mark};
\fill[popblueL] (0,-0.8) rectangle (12.6,0);
\node[anchor=west] at (0.15,-0.4) {\textbf{1. Task and length}};
\node[anchor=west, align=left] at (5.6,-0.4) {Answers the topic (difficulties of ICTs);\\ about 250--300 words};
\node[anchor=east, font=\bfseries] at (12.45,-0.4) {/ 4};
\fill[popcream] (0,-1.6) rectangle (12.6,-0.8);
\node[anchor=west] at (0.15,-1.2) {\textbf{2. Organisation}};
\node[anchor=west, align=left] at (5.6,-1.2) {Clear introduction, paragraphs, conclusion;\\ varied connectors};
\node[anchor=east, font=\bfseries] at (12.45,-1.2) {/ 4};
\fill[popblueL] (0,-2.4) rectangle (12.6,-1.6);
\node[anchor=west] at (0.15,-2.0) {\textbf{3. Vocabulary}};
\node[anchor=west, align=left] at (5.6,-2.0) {Rich, precise ICT vocabulary;\\ no repetition, no French words};
\node[anchor=east, font=\bfseries] at (12.45,-2.0) {/ 4};
\fill[popcream] (0,-3.2) rectangle (12.6,-2.4);
\node[anchor=west] at (0.15,-2.8) {\textbf{4. Grammar and spelling}};
\node[anchor=west, align=left] at (5.6,-2.8) {Correct tenses, agreements, articles;\\ accurate spelling};
\node[anchor=east, font=\bfseries] at (12.45,-2.8) {/ 4};
\fill[popblueL] (0,-4.0) rectangle (12.6,-3.2);
\node[anchor=west] at (0.15,-3.6) {\textbf{5. Originality and examples}};
\node[anchor=west, align=left] at (5.6,-3.6) {Personal or local examples (school,\\ cybercafé, power cuts); personal voice};
\node[anchor=east, font=\bfseries] at (12.45,-3.6) {/ 4};
\fill[popgold] (0,-4.8) rectangle (12.6,-4.0);
\node[anchor=west, font=\bfseries] at (0.15,-4.4) {TOTAL};
\node[anchor=east, font=\bfseries] at (12.45,-4.4) {/ 20};
\draw[popline] (0,0) rectangle (12.6,-4.8);
\draw[popline] (5.45,0) -- (5.45,-4.8);
\end{tikzpicture}
\legende{Grille d'évaluation type d'une composition au Baccalauréat : 5 critères $\times$ 4 points = 20 points.}
\end{popfigure}

\begin{aretenir}[Ce que la grille vous apprend]
Aucun critère ne vaut plus qu'un autre : un élève au vocabulaire riche mais désorganisé perd autant qu'un élève organisé mais plein de fautes. La stratégie gagnante est l'\cle{équilibre} : répondez au sujet, structurez en trois parties, utilisez le vocabulaire de la leçon, soignez les temps, et illustrez avec un exemple vécu (une coupure de courant pendant un cours en ligne à Douala, un cybercafé trop cher à Bamenda).
\end{aretenir}

\begin{savaistu}[La présentation compte aussi]
Écrire 250 mots lisibles vaut mieux que 400 mots pleins de ratures. Un correcteur qui lit des centaines de copies apprécie une écriture claire, des paragraphes aérés et peu de ratures. Si vous vous trompez, barrez proprement d'un trait et continuez. Une copie soignée et un peu courte obtient souvent une meilleure note qu'une copie longue et illisible : la qualité prime sur la quantité.
\end{savaistu}

\courssub{The C-O-L proofreading strategy}

Il ne suffit pas de bien écrire : il faut aussi bien \cle{relire}. Réservez toujours les dix dernières minutes à une relecture méthodique en trois passes, dans cet ordre précis.

\begin{methode}[La relecture C-O-L : Content, Organisation, Language]
\begin{enumerate}
\item \textbf{C — Content (le contenu)} : ai-je vraiment répondu au sujet ? Ai-je parlé des \emph{difficulties} des TIC et non de leurs avantages ? Chaque paragraphe apporte-t-il une idée liée au sujet ? Sinon, supprimez ou réorientez.
\item \textbf{O — Organisation} : y a-t-il une introduction avec problématique, des paragraphes bien délimités, une conclusion ? Ai-je utilisé des connecteurs variés (\emph{first, moreover, however, as a result, to conclude}) ? Ai-je évité \emph{because... so} et les répétitions de \emph{but} ?
\item \textbf{L — Language (la langue)} : les temps sont-ils cohérents (présent pour les faits généraux, prétérit pour les anecdotes) ? Les verbes s'accordent-ils (\emph{he has}, pas \emph{he have}) ? Les indénombrables sont-ils sans -s (\emph{information, advice, news}) ? L'orthographe des mots pièges est-elle correcte (\emph{which, because, technology, receive}) ?
\end{enumerate}
\end{methode}

\begin{exercice}[À vous de jouer]
Reprenez la copie faible du début de cette section. 1) Listez ses dix erreurs en les classant selon les trois passes C-O-L. 2) Réécrivez le deuxième paragraphe en anglais correct, avec des connecteurs variés. 3) Rédigez la conclusion qui manque, en deux ou trois phrases.
\end{exercice}

\begin{corrige}[Exercice — éléments de réponse]
1) \textbf{Content} : dernier paragraphe hors sujet (avantages au lieu de difficultés). \textbf{Organisation} : pas de conclusion ; répétition de \emph{But} ; \emph{Because... so} ensemble. \textbf{Language} : \emph{I am agree} $\rightarrow$ \emph{I agree} ; \emph{informations} $\rightarrow$ \emph{information} ; \emph{advices} $\rightarrow$ \emph{advice} ; \emph{assisted at} $\rightarrow$ \emph{attended} ; \emph{there is many problems} $\rightarrow$ \emph{there are many problems} ; \emph{informatics tools} $\rightarrow$ \emph{computer tools}.

2) Exemple de paragraphe réécrit : \eng{The connection is very slow, so students cannot download their lessons. Moreover, teachers give us a lot of information and advice, but we cannot always use them well. Last week, for instance, I attended a computer lesson and the electricity went off after ten minutes. However, the teacher continued, although we could not do the exercises. It was extremely frustrating.} « La connexion est très lente, donc les élèves ne peuvent pas télécharger leurs leçons. De plus, les professeurs nous donnent beaucoup d'informations et de conseils, mais nous ne pouvons pas toujours bien les utiliser. La semaine dernière, par exemple, j'ai suivi un cours d'informatique et le courant s'est coupé au bout de dix minutes. Cependant, le professeur a continué, bien que nous ne puissions pas faire les exercices. C'était extrêmement frustrant. »

3) Exemple de conclusion : \eng{To conclude, ICTs can help students learn, but slow connections, power cuts and lack of training make them difficult to use. Schools should therefore invest in better equipment and train both teachers and students.} « Pour conclure, les TIC peuvent aider les élèves à apprendre, mais les connexions lentes, les coupures de courant et le manque de formation les rendent difficiles à utiliser. Les écoles devraient donc investir dans de meilleurs équipements et former les enseignants comme les élèves. »
\end{corrige}

\begin{aretenir}[Bilan de la section]
Pour maximiser votre note : respectez le sujet (les \emph{difficulties}, pas les avantages), structurez en introduction-développement-conclusion, variez les connecteurs, méfiez-vous des indénombrables et des calques du français, et consacrez toujours dix minutes à la relecture C-O-L. La grille sur 20 points récompense l'équilibre : contenu, organisation, vocabulaire, grammaire et originalité comptent chacun pour un cinquième de la note.
\end{aretenir}

\newpage

\coursec{Activité d'intégration}

\coursec{Lesson 60 — Writing: Writes compositions on difficulties with the use of ICTs at work or for learning}

\courssub{Objectifs de l'activité}
\begin{itemize}
  \item Mobiliser l'ensemble des acquis de la leçon (structure de composition, connecteurs logiques, vocabulaire thématique, temps verbaux) pour produire un texte cohérent et argumenté.
  \item Rédiger une composition de \textbf{250 à 300 mots} sur les difficultés liées à l'utilisation des TIC pour l'apprentissage, dans un contexte camerounais réaliste.
  \item Appliquer une démarche d'écriture processuelle : planification, rédaction, relecture autonome (méthode C-O-L), correction croisée par les pairs à l'aide d'une grille d'évaluation.
  \item Développer l'autonomie et l'esprit critique face à sa propre production et à celle des autres.
\end{itemize}

\courssub{Situation}
\begin{textebox}{School Magazine Call for Articles — \eng{Government Bilingual High School Yaoundé, English Section}}
\eng{GBHS Yaoundé — English Section Magazine }\\\eng{Vol. 12, Issue 3 — March 2025}

\eng{\textbf{CALL FOR ARTICLES: “ICTs in Our Classrooms: The Real Struggles”}}

\eng{Dear Students,}

\eng{The Editorial Board of the \emph{GBHS Yaoundé English Section Magazine} invites you to contribute an article for our next special edition on \textbf{Technology and Education}. We want to hear \emph{your} voice — the daily reality of using Information and Communication Technologies (ICTs) for learning in our school and at home.}

\eng{Topic: \textbf{“The difficulties I face when using ICTs for learning”}}

\eng{Guidelines:}
\begin{itemize}
  \item Write a composition of \textbf{250–300 words}.
  \item Structure your text clearly: \emph{introduction, development (two or three paragraphs), conclusion}.
  \item Use appropriate linking words and topic vocabulary.
  \item Illustrate your points with concrete examples from your experience (power cuts, poor connectivity, lack of equipment, digital skills gaps, distractions, cost of data, etc.).
  \item Submit your draft to your English teacher by \textbf{Friday, 28 March 2025}.
\end{itemize}

\eng{The best three articles will be published in the magazine and displayed on the \emph{English Club} notice board. Certificates and small prizes (data bundles, USB keys) will be awarded during the \emph{English Day} ceremony.}

\eng{Let your pen (or keyboard) speak!}

\eng{\textit{The Editorial Team}}
\end{textebox}

\courssub{Consignes}
\begin{enumerate}
  \item \textbf{Analyse du sujet (5 min).} Lisez l'appel à contributions ci-dessus. Identifiez : le destinataire (lecteurs du magazine), le but (informer, témoigner, persuader), le ton (personnel mais structuré), la contrainte de longueur (250–300 mots) et le format (composition en trois parties).
  
  \item \textbf{Brainstorming et planification (10 min).} Sur une feuille de brouillon, notez toutes les difficultés qui vous viennent à l'esprit (ex. : coupures d'électricité, connexion lente/coûteuse, ordinateurs insuffisants au labo, manque de formation, distractions réseaux sociaux, fatigue oculaire, piratage de comptes, etc.). Sélectionnez les \textbf{3 ou 4 idées les plus fortes} et organisez-les en un plan détaillé :
  \begin{itemize}
    \item \textbf{Introduction} : accroche, définition rapide des TIC à l'école, annonce du plan.
    \item \textbf{Développement} : 2 ou 3 paragraphes, chacun centré sur une difficulté majeure (ex. : infrastructure, coût/accès, compétences/distractions). Chaque paragraphe = phrase sujet + exemples concrets + connecteurs.
    \item \textbf{Conclusion} : synthèse, ouverture (solutions, appel à l'action), phrase de clôture marquante.
  \end{itemize}
  
  \item \textbf{Rédaction (35 min).} Rédigez votre composition \textbf{au propre} sur la copie d'examen en respectant votre plan. Veillez à :
  \begin{itemize}
    \item Utiliser au moins \textbf{8 connecteurs logiques} variés (addition, contraste, cause, conséquence, illustration, conclusion).
    \item Intégrer au moins \textbf{10 mots/expressions du vocabulaire thématique} vus en leçon (\cle{connectivity issues}, \cle{power outage}, \cle{digital divide}, \cle{bandwidth}, \cle{computer literacy}, \cle{distraction}, \cle{eye strain}, \cle{cybersecurity}, \cle{affordability}, \cle{infrastructure}, etc.).
    \item Varier les temps verbaux de façon pertinente (present simple pour faits généraux, present perfect pour expérience vécue, past simple pour anecdotes datées, modaux pour recommandations).
    \item Soigner la ponctuation et la paragraphation (un saut de ligne par paragraphe).
  \end{itemize}
  
  \item \textbf{Relecture autonome — Méthode C-O-L (10 min).} Relisez votre texte trois fois, à chaque fois avec un focus différent :
  \begin{itemize}
    \item \textbf{C — Cohérence \& Cohésion} : Les idées s'enchaînent-elles logiquement ? Les connecteurs sont-ils bien choisis ? Les paragraphes sont-ils équilibrés ?
    \item \textbf{O — Orthographe \& Grammaire} : Accords (sujet/verbe, adjectifs), temps verbaux, prépositions, articles, faux amis (\cle{actually} vs \cle{actuellement}, \cle{eventually} vs \cle{éventuellement}), majuscules (\eng{I, Monday, Cameroon, English}).
    \item \textbf{L — Lexique \& Longueur} : Le vocabulaire est-il précis et varié ? Avez-vous évité les répétitions ? Comptez les mots (250–300). Ajustez si nécessaire.
  \end{itemize}
  
  \item \textbf{Correction croisée par les pairs (15 min).} Échangez votre copie avec votre binôme. Utilisez la \textbf{grille d'évaluation} fournie par l'enseignant (voir \emph{Ressources à mobiliser}) pour annoter le texte de votre camarade : points forts (cochez), points à améliorer (surlignez + code d'erreur), note sur 20. Rendez la copie annotée à son auteur.
  
  \item \textbf{Rédaction finale et dépôt (5 min).} Intégrez les corrections pertinentes, faites une dernière relecture rapide, puis remettez votre composition finale à l'enseignant. Les trois meilleures copies seront lues en classe et affichées au tableau du \emph{English Club}.
\end{enumerate}

\courssub{Ressources à mobiliser}
\begin{propriete}[Rappel — Structure type d'une composition]
\begin{center}
\begin{tabularx}{\linewidth}{|X|X|X|}
\hline
\rowcolor{popblueL} \textbf{Partie} & \textbf{Contenu} & \textbf{Expressions utiles} \\
\hline
Introduction & Accroche (fait, chiffre, question, citation) + définition du sujet + annonce du plan & \eng{Nowadays, ICTs have become... / It is undeniable that... / Many students face... / This essay will examine... / Firstly,... secondly,... finally...} \\
\hline
Développement (¶2–4) & 1 idée principale par paragraphe = phrase sujet + arguments + exemples concrets + connecteurs & \eng{To begin with, / A major difficulty is... / For instance, / In addition, / However, / As a result, / Moreover, / Another challenge lies in...} \\
\hline
Conclusion & Synthèse des points clés + ouverture (solution, recommandation, avis personnel) + phrase de clôture & \eng{To sum up, / In conclusion, / All things considered, / It is high time the authorities... / I strongly believe that... / Only then will we...} \\
\hline
\end{tabularx}
\end{center}
\end{propriete}

\begin{propriete}[Connecteurs logiques — catégories et exemples]
\begin{center}
\begin{tabularx}{\linewidth}{|l|X|}
\hline
\rowcolor{popblueL} \textbf{Fonction} & \textbf{Connecteurs (niveau Terminale)} \\
\hline
Addition / Enchaînement & \eng{furthermore, moreover, additionally, besides, what is more, also, too} \\
\hline
Contraste / Opposition & \eng{however, nevertheless, on the other hand, whereas, while, despite / in spite of, although / even though} \\
\hline
Cause / Raison & \eng{because, since, as, due to, owing to, the reason why... is that...} \\
\hline
Conséquence / Résultat & \eng{therefore, consequently, as a result, thus, hence, so} \\
\hline
Illustration / Exemple & \eng{for example, for instance, such as, namely, a case in point is...} \\
\hline
Conclusion / Résumé & \eng{to conclude, in conclusion, to sum up, all in all, ultimately} \\
\hline
Condition / Hypothèse & \eng{if, unless, provided that, as long as, suppose (that)} \\
\hline
\multicolumn{2}{|X|}{\emph{Attention :} Ne commencez \textbf{jamais} une phrase par \eng{and, but, so, because} à l'écrit formel. Utilisez \eng{In addition, However, Therefore, Since}.} \\
\hline
\end{tabularx}
\end{center}
\end{propriete}

\begin{definition}[Vocabulaire thématique — ICT difficulties (rappel)]
\begin{multicols}{2}
\begin{itemize}
  \item \cle{connectivity issues} — problèmes de connexion
  \item \cle{power outage / blackout} — coupure d'électricité
  \item \cle{unreliable internet} — internet instable
  \item \cle{low bandwidth} — faible bande passante
  \item \cle{digital divide} — fracture numérique
  \item \cle{computer literacy / digital skills} — compétences numériques
  \item \cle{affordability / cost of data} — coût des données / accessibilité financière
  \item \cle{insufficient equipment} — équipement insuffisant
  \item \cle{outdated hardware / software} — matériel / logiciels obsolètes
  \item \cle{distraction / procrastination} — distraction / procrastination
  \item \cle{eye strain / screen fatigue} — fatigue oculaire
  \item \cle{cybersecurity threats / hacking} — menaces de cybersécurité / piratage
  \item \cle{technical glitch / bug} — incident technique / bogue
  \item \cle{maintenance / technical support} — maintenance / support technique
  \item \cle{e-learning / online platform} — apprentissage en ligne / plateforme
\end{itemize}
\end{multicols}
\end{definition}

\begin{propriete}[Temps verbaux recommandés pour ce sujet]
\begin{center}
\begin{tabular}{|l|l|l|}
\hline
\rowcolor{popblueL} \textbf{Temps} & \textbf{Emploi} & \textbf{Exemples} \\
\hline
Present Simple & Faits généraux, habitudes, vérités permanentes & \eng{Many students \textbf{lack} access to computers. / Power cuts \textbf{occur} daily.} \\
\hline
Present Perfect & Expérience vécue (non datée), conséquence présente & \eng{I \textbf{have faced} connectivity issues. / The school \textbf{has not provided} enough laptops.} \\
\hline
Past Simple & Anecdotes datées, actions terminées dans le passé & \eng{Last week, the network \textbf{crashed} during the exam. / Yesterday, I \textbf{could not} submit my assignment.} \\
\hline
Present Continuous & Situation en cours, évolution & \eng{Internet costs \textbf{are rising}. / More students \textbf{are using} smartphones for learning.} \\
\hline
Modaux (\eng{should, must, have to, could, might}) & Recommandations, obligation, possibilité, hypothèse & \eng{The government \textbf{should invest} in infrastructure. / We \textbf{must} learn digital skills. / Better connectivity \textbf{could} improve results.} \\
\hline
Passive voice & Focus sur l'objet / impersonnalisation & \eng{Computers \textbf{are shared} among five students. / Data bundles \textbf{cannot be afforded} by all.} \\
\hline
\end{tabular}
\end{center}
\end{propriete}

\begin{methode}[Méthode C-O-L — Relecture efficace en 3 passages]
\begin{enumerate}
  \item \textbf{C — Cohérence \& Cohésion} : Lisez \emph{à voix haute} (murmurez). Vérifiez : phrase d'accroche ? thèse claire ? plan annoncé ? chaque paragraphe = 1 idée ? connecteurs variés et justes ? transitions fluides ? conclusion = synthèse + ouverture ?
  \item \textbf{O — Orthographe \& Grammaire} : Lisez \emph{à l'envers} (dernière phrase → première) pour couper le sens et voir les formes. Traquez : accords S/V (\eng{he/she/it + s}), \eng{a/an/the}, prépositions (\eng{good \textbf{at}, interested \textbf{in}, depend \textbf{on}}), faux amis, majuscules (\eng{I, Cameroon, English, Monday}), ponctuation (pas de virgule avant \eng{that} en relatif).
  \item \textbf{L — Lexique \& Longueur} : Surveillez les répétitions (remplacez par synonymes/pronoms). Vérifiez le vocabulaire technique (au moins 10 items). Comptez les mots : 250–300. Si < 250 : développez un exemple. Si > 300 : coupez les phrases redondantes.
\end{enumerate}
\end{methode}

\begin{propriete}[Grille d'évaluation — Composition ICT difficulties (sur 20)]
\begin{center}
\begin{tabularx}{\linewidth}{|X|c|X|}
\hline
\rowcolor{popblueL} \textbf{Critère} & \textbf{Points} & \textbf{Indicateurs de réussite} \\
\hline
\cle{Task achievement} (respect du sujet, longueur) & 4 & Sujet traité pleinement, 250–300 mots, ton approprié (magazine scolaire), exemples personnels concrets \\
\hline
\cle{Organisation \& Cohérence} (structure, paragraphes, connecteurs) & 4 & Intro/développement/conclusion clairs, 3–4 paragraphes, connecteurs variés (≥8), progression logique \\
\hline
\cle{Vocabulaire} (thématique, variété, précision) & 4 & ≥10 mots du thème ICT, registres variés, pas de répétitions lourdes, collocations correctes (\eng{face difficulties, bridge the gap, raise awareness}) \\
\hline
\cle{Grammaire \& Syntaxe} (temps, structures, exactitude) & 4 & Temps verbaux justifiés, phrases complexes (relatives, subordonnées), voix passive, modaux, peu d'erreurs bloquantes \\
\hline
\cle{Orthographe \& Ponctuation} (lisibilité, conventions) & 4 & Orthographe soignée, majuscules correctes, ponctuation anglaise (point, virgule, apostrophe), paragraphation visible \\
\hline
\textbf{Total} & \textbf{20} & \\
\hline
\end{tabularx}
\end{center}
\end{propriete}

\courssub{Proposition de production et corrigé commenté}
\begin{exoresolu}{Modèle de composition — \eng{“The difficulties I face when using ICTs for learning”} (≈ 280 mots)}
\eng{Nowadays, Information and Communication Technologies have become indispensable tools in education. In Cameroon, the government has launched several programmes to integrate ICTs into schools, yet the daily reality for many students remains far from the digital ideal. This essay examines the main obstacles I encounter when using ICTs for learning, focusing on infrastructure gaps, financial barriers and the lack of digital skills.}

\eng{To begin with, unreliable infrastructure is a major hindrance. Frequent power outages disrupt online classes and make it impossible to charge devices. At Government Bilingual High School Yaoundé, the computer laboratory has only twenty functional desktops for over two hundred students in the English section. Consequently, we share terminals in groups of five, which limits hands-on practice. Moreover, the internet connection is painfully slow; low bandwidth means that downloading a simple PDF can take half an hour.}

\eng{In addition, the cost of data creates a digital divide among learners. Many of my classmates cannot afford daily bundles to access platforms like Google Classroom or Khan Academy. Since most educational resources are now online, students from low-income families fall behind. I have personally missed deadlines because my data expired mid-research. This affordability issue undermines equal opportunities.}

\eng{Another challenge lies in insufficient computer literacy. Although we are called “digital natives”, many peers struggle with basic tasks: formatting a document, citing sources, or protecting their accounts from hacking. Distractions are also rampant; social media notifications constantly tempt us to procrastinate. As a result, screen time often yields more eye strain than academic progress.}

\eng{To sum up, ICTs hold great promise for Cameroonian education, but infrastructure deficits, high data costs and skills gaps turn that promise into daily frustration. It is high time the authorities invested in stable electricity, subsidised student data plans and compulsory digital literacy training. Only then will technology truly empower every learner, regardless of background.}
\tcblower
\textbf{Commentaire des choix de langue (pour l'enseignant / auto-évaluation) :}
\begin{itemize}
  \item \textbf{Structure} : Introduction (accroche + contexte + thèse + plan), 3 paragraphes de développement (infrastructure, coût, compétences/distractions), conclusion (synthèse + recommandations + phrase forte). Paragraphation visible.
  \item \textbf{Connecteurs} (12 utilisés) : \eng{yet, To begin with, Moreover, Consequently, In addition, Since, This affordability issue, Another challenge lies in, As a result, To sum up, but, Only then}. Variété : addition, conséquence, contraste, illustration, conclusion, inversion (\eng{Only then will...}).
  \item \textbf{Vocabulaire thématique} (14 items) : \eng{indispensable, integrate, digital ideal, infrastructure gaps, financial barriers, power outages, functional, hands-on practice, bandwidth, digital divide, affordability, equal opportunities, computer literacy, hacking, procrastinate, eye strain, empower, subsidised, compulsory}. Collocations naturelles : \eng{fall behind, miss deadlines, hold promise, turn into, invest in}.
  \item \textbf{Grammaire / Temps} : Present simple (\eng{have become, remains, disrupt, make, has, share, limits, means, creates, undermines, struggle, tempt, yields}), Present perfect (\eng{has launched, have personally missed}), Past simple (\eng{expired}), Modaux (\eng{should invest — implicite dans "It is high time... invested", must learn, could improve}), Passive (\eng{are called, are now online, is undermined}), Subjonctif mandatif après \eng{It is high time} (\eng{invested}), Inversion après \eng{Only then} (\eng{will technology truly empower}).
  \item \textbf{Contexte camerounais} : Références concrètes (GBHS Yaoundé, English section, 200 élèves, 20 ordinateurs, Google Classroom, Khan Academy, forfaits data, coupures d'électricité) — authenticité et ancrage local.
  \item \textbf{Registre} : Formel mais accessible (magazine scolaire), pas de familier, pas de contractions (\eng{cannot, do not, it is}), phrases complètes.
  \item \textbf{Longueur} : 282 mots (dans la fourchette 250–300).
\end{itemize}
\end{exoresolu}

\courssub{Bilan de l'activité}
\begin{aretenir}[Ce que vous devez retenir de cette activité d'intégration]
\begin{itemize}
  \item \textbf{L'écriture est un processus} : planifier $\rightarrow$ rédiger $\rightarrow$ relire (C-O-L) $\rightarrow$ corriger $\rightarrow$ finaliser. Sauter une étape fragilise le résultat final.
  \item \textbf{La structure fait la note} : une composition sans introduction claire, sans paragraphes thématiques ou sans conclusion synthétique ne peut pas obtenir la moyenne, même avec un bon vocabulaire.
  \item \textbf{Les connecteurs sont le « ciment » du texte} : ils guident le lecteur, montrent la logique de votre pensée et prouvent votre maîtrise de la cohésion (critère noté sur 4).
  \item \textbf{Le vocabulaire thématique doit être précis} : évitez les termes vagues (\eng{problem, thing, bad}) au profit de mots techniques (\eng{connectivity issues, bandwidth, digital divide, affordability}). C'est le marqueur du niveau Terminale.
  \item \textbf{La relecture C-O-L change tout} : 10 minutes concentrées sur Cohérence, Orthographe/Grammaire, Lexique/Longueur éliminent 80\% des erreurs évitables.
  \item \textbf{La correction par les pairs est un apprentissage double} : en corrigeant l'autre, vous aiguisez votre propre regard critique ; en acceptant les remarques, vous progressez.
  \item \textbf{Le contexte camerounais est votre force} : les correcteurs valorisent les exemples ancrés dans votre réalité (GBHS Yaoundé, \eng{bendskin} pour aller au cyber, forfait MTN/Orange, coupures ENEO, labo surchargé). Soyez concret, pas abstrait.
\end{itemize}
\end{aretenir}

\begin{experience}[Prolongement — English Club : « ICT Solutions Hackathon »]
En binôme ou trinôme, imaginez \textbf{trois solutions concrètes} pour réduire les difficultés identifiées dans vos compositions (ex. : “Solar-powered charging station in the school yard”, “Negotiated student data bundle with MTN/Orange”, “Peer-to-peer digital literacy workshops every Friday”). Présentez votre projet sous forme d'un \textbf{pitch de 2 minutes} (oral) + une \textbf{affiche A4} (visuel + arguments clés en anglais). Les meilleurs projets seront transmis à l'administration et au club des parents d'élèves (APEE). Cette activité valide les compétences \cle{Speaking} (presentation) et \cle{Writing} (poster), et renforce l'engagement citoyen.
\end{experience}

\newpage

\coursec{Méthodes et exercices résolus}

\coursec{Lesson 60 -- Writing: Writes compositions on difficulties with the use of ICTs at work or for learning}

\courssub{Understanding the Task and Brainstorming Ideas}

\begin{methode}[Analyser le sujet et faire un remue-méninges efficace]
    \textbf{Objectif :} Comprendre exactement ce qui est demandé et générer des idées pertinentes avant d'écrire.
    \begin{enumerate}
        \item \textbf{Lisez le sujet deux fois.} Soulignez les \cle{mots-clés} : le thème (ICTs), le contexte (work \textbf{ou} learning), l'angle (difficulties \textbf{/} challenges), le format (composition / essay).
        \item \textbf{Identifiez le type de composition.} S'agit-il d'un \eng{for and against essay} (avantages/inconvénients), d'un \eng{problem-solution essay} (problèmes/solutions) ou d'un \eng{opinion essay} ? Ici, le sujet impose souvent une approche \textbf{problèmes / solutions} ou \textbf{causes / effets}.
        \item \textbf{Faites un « brain dump » (vidage de cerveau) pendant 3--4 minutes.} Notez \textbf{tout} ce qui vous vient à l'esprit en vrac (mots, expressions, exemples camerounais : coupures d'électricité, coût des données, formation insuffisante, cybercriminalité, distraction des réseaux sociaux). Ne vous censurez pas.
        \item \textbf{Classez vos idées en catégories.} Regroupez-les en 2 ou 3 grands axes logiques pour former les paragraphes du développement.
            \begin{itemize}
                \item \textit{Axe 1 : Obstacles techniques et infrastructurels} (connectivité, électricité, matériel vétuste).
                \item \textit{Axe 2 : Obstacles humains et financiers} (coût, illettrisme numérique, résistance au changement, cybersécurité).
                \item \textit{Axe 3 : Pédagogiques / Professionnels} (distraction, surcharge d'informations, isolation, adaptation des programmes).
            \end{itemize}
        \item \textbf{Sélectionnez les 2 meilleures idées par axe} avec un exemple concret (statistique inventée plausible, anecdote, fait d'actualité) pour les illustrer.
    \end{enumerate}
\end{methode}

\begin{exoresolu}[Brainstorming et plan détaillé pour un sujet type Bac]
    \textbf{Sujet :} \eng{``Write a composition of about 250--300 words on the difficulties students face when using ICTs for learning in Cameroon. Suggest possible solutions.''}
    \tcblower
    \textbf{Étape 1 : Analyse du sujet}
    \begin{itemize}
        \item \textbf{Theme :} ICTs in education (ICTs pour l'apprentissage).
        \item \textbf{Target population :} Students (élèves/étudiants).
        \item \textbf{Context :} Cameroon (contexte camerounais indispensable).
        \item \textbf{Task :} 1. Discuss difficulties (développement 1). 2. Suggest solutions (développement 2). Format : Problem-Solution Essay.
    \end{itemize}

    \textbf{Étape 2 : Brainstorming (extrait)}
    \begin{center}
    \begin{tabularx}{\linewidth}{|X|X|}\hline
    \rowcolor{popblueL} \textbf{Difficulties (Problèmes)} & \textbf{Solutions} \\ \hline
    \textbf{Infrastructure :} Frequent power cuts (load shedding), poor internet connectivity in rural areas (Bertoua, Maroua vs Yaoundé), high cost of data bundles (MTN, Orange, Nexttel). & Solar panels in schools, government subsidies for student data plans, offline digital libraries (Kolibri, Kiwix). \\ \hline
    \textbf{Human / Skills :} Low digital literacy (teachers \& students), lack of maintenance culture, technophobia among older teachers. & Compulsory ICT training in ENIEG/ENS, peer-to-peer student clubs (Google Developer Student Clubs), maintenance budgets. \\ \hline
    \textbf{Pedagogical / Social :} Distraction (Facebook, TikTok, WhatsApp during class), plagiarism (copy-paste), cyberbullying, eye strain. & Digital citizenship curriculum, firewalls/MDM on school devices, awareness campaigns, ergonomic breaks. \\ \hline
    \end{tabularx}
    \end{center}

    \textbf{Étape 3 : Plan détaillé (Outline)}
    \begin{itemize}
        \item \textbf{Introduction :} Hook (COVID-19 a accéléré l'adoption), Context (Cameroon's digital ambition), Thesis: \eng{``While ICTs offer immense educational opportunities, Cameroonian students grapple with infrastructural, financial and pedagogical hurdles that require urgent multi-stakeholder solutions.''}
        \item \textbf{Body Paragraph 1 (Infrastructural \& Financial):} Topic sentence: \eng{``The primary barrier remains inadequate infrastructure.''} Details: Electricity, connectivity, cost of devices/data. Example: \eng{``A student in Bamenda may walk kilometres to find a charged laptop and a stable signal.''} Connector: \eng{Moreover / Furthermore}.
        \item \textbf{Body Paragraph 2 (Human \& Pedagogical):} Topic sentence: \eng{``Beyond hardware, the human factor poses significant challenges.''} Details: Digital literacy gap, distraction, plagiarism. Example: \eng{``Many teachers use the blackboard only because they lack training.''} Connector: \eng{In addition / Besides}.
        \item \textbf{Body Paragraph 3 (Solutions):} Topic sentence: \eng{``Addressing these issues demands a holistic approach.''} Solutions: Govt policy (subsidies, curriculum), School level (training, offline resources), Individual (discipline, time management). Connector: \eng{Consequently / To remedy this}.
        \item \textbf{Conclusion :} Restate thesis (other words), Summary of main points, Final thought / Call to action: \eng{``Bridging the digital divide is not a luxury but a necessity for Cameroon's emergence by 2035.''}
    \end{itemize}
    \textbf{Conclusion de l'exercice :} Ce plan assure une progression logique (Problèmes $\rightarrow$ Solutions) et ancre le propos dans la réalité camerounaise (exemples locaux), ce qui est valorisé à l'examen.
\end{exoresolu}

\courssub{Structuring the Composition: Introduction, Body, Conclusion}

\begin{methode}[Construire une architecture solide en trois parties]
    \textbf{Objectif :} Produire un texte cohérent, bien paragraphe, respectant les normes de la composition académique anglophone.
    \begin{enumerate}
        \item \textbf{L'Introduction (environ 10\% du nombre de mots).} Elle suit l'entonnoir :
            \begin{itemize}
                \item \textit{Hook (Accroche) :} Une phrase choc, une citation, un fait marquant, une question rhétorique. \eng{``In 2020, the COVID-19 pandemic turned every kitchen table in Yaoundé into a potential classroom.''}
                \item \textit{Background (Contexte) :} 2--3 phrases pour situer le débat (politique nationale, enjeu mondial).
                \item \textit{Thesis Statement (Phrase thèse) :} \textbf{La phrase la plus importante.} Elle annonce le plan (souvent 2 ou 3 points). \eng{``This essay examines the infrastructural and pedagogical obstacles hindering e-learning in Cameroon and proposes concrete remedies.''}
            \end{itemize}
        \item \textbf{Le Développement (Body Paragraphs -- 80\%).} \textbf{Un paragraphe = Une idée principale.} Structure \textbf{PEEL} :
            \begin{itemize}
                \item \textbf{P}oint (Topic Sentence) : Annonce l'idée du paragraphe.
                \item \textbf{E}xplanation : Développe l'idée (pourquoi ? comment ?).
                \item \textbf{E}vidence / Example : Donne un fait, un chiffre, un exemple camerounais.
                \item \textbf{L}ink : Relie au paragraphe suivant ou à la thèse.
            \end{itemize}
            \textit{Ordre logique :} Du plus général au plus spécifique, ou Problèmes puis Solutions, ou Causes puis Effets.
        \item \textbf{La Conclusion (environ 10\%).} Elle suit l'entonnoir inverse :
            \begin{itemize}
                \item \textit{Restatement :} Reformulez la thèse (synonymes, structure différente). \textbf{Never copy-paste.}
                \item \textit{Summary :} Résumez les 2--3 arguments principaux en une phrase chacun.
                \item \textit{Final Thought / Call to Action :} Ouvrez sur l'avenir, une recommandation forte, une citation. \eng{``Empowering the youth with reliable digital tools is investing in Cameroon's future.''}
            \end{itemize}
        \item \textbf{Paragraphing visuel :} Sautez une ligne claire entre chaque paragraphe. Indentez (rentrez) la première ligne de chaque paragraphe (sauf si consigne contraire).
    \end{enumerate}
\end{methode}

\begin{exoresolu}[Reconnaissance de structure et rédaction d'une introduction]
    \textbf{Exercice A :} Reconstituez l'ordre logique de ces phrases pour former une introduction correcte.
    \begin{enumerate}
        \item \eng{Therefore, understanding these barriers is crucial for policy makers.}
        \item \eng{The integration of Information and Communication Technologies in education has become a global imperative.}
        \item \eng{This essay will analyse the major infrastructural and human obstacles faced by learners in Buea and propose sustainable solutions.}
        \item \eng{However, in many Cameroonian secondary schools, this integration remains more theoretical than practical.}
    \end{enumerate}
    \textbf{Exercice B :} Rédigez une introduction complète (Hook, Background, Thesis) pour le sujet : \eng{``The use of smartphones in class: a distraction or a learning tool?''}
    \tcblower
    \textbf{Solution Exercice A :} Ordre correct : \textbf{2 -- 4 -- 1 -- 3}.
    \begin{itemize}
        \item \eng{2} = Hook / General Context (Global imperative).
        \item \eng{4} = Background / Contrast (Cameroon reality vs theory).
        \item \eng{1} = Link / Importance (Why write this essay?).
        \item \eng{3} = Thesis Statement (Announces plan: obstacles in Buea + solutions).
    \end{itemize}

    \textbf{Solution Exercice B : Proposition de rédaction}
    \begin{textebox}[Model Introduction]
    \eng{``A recent survey by the Ministry of Secondary Education revealed that over 70\% of students in Douala own a smartphone. (Hook: Chiffre choc / Contexte local). While these devices are primarily seen as entertainment gadgets, their potential as portable libraries and research tools is undeniable. (Background: Débat). This composition argues that, despite the risks of distraction and cyberbullying, smartphones can become powerful educational allies if integrated within a strict pedagogical framework. (Thesis: Position nuancée + Plan implicite: risques vs avantages + condition).''}
    \end{textebox}
    \textbf{Analyse :} L'introduction est courte (~55 mots), ancre le sujet au Cameroun (Douala, MINESEC), présente les deux faces (distraction vs outil) et prend position claire (thesis) qui guide le développement.
\end{exoresolu}

\courssub{Mastering Connectors, Vocabulary and Register}

\begin{methode}[Utiliser les connecteurs logiques et le lexique thématique avec précision]
    \textbf{Objectif :} Assurer la fluidité (cohérence) et montrer la maîtrise lexicale (ressource lexicale), deux critères majeurs de notation.
    \begin{enumerate}
        \item \textbf{Maîtrisez les familles de connecteurs (Linking Words).} Ne les saupoudrez pas ; utilisez-les pour \textbf{structurer la pensée}.
            \begin{center}
            \begin{tabularx}{\linewidth}{|l|X|X|}\hline
            \rowcolor{popblueL} \textbf{Fonction} & \textbf{Connecteurs (Variété Bac+)} & \textbf{Position / Ponctuation} \\ \hline
            \textbf{Addition / Listing} & \eng{Furthermore, Moreover, In addition, Additionally, Not only ... but also} & Début de phrase + virgule. \\ \hline
            \textbf{Contrast / Concession} & \eng{However, Nevertheless, On the one hand ... on the other hand, Although / Even though + clause, Despite / In spite of + N/V-ing} & \eng{However,} (début). \eng{Although} (début ou milieu, pas de virgule avant). \eng{Despite} + nom/V-ing. \\ \hline
            \textbf{Cause / Reason} & \eng{Because / Since / As + clause, Due to / Owing to / Because of + N/V-ing, The reason why ... is that ...} & \eng{Due to} (souvent après \eng{be}). Évitez \eng{Because} en tout début de phrase formelle (préférez \eng{Since/As}). \\ \hline
            \textbf{Consequence / Result} & \eng{Therefore, Consequently, As a result, Thus, Hence, So (trop oral)} & Début de phrase + virgule. \\ \hline
            \textbf{Example / Illustration} & \eng{For instance, For example, Such as, A case in point is ..., To illustrate} & \eng{For instance,} (début). \eng{Such as} (milieu, pas de virgule avant). \\ \hline
            \textbf{Conclusion} & \eng{In conclusion, To sum up, Ultimately, All things considered} & Début de phrase + virgule. \\ \hline
            \end{tabularx}
            \end{center}
        \item \textbf{Enrichissez votre « Vocabulaire ICT / Difficultés ».} Remplacez les mots simples (\eng{problem, bad, difficult, use}) par du vocabulaire précis (\cle{Academic Vocabulary}).
            \begin{center}
            \begin{tabularx}{\linewidth}{|X|X|X|}\hline
            \rowcolor{popblueL} \textbf{Mot simple (Éviter)} & \textbf{Vocabulaire ciblé (Cibler)} & \textbf{Collocations utiles} \\ \hline
            \eng{Problem / Difficulty} & \eng{Hurdle, obstacle, barrier, bottleneck, setback, challenge, drawback} & \eng{Overcome a hurdle, address a bottleneck} \\ \hline
            \eng{Lack / Not enough} & \eng{Shortage, deficit, scarcity, insufficiency, dearth} & \eng{A shortage of devices, a dearth of skilled trainers} \\ \hline
            \eng{Internet / Connection} & \eng{Connectivity, bandwidth, broadband, network coverage, latency} & \eng{Poor connectivity, limited bandwidth, unstable network} \\ \hline
            \eng{Electricity / Power cut} & \eng{Power outage, load shedding, unreliable power supply, grid instability} & \eng{Frequent power outages disrupt online classes} \\ \hline
            \eng{Skill / Knowledge} & \eng{Digital literacy, proficiency, competence, expertise, tech-savvy (adj)} & \eng{Low digital literacy, enhance proficiency} \\ \hline
            \eng{Expensive / Cost} & \eng{Prohibitive cost, financial burden, affordability issue, exorbitant data prices} & \eng{The prohibitive cost of data bundles} \\ \hline
            \eng{Distraction} & \eng{Diversion, loss of focus, multitasking (souvent négatif ici), attention span} & \eng{Shortened attention span, constant notifications} \\ \hline
            \end{tabularx}
            \end{center}
        \item \textbf{Adoptez le registre formel (Formal Register).}
            \begin{itemize}
                \item Pas de contractions : \eng{do not} (pas \eng{don't}), \eng{it is} (pas \eng{it's}).
                \item Préférez le passif ou l'impersonnel : \eng{It is argued that...}, \eng{Students are faced with...} (vs \eng{Students face...}).
                \item Verbes précis : \eng{hamper, hinder, impede} (plutôt que \eng{stop/slow}), \eng{facilitate, enable, empower} (plutôt que \eng{help}), \eng{mitigate, alleviate} (plutôt que \eng{reduce/solve}).
            \end{itemize}
    \end{enumerate}
\end{methode}

\begin{exoresolu}[Transformation de registre et insertion de connecteurs]
    \textbf{Exercice A (Registre) :} Réécrivez ces phrases en style formel académique.
    \begin{enumerate}
        \item \eng{Students can't learn well because the net is bad.}
        \item \eng{Teachers don't know how to use computers, so it's a big problem.}
        \item \eng{Buying data is too expensive for most parents.}
    \end{enumerate}
    \textbf{Exercice B (Connecteurs) :} Complétez le paragraphe avec les connecteurs appropriés parmi : \eng{However, For instance, Furthermore, Consequently, In conclusion}.
    \eng{``Many schools in rural Cameroon lack computers. \_\_\_\_\_, even when hardware is donated, there is often no electricity to run it. \_\_\_\_\_, a school in the East Region received 20 laptops last year but they remain unused due to load shedding. \_\_\_\_\_, the digital divide widens between urban and rural learners. \_\_\_\_\_, urgent investment in solar energy is needed.''}
    \tcblower
    \textbf{Solution Exercice A :}
    \begin{enumerate}
        \item \eng{Students \textbf{are unable to} learn effectively \textbf{due to poor / unreliable internet connectivity}.} (\eng{cannot} $\rightarrow$ \eng{are unable to}; \eng{bad} $\rightarrow$ \eng{poor/unreliable}; \eng{because} $\rightarrow$ \textbf{due to} + nom).
        \item \eng{Teachers' \textbf{lack of digital proficiency} \textbf{constitutes a major obstacle / poses a significant challenge}.} (\eng{don't know how to use} $\rightarrow$ \eng{lack proficiency}; \eng{so it's a big problem} $\rightarrow$ verbe fort \eng{constitutes/poses} + nom précis).
        \item \eng{The \textbf{prohibitive cost of data bundles} \textbf{places a heavy financial burden on} most households.} (\eng{too expensive} $\rightarrow$ \eng{prohibitive cost}; \eng{buying} $\rightarrow$ \eng{cost}; structure causative \eng{places a burden on}).
    \end{enumerate}

    \textbf{Solution Exercice B :}
    \eng{``Many schools in rural Cameroon lack computers. \textbf{Furthermore}, even when hardware is donated, there is often no electricity to run it. \textbf{For instance}, a school in the East Region received 20 laptops last year but they remain unused due to load shedding. \textbf{Consequently}, the digital divide widens between urban and rural learners. \textbf{In conclusion}, urgent investment in solar energy is needed.''}
    \textbf{Justifications :} \eng{Furthermore} (ajout d'argument), \eng{For instance} (exemple précis), \eng{Consequently} (conséquence logique), \eng{In conclusion} (synthèse/ouverture solution).
\end{exoresolu}

\courssub{Grammar Focus: Verb Tenses, Modals and Passive Voice}

\begin{methode}[Choisir les bons temps, modaux et voix pour l'essai argumentatif]
    \textbf{Objectif :} Éviter les erreurs de concordance des temps et utiliser la grammaire pour nuancer l'argumentation (modalité, objectivité).
    \begin{enumerate}
        \item \textbf{Present Simple (Le roi de l'essai).} Pour les faits généraux, vérités actuelles, situations permanentes.
            \eng{``Many students \textbf{lack} access to reliable electricity.''} \eng{``The curriculum \textbf{does not integrate} coding yet.''}
        \item \textbf{Present Perfect (Le lien passé-présent).} Pour des évolutions récentes, des expériences vécues, des statistiques récentes (\eng{recently, since 2020, over the past decade}).
            \eng{``Internet penetration \textbf{has increased} significantly \textbf{since} the launch of the National ICT Strategic Plan.''} \eng{``The government \textbf{has invested} in fibre optics.''}
        \item \textbf{Modaux de recommandation / obligation / possibilité (Crucial pour les solutions).}
            \begin{center}
            \begin{tabular}{|l|l|l|}\hline
            \rowcolor{popblueL} \textbf{Modal} & \textbf{Nuance / Force} & \textbf{Exemple solution} \\ \hline
            \eng{Must / Have to} & Obligation forte, impératif moral/légal & \eng{The State \textbf{must} subsidise data for students.} \\ \hline
            \eng{Should / Ought to} & Conseil fort, recommandation (standard) & \eng{Schools \textbf{should} integrate digital literacy modules.} \\ \hline
            \eng{Could / Might} & Possibilité, suggestion moins directe & \eng{Partnerships with telcos \textbf{could} lower costs.} \\ \hline
            \eng{Need to} & Nécessité objective & \eng{Teachers \textbf{need to} undergo continuous training.} \\ \hline
            \end{tabular}
            \end{center}
            \textbf{Attention :} \eng{Should} n'est pas suivi de \eng{to} (\eng{should go}, pas \eng{should to go}). \eng{Must} n'a pas de passé (\eng{had to}).
        \item \textbf{Passive Voice (Voix passive) pour l'objectivité.} Déplace le focus sur l'action/objet, pas sur l'acteur (souvent inconnu ou « they »).
            \eng{Active: ``The government should build labs.'' $\rightarrow$ Passive: ``Labs \textbf{should be built}.''}
            \eng{Active: ``They waste a lot of time.'' $\rightarrow$ Passive: ``A lot of time \textbf{is wasted}.''}
            \textit{Structure :} Sujet + \eng{be} (temps requis) + \textbf{Participe Passé (V3)} + (by agent).
        \item \textbf{Conditionnelles (Si clauses) pour argumenter.}
            \eng{First Conditional (Réaliste) :} \eng{If the government \textbf{invests} in solar power, schools \textbf{will have} stable electricity.}
            \eng{Second Conditional (Hypothétique/Irréel présent) :} \eng{If students \textbf{had} tablets, they \textbf{would access} offline libraries.}
    \end{enumerate}
\end{methode}

\begin{exoresolu}[Correction grammaticale et transformation voix passive / modaux]
    \textbf{Exercice A :} Mettez les verbes entre parenthèses au temps/forme corrects.
    \eng{``Over the last five years, the number of internet users in Cameroon (rise) \_\_\_\_\_ sharply. However, the quality of service (not improve) \_\_\_\_\_ at the same pace. If the regulatory body (enforce) \_\_\_\_\_ stricter standards, consumers (benefit) \_\_\_\_\_. It is high time the authorities (take) \_\_\_\_\_ action.''}
    \textbf{Exercice B :} Transformez à la voix passive (conservez le modal/temps).
    \begin{enumerate}
        \item \eng{The ministry must allocate more funds to ICT infrastructure.}
        \item \eng{They are currently installing fibre optic cables in Buea.}
        \item \eng{Students should not use phones during exams.}
    \end{enumerate}
    \textbf{Exercice C :} Corrigez les erreurs de modaux/temps dans ce paragraphe.
    \eng{``To solve the problem, the government must to provide free wifi. Teachers should to be trained. If we will have good internet, we can learn better. The computers are not used since two years.''}
    \tcblower
    \textbf{Solution Exercice A :}
    \eng{``Over the last five years, the number of internet users in Cameroon \textbf{has risen} sharply. However, the quality of service \textbf{has not improved} at the same pace. If the regulatory body \textbf{enforces} stricter standards, consumers \textbf{will benefit}. It is high time the authorities \textbf{took} action.''}
    \begin{itemize}
        \item \eng{has risen} : Present Perfect (évolution sur période \eng{over the last five years}). Verbe irrégulier \eng{rise - rose - risen}.
        \item \eng{has not improved} : Present Perfect (négatif, période non finie).
        \item \eng{enforces / will benefit} : \textbf{First Conditional} (If + Present Simple, Will + BV). Réalisme.
        \item \eng{took} : \eng{It is high time} + \textbf{Prétérit modal (Subjonctif)}. Pas \eng{take} ni \eng{have taken}.
    \end{itemize}

    \textbf{Solution Exercice B (Passive) :}
    \begin{enumerate}
        \item \eng{More funds \textbf{must be allocated} to ICT infrastructure (by the ministry).}
        \item \eng{Fibre optic cables \textbf{are currently being installed} in Buea.} (Present Continuous Passive : \eng{be being} + V3).
        \item \eng{Phones \textbf{should not be used} during exams (by students).} (Modal Passive : \eng{should be} + V3).
    \end{enumerate}

    \textbf{Solution Exercice C (Correction) :}
    \eng{``To solve the problem, the government \textbf{must provide} (pas \eng{to}) free wifi. Teachers \textbf{should be trained} (pas \eng{to be}, passif modal). If we \textbf{have} (pas \eng{will have}, First Conditional) good internet, we \textbf{can / will be able to} learn better. The computers \textbf{have not been used for} (pas \eng{since} + durée, Present Perfect Passive) two years.''}
    \textbf{Rappel :} \eng{Since} + point de départ (date) ; \eng{For} + durée.
\end{exoresolu}

\courssub{Proofreading: The Final Polish (Self-Correction Checklist)}

\begin{methode}[Relire et corriger efficacement en 5 minutes chrono]
    \textbf{Objectif :} Gagner les points faciles (orthographe, accord, ponctuation, mots oubliés) qui font la différence entre 14 et 17.
    \begin{enumerate}
        \item \textbf{Lisez à l'envers (phrase par phrase, de la fin au début).} Cela casse le sens global et force l'œil à voir la \textbf{forme} (orthographe, grammaire) et non le fond.
        \item \textbf{Check-list « COPS » (Capitals, Organisation, Punctuation, Spelling/Grammar) :}
            \begin{itemize}
                \item \textbf{C -- Capitals :} Majuscules en début de phrase, \textbf{I} (pronom), noms propres (\eng{Cameroon, Yaoundé, MTN, MINESEC, English, Monday}), adjectifs de nationalité (\eng{Cameroonian, French}).
                \item \textbf{O -- Organisation :} Paragraphes distincts (saut de ligne) ? Introduction / Corps / Conclusion présents ? Connecteurs variés ? Thesis statement respectée ?
                \item \textbf{P -- Punctuation :} Points (pas de \eng{...} à la fin), virgules après connecteurs initiaux (\eng{However,}), pas de virgule avant \eng{that} dans les relatives, apostrophes (\eng{students' needs, it's} $\rightarrow$ interdire \eng{it's}).
                \item \textbf{S -- Spelling \& Grammar (Mots pièges) :}
                    \begin{itemize}
                        \item \textbf{Accords :} \eng{One of the \textbf{main problems} is...} (pluriel après \eng{one of the}), \eng{The number of students \textbf{is}...} (singulier), \eng{A number of students \textbf{are}...} (pluriel).
                        \item \textbf{Ordre des mots :} \eng{Explain \textbf{me} the rule} $\rightarrow$ \eng{Explain the rule \textbf{to me}}. \eng{I \textbf{am agree}} $\rightarrow$ \eng{I \textbf{agree}}.
                        \item \textbf{Prépositions :} \eng{Good \textbf{at} using}, \eng{Interested \textbf{in}}, \eng{Depend \textbf{on}}, \eng{Access \textbf{to}} (pas \eng{access to internet} $\rightarrow$ \eng{access \textbf{to} the internet}), \eng{Consist \textbf{of}}.
                        \item \textbf{Orthographe « Bac » :} \eng{Accommodation} (2c, 2m), \eng{Necessary} (1c, 2s), \eng{Receive} (i avant e sauf après c), \eng{Separate} (a rat), \eng{Definitely} (finite), \eng{Government} (n), \eng{Environment}, \eng{Privilege}.
                    \end{itemize}
            \end{itemize}
        \item \textbf{Comptez les mots (approximatif).} 250--300 mots = environ 1.5 à 2 pages manuscrites (écriture moyenne). Si trop court : développez un exemple. Si trop long : coupez les répétitions.
    \end{enumerate}
\end{methode}

\begin{exoresolu}[Atelier de relecture : Chasse aux erreurs]
    \textbf{Consigne :} Le paragraphe ci-dessous contient \textbf{12 erreurs} (orthographe, grammaire, vocabulaire, ponctuation, registre). Identifiez-les, corrigez-les et réécrivez la version finale.
    \eng{``In Cameroon, the use of ICTs at school is a big challenge. First, their is a lack of equipments in many establishments. Teachers are not formed to use digital tools, so they don't use it. Furthermore the cost of internet is to high for parents. Students uses phones for facebook and games instead of learning. However, solutions exist. The government should to provide tablets. If we had electricity, we will study better. To conclude, ICTs are necessaries for the developpement of our country.''}
    \tcblower
    \textbf{Analyse et Corrections (12 erreurs + améliorations style) :}
    \begin{enumerate}
        \item \eng{their} $\rightarrow$ \textbf{there} (existence). \textit{Homophone.}
        \item \eng{equipments} $\rightarrow$ \textbf{equipment} (indénombrable, pas de s). \textit{Erreur de compte/incountable.}
        \item \eng{formed} $\rightarrow$ \textbf{trained} (faux ami : \eng{formed} = façonné ; \eng{trained} = formé). \textit{Faux ami.}
        \item \eng{they don't use it} $\rightarrow$ \textbf{they do not use them} (pluriel \eng{tools} ; pas de contraction). \textit{Accord + Registre.}
        \item \eng{Furthermore the cost} $\rightarrow$ \textbf{Furthermore, the cost} (virgule après connecteur initial). \textit{Ponctuation.}
        \item \eng{to high} $\rightarrow$ \textbf{too high} (\eng{too} = trop ; \eng{to} = préposition). \textit{Homophone / Grammaire.}
        \item \eng{Students uses} $\rightarrow$ \textbf{Students use} (sujet pluriel). \textit{Accord S 3e pers.}
        \item \eng{facebook} $\rightarrow$ \textbf{Facebook} (majuscule nom propre). \textit{Majuscule.}
        \item \eng{should to provide} $\rightarrow$ \textbf{should provide} (modal sans \eng{to}). \textit{Grammaire modale.}
        \item \eng{If we had ... we will} $\rightarrow$ \textbf{If we had ... we would} (Second Conditional : hypothétique). \textit{Concordance temps.}
        \item \eng{necessaries} $\rightarrow$ \textbf{necessary} (adjectif invariable). \textit{Accord adjectif.}
        \item \eng{developpement} $\rightarrow$ \textbf{development} (orthographe anglaise : \eng{p}, pas \eng{pp}). \textit{Orthographe.}
    \end{enumerate}
    \textbf{Version corrigée et améliorée (Registre formel) :}
    \begin{textebox}[Corrected Version]
    \eng{``In Cameroon, the integration of ICTs in education remains a major challenge. Firstly, there is a \textbf{shortage of equipment} in many institutions. Teachers \textbf{lack adequate training} in digital tools; \textbf{consequently}, they \textbf{rarely integrate them} into their lessons. \textbf{Moreover}, the \textbf{prohibitive cost of internet connectivity} \textbf{places a heavy burden on} parents. \textbf{Instead of leveraging} smartphones for research, many students \textbf{use them for social media and entertainment}. \textbf{Nevertheless}, viable solutions exist. The government \textbf{must subsidise} devices and data for learners. \textbf{Should} stable electricity \textbf{be provided}, students \textbf{would benefit} from uninterrupted access. \textbf{In conclusion}, ICTs are \textbf{indispensable for} Cameroon's \textbf{development}.''}
    \end{textebox}
    \textbf{Gains :} Vocabulaire précis (\eng{shortage, prohibitive cost, leverage, subsidise}), grammaire irréprochable (passif modal \eng{should be provided}, inversion conditionnelle \eng{Should ... be}), connecteurs variés, registre formel.
\end{exoresolu}

\courssub{Time Management and Exam Strategy (Simulation)}

\begin{methode}[Gérer les 2 heures de l'épreuve de Writing au Bac]
    \textbf{Objectif :} Terminer une composition propre, complète et relue dans le temps imparti.
    \begin{enumerate}
        \item \textbf{0h00 -- 0h15 (15 min) : Analyse \& Planification.}
            \begin{itemize}
                \item Lire le sujet, souligner mots-clés, choisir l'angle (si choix).
                \item Brainstorming 5 min (face cachée de la copie).
                \item Structurer : Thesis Statement + 2--3 Topic Sentences + 1 idée de conclusion. Écrire le plan au brouillon.
            \end{itemize}
        \item \textbf{0h15 -- 1h15 (60 min) : Rédaction (Drafting).}
            \begin{itemize}
                \item Écrire \textbf{sur la copie d'examen} (pas de recopie finale, pas le temps).
                \item Suivre le plan \textbf{strictement}. Ne pas chercher le mot parfait : écrire, laisser une case vide \eng{[...]} si blocage, continuer.
                \item Viser 280 mots $\pm$ 10\%.
            \end{itemize}
        \item \textbf{1h15 -- 1h45 (30 min) : Révision \& Correction (Le moment clé).}
            \begin{itemize}
                \item Relire une 1ère fois pour le \textbf{sens et la structure} : Les paragraphes suivent-ils le plan ? Les connecteurs sont-ils logiques ? L'argumentation est-elle claire ?
                \item Relire une 2ème fois pour la \textbf{forme (COPS)} : Orthographe, accords, temps, prépositions, majuscules, ponctuation. Utiliser la méthode « lecture à l'envers ».
                \item Compter les mots (lignes $\times$ mots/ligne).
            \end{itemize}
        \item \textbf{1h45 -- 2h00 (15 min) : Marge de sécurité / Propreté.}
            \begin{itemize}
                \item Ratures propres (trait unique), écriture lisible.
                \item Vérifier que le nom/numéro est sur la copie.
            \end{itemize}
    \end{enumerate}
    \textbf{Astuce :} Si vous bloquez sur un mot anglais, écrivez-le en français entre crochets \eng{[...]} et revenez-y à la relecture. Ne restez pas planté 5 min sur un mot.
\end{methode}

\begin{exoresolu}[Simulation chronométrée : Sujet complet]
    \textbf{Sujet d'examen blanc (2h) :}
    \eng{``Remote work has become common since the pandemic. Write a composition discussing the difficulties employees face when working from home in Cameroon and suggest ways to improve their productivity.''}
    \tcblower
    \textbf{Simulation guidée (Ce que l'élève doit produire au brouillon en 15 min) :}

    \textbf{1. Mots-clés :} Remote work / Work from home (WFH) / Since pandemic / Difficulties / Employees / Cameroon / Improve productivity.

    \textbf{2. Brainstorming (Catégorisé) :}
    \begin{itemize}
        \item \textbf{Difficultés :}
            \begin{itemize}
                \item \textit{Infra :} Coupures courant (ENEO), Internet instable/cher, espace de travail inadapté (bruit, famille, promiscuité).
                \item \textit{Humain/Org :} Isolement, difficulté séparation vie pro/perso, management à distance (micromanagement ou abandon), manque matériel (imprimante, scanner), procrastination.
            \end{itemize}
        \item \textbf{Solutions :}
            \begin{itemize}
                \item \textit{Employeur :} Fournir routeurs 4G de secours, data illimitée, chèque équipement bureau (chaise, écran), formation gestion temps/outils collaboratifs (Teams, Slack, Trello).
                \item \textit{État :} Incitations fiscales entreprises équipant télétravail, amélioration réseau national.
                \item \textit{Employé :} Routine stricte, espace dédié, technique Pomodoro, communication proactive.
            \end{itemize}
    \end{itemize}

    \textbf{3. Thesis Statement :}
    \eng{``Although remote work offers flexibility, Cameroonian employees struggle with infrastructural deficits, inadequate home environments, and managerial gaps, requiring coordinated efforts from the state, employers, and workers themselves to boost productivity.''}

    \textbf{4. Topic Sentences (Plan) :}
    \begin{enumerate}
        \item \eng{Infrastructural hurdles, notably unstable electricity and costly internet, severely hamper remote work efficiency.}
        \item \eng{The unsuitability of home environments and the blurring of professional boundaries further erode focus and well-being.}
        \item \eng{To mitigate these issues, a tripartite approach involving state regulation, employer support, and self-discipline is essential.}
    \end{enumerate}

    \textbf{Conseil final :} Entraînez-vous à faire ce plan en 12 minutes chrono. Le jour J, le plan est votre assurance-vie contre le hors-sujet et la page blanche.
\end{exoresolu}

\begin{attention}[Les 5 erreurs qui coûtent des points au Bac]
    \textbf{1. L'orthographe « tueuse » (Lexical Spelling) -- \eng{0,5 à 1 pt} par erreur récurrente.}
    \begin{itemize}
        \item \eng{\textbf{Accommodation} (2c, 2m) $\neq$ accomodation.}
        \item \eng{\textbf{Necessary} (1c, 2s) $\neq$ neccessary / neccesary.}
        \item \eng{\textbf{Receive} (i avant e après c) $\neq$ recieve.}
        \item \eng{\textbf{Separate} (a rat inside) $\neq$ seperate.}
        \item \eng{\textbf{Government} (n muet) $\neq$ gouvernment.}
        \item \eng{\textbf{Development} (1p) $\neq$ developpement (fr).}
        \item \eng{\textbf{Privilege} (pas d après g) $\neq$ privilège.}
        \item \eng{\textbf{Environment} (n avant m) $\neq$ environement.}
        \item \eng{\textbf{Occurrence} (2c, 2r) $\neq$ occurence / ocurrence.}
        \item \eng{\textbf{Committee} (2m, 2t, 2e) $\neq$ comitee / commitee.}
        \item \eng{\textbf{Recommend} (1c, 2m) $\neq$ recomend / reccommend.}
        \item \eng{\textbf{Referred} (2r) $\neq$ refered.}
    \end{itemize}
    \textit{Astuce :} Apprenez par cœur cette liste des « \cle{Dirty Dozen} » (la douzaine sale) avant l'examen.

    \textbf{2. Les Faux Amis (False Friends) -- Perte de sens + point grammaire/vocab.}
    \begin{center}
    \begin{tabular}{|l|l|l|}\hline
    \rowcolor{poporangeL} \textbf{Mot anglais} & \textbf{Piège français} & \textbf{Traduction correcte / Alternative} \\ \hline
    \eng{Actually} & Actuellement & \eng{Currently / At present} (\eng{Actually} = en fait, en réalité) \\ \hline
    \eng{Eventually} & Éventuellement & \eng{Finally / In the end} (\eng{Eventually} = finalement, après un moment) \\ \hline
    \eng{Sensible} & Sensible & \eng{Reasonable / Sensitive} (\eng{Sensible} = raisonnable ; \eng{Sensitive} = sensible) \\ \hline
    \eng{Formed} & Formé (éduqué) & \eng{Trained / Educated} (\eng{Formed} = constitué, façonné) \\ \hline
    \eng{Important} (pour quantité) & Important (grand) & \eng{Significant / Substantial / Considerable} \\ \hline
    \eng{Opportunity} & Opportunité & \eng{Chance / Possibility} (\eng{Opportunity} = occasion favorable) \\ \hline
    \end{tabular}
    \end{center}

    \textbf{3. Les Calques Syntaxiques du Français (Word Order / Structure) -- \eng{0,5 pt} à chaque fois.}
    \begin{itemize}
        \item \eng{\textcolor{red}{I am agree}} $\rightarrow$ \eng{\textcolor{popgreen}{I agree}} (Verbe d'état, pas \eng{be} + adj).
        \item \eng{\textcolor{red}{I am waiting you}} $\rightarrow$ \eng{\textcolor{popgreen}{I am waiting \textbf{for} you}} (Verbe + préposition obligatoire).
        \item \eng{\textcolor{red}{Explain me the rule}} $\rightarrow$ \eng{\textcolor{popgreen}{Explain the rule \textbf{to me}}} (Ordre V + OD + \eng{to} + OI).
        \item \eng{\textcolor{red}{The difficulties of the use}} $\rightarrow$ \eng{\textcolor{popgreen}{The difficulties \textbf{in using} / \textbf{of using}}} (Gérondif après préposition / nom).
        \item \eng{\textcolor{red}{Permit / Allow to do}} $\rightarrow$ \eng{\textcolor{popgreen}{Allow \textbf{someone} to do / Permit \textbf{someone} to do}} (Objet obligatoire).
        \item \eng{\textcolor{red}{Help me to do}} $\rightarrow$ \eng{\textcolor{popgreen}{Help me (to) do}} (\eng{to} optionnel, souvent omis).
    \end{itemize}

    \textbf{4. Les Temps Verbaux : Present Perfect vs Past Simple / Conditionnelles -- \eng{1 pt} par erreur de sens.}
    \begin{itemize}
        \item \eng{\textcolor{red}{I have seen him yesterday.}} $\rightarrow$ \eng{\textcolor{popgreen}{I \textbf{saw} him yesterday.}} (Marqueur de temps passé $\rightarrow$ Past Simple).
        \item \eng{\textcolor{red}{Since two years, I learn English.}} $\rightarrow$ \eng{\textcolor{popgreen}{I \textbf{have been learning} English \textbf{for} two years.}} (\eng{Since} + date, \eng{For} + durée ; Present Perfect Continuous pour durée action continue).
        \item \eng{\textcolor{red}{If I will have time, I will do it.}} $\rightarrow$ \eng{\textcolor{popgreen}{If I \textbf{have} time, I \textbf{will do} it.}} (Jamais \eng{will} après \eng{if} dans conditionnelle 1).
        \item \eng{\textcolor{red}{If I would know, I would tell you.}} $\rightarrow$ \eng{\textcolor{popgreen}{If I \textbf{knew}, I \textbf{would tell} you.}} (Conditionnelle 2 : \eng{if} + Prétérit, \eng{would} + BV).
    \end{itemize}

    \textbf{5. Prépositions et Articles (Les « petits mots » qui fâchent) -- \eng{0,25 pt} mais s'accumulent.}
    \begin{itemize}
        \item \eng{Good \textbf{in} English} $\rightarrow$ \eng{Good \textbf{at} English / \textbf{in} class}.
        \item \eng{Interested \textbf{by}} $\rightarrow$ \eng{Interested \textbf{in}}.
        \item \eng{Difficulties \textbf{to} do} $\rightarrow$ \eng{Difficulties \textbf{in doing} / \textbf{to do} (après adjectif \eng{difficult}, mais \eng{difficulty} + \eng{in V-ing})}.
        \item \eng{Access \textbf{internet}} $\rightarrow$ \eng{Access \textbf{to} the internet} (\eng{Access} verbe transitif direct : \eng{access the internet} ; nom : \eng{access to}).
        \item \eng{Depend \textbf{of}} $\rightarrow$ \eng{Depend \textbf{on}}.
        \item \eng{\textbf{The} ICTs} $\rightarrow$ \eng{\textbf{$\emptyset$} ICTs} (Pas de \eng{the} devant sigle pluriel général, sauf spécifié : \eng{The ICTs \textbf{used in our school}}).
        \item \eng{\textbf{A} equipment / \textbf{An} information} $\rightarrow$ \eng{\textbf{$\emptyset$} equipment / \textbf{$\emptyset$} information} (Indénombrables = pas d'article indéfini). Utilisez \eng{a piece of equipment / an item of information}.
    \end{itemize}
    \textbf{Conseil final :} Relisez \textbf{spécifiquement} ces 5 points lors de votre dernière relecture (1h45--2h00). C'est là que se joue la mention.
\end{attention}

\newpage

\coursec{Exercices}

\coursec{Lesson 60 — Writing: Writes compositions on difficulties with the use of ICTs at work or for learning}

\courssub{Understanding the Task and Brainstorming}

\begin{methode}[Analyser le sujet et faire un brainstorming]
\begin{enumerate}
    \item \textbf{Identify the keywords:} \eng{difficulties}, \eng{ICTs} (Information and Communication Technologies), \eng{school} / \eng{workplace}, \eng{solutions}.
    \item \textbf{Determine the format:} A \eng{composition} (essai argumentatif) — ton formel, pas de contractions, structure \eng{Introduction / Body / Conclusion}.
    \item \textbf{Brainstorm ideas (Mind-map):}
    \begin{itemize}
        \item \textbf{Infrastructure:} poor connectivity, outdated hardware, power outages, lack of maintenance.
        \item \textbf{Human factor:} insufficient training, resistance to change, weak cybersecurity awareness, digital divide.
        \item \textbf{Solutions:} government investment (fibre optic), public-private partnerships, continuous teacher training, preventive maintenance policies, cybersecurity protocols.
    \end{itemize}
    \item \textbf{Select and group:} Keep 2-3 strong main points for body paragraphs. Order them logically (e.g. problems $\rightarrow$ causes $\rightarrow$ solutions).
\end{enumerate}
\end{methode}

\courssub{The Structure of a Composition}

\begin{propriete}[Structure standard d'une composition (Essay)]
Une composition formelle comporte trois parties obligatoires~:
\begin{center}
\begin{tabularx}{\linewidth}{|X|X|X|}
\hline
\rowcolor{popblueL}
\textbf{Part} & \textbf{Content} & \textbf{Key Elements} \\
\hline
\textbf{Introduction} ($\approx$ 10\%) & \eng{Hook} (accroche), \eng{Background} (contexte), \textbf{Thesis Statement} (énoncé de thèse + plan) & \eng{In today's digital age...}, \eng{However...}, \eng{This essay examines... and suggests...} \\
\hline
\textbf{Body Paragraphs} ($\approx$ 80\%) & 2 à 3 paragraphes distincts. Chaque paragraphe = 1 idée principale. & \textbf{Topic Sentence}, \eng{Explanation}, \eng{Evidence/Example} (\eng{For instance}), \eng{Linkers} \\
\hline
\textbf{Conclusion} ($\approx$ 10\%) & \eng{Restatement} de la thèse (mots différents), \eng{Summary} des points clés, \eng{Final thought / Recommendation} (ouverture) & \eng{To sum up...}, \eng{Ultimately...}, \eng{It is crucial that...} \\
\hline
\end{tabularx}
\end{center}
\end{propriete}

\begin{exemplebox}[Outline exemple : « Difficulties of using ICTs in Cameroonian schools »]
\textbf{Introduction:} Hook on digital transformation. Context: Cameroonian schools struggle. Thesis: \eng{This essay discusses connectivity and equipment gaps, lack of training, and proposes investment and maintenance as solutions.}

\textbf{Body 1 (Infrastructure):} Topic Sentence: \eng{First, poor infrastructure hinders ICT integration.} Details: weak bandwidth, old computers, power cuts. Connector: \eng{For instance}, rural schools rely on mobile data. Result: \eng{Consequently}, e-learning impossible.

\textbf{Body 2 (Human Factor):} Topic Sentence: \eng{Furthermore, the human factor exacerbates technical issues.} Details: untrained teachers, no IT support, weak passwords. Example: \eng{A teacher in Buea cannot use the new platform.}

\textbf{Body 3 (Solutions):} Topic Sentence: \eng{To address these challenges, concrete measures are needed.} Solutions: Gov investment in fibre, equipment renewal, mandatory continuous training, maintenance contracts.

\textbf{Conclusion:} \eng{To sum up}, while ICTs offer opportunities, the digital divide is real. \eng{Ultimately}, political will and funding are essential to modernise education.
\end{exemplebox}

\courssub{Useful Vocabulary and Connectors}

\begin{aretenir}[Vocabulaire thématique : ICT Difficulties (Niveau Terminale)]
\begin{center}
\begin{tabularx}{\linewidth}{|l|X|X|}
\hline
\rowcolor{popblueL}
\textbf{Category} & \textbf{English Term} & \textbf{French Translation / Context} \\
\hline
\textbf{Connectivity} & \eng{bandwidth} & bande passante (débit max) \\
& \eng{connectivity issues} & problèmes de connexion \\
& \eng{power outage / power cut} & coupure de courant \\
& \eng{unstable / unreliable connection} & connexion instable / peu fiable \\
& \eng{ADSL / fibre optic} & ADSL / fibre optique \\
& \eng{network congestion} & congestion du réseau \\
\hline
\textbf{Hardware / Software} & \eng{outdated / obsolete equipment} & matériel obsolète / dépassé \\
& \eng{legacy system} & système hérité (ancien) \\
& \eng{glitch / bug} & bogue, dysfonctionnement mineur \\
& \eng{crash / freeze} & plantage / gel de l'écran \\
& \eng{hardware failure} & panne matérielle \\
& \eng{maintenance / upkeep} & maintenance / entretien \\
& \eng{student-computer ratio} & ratio élèves/ordinateur \\
\hline
\textbf{Cybersecurity} & \eng{cybersecurity} & cybersécurité \\
& \eng{phishing} & hameçonnage \\
& \eng{ransomware} & rançongiciel \\
& \eng{malware} & logiciel malveillant \\
& \eng{firewall} & pare-feu \\
& \eng{encryption / backup} & chiffrement / sauvegarde \\
& \eng{weak / strong password} & mot de passe faible / robuste \\
& \eng{data breach} & violation de données \\
\hline
\textbf{Human / Admin} & \eng{digital divide} & fracture numérique \\
& \eng{digital literacy} & culture / littératie numérique \\
& \eng{continuous training / upskilling} & formation continue / perfectionnement \\
& \eng{resistance to change} & résistance au changement \\
& \eng{budget constraints} & contraintes budgétaires \\
& \eng{stakeholder} & partie prenante \\
\hline
\end{tabularx}
\end{center}
\end{aretenir}

\begin{propriete}[Connecteurs logiques (Linkers) — Classification par fonction]
\begin{center}
\begin{tabularx}{\linewidth}{|l|X|}
\hline
\rowcolor{popblueL}
\textbf{Function} & \textbf{Connectors (Formal)} \\
\hline
\textbf{Addition} & \eng{Furthermore}, \eng{Moreover}, \eng{In addition}, \eng{Besides}, \eng{What is more} \\
\hline
\textbf{Contrast / Concession} & \eng{However}, \eng{Nevertheless}, \eng{On the other hand}, \eng{Conversely}, \eng{Although / Even though}, \eng{Despite / In spite of} (+ V-ing / noun) \\
\hline
\textbf{Cause / Reason} & \eng{Because / Since / As}, \eng{Due to / Owing to} (+ noun), \eng{The reason why... is that...} \\
\hline
\textbf{Effect / Result} & \eng{Therefore}, \eng{Consequently}, \eng{As a result}, \eng{Thus}, \eng{Hence}, \eng{So} (moins formel) \\
\hline
\textbf{Example} & \eng{For instance}, \eng{For example}, \eng{Such as}, \eng{Namely}, \eng{A case in point is...} \\
\hline
\textbf{Sequencing} & \eng{First / Firstly}, \eng{Second / Secondly}, \eng{Finally / Lastly}, \eng{Subsequently} \\
\hline
\textbf{Conclusion} & \eng{To sum up}, \eng{In conclusion}, \eng{Ultimately}, \eng{All things considered}, \eng{In a nutshell} (semi-formel) \\
\hline
\end{tabularx}
\end{center}
\end{propriete}

\begin{attention}[Pièges fréquents : Connecteurs]
\begin{itemize}
    \item \textbf{Despite / In spite of} + \textbf{Noun / V-ing} (jamais de clause avec sujet + verbe). \eng{Despite the rain, we worked.} / \eng{Despite being tired...} $\neq$ \eng{Despite it rained}.
    \item \textbf{However} = adverbe (point-virgule ou nouveau paragraphe, virgule après). \eng{The net is slow; however, we work.} $\neq$ \eng{The net is slow, however we work.} (Run-on sentence).
    \item \textbf{Because of / Due to} + Noun. \textbf{Because / Since / As} + Clause. \eng{Due to the strike, we stopped.} / \eng{Because the workers struck, we stopped.}
    \item \textbf{So that} = but (finalité) + \eng{can/could/will/would}. \eng{So} = conséquence. \eng{It was slow, so we waited.}
\end{itemize}
\end{attention}

\courssub{Verb Tenses and Grammar for this Essay}

\begin{propriete}[Temps verbaux et structures clés pour la composition]
\begin{center}
\begin{tabularx}{\linewidth}{|l|l|X|}
\hline
\rowcolor{popblueL}
\textbf{Structure} & \textbf{Usage} & \textbf{Exemples (traduits)} \\
\hline
\textbf{Present Simple} & Faits généraux, vérités permanentes, situation actuelle. & \eng{Slow internet \textbf{affects} productivity.} (Une connexion lente affecte la productivité.) \\
\hline
\textbf{Present Perfect} & Passé récent avec conséquence présente, expérience, évolution. & \eng{The government \textbf{has promised} fibre optic.} (Le gvt a promis la fibre.) \\
\hline
\textbf{Passive Voice} & Objectivité, focus sur l'objet, agent inconnu/secondaire. & \eng{Data \textbf{is stored} on servers.} (Les données sont stockées...) \\
&  & \eng{Computers \textbf{were donated} last year.} (Des ordis ont été donnés l'an dernier.) \\
\hline
\textbf{1st Conditional} & Situation réelle possible dans le futur (If + Present, Will/Can/Must). & \eng{If the school \textbf{invests} in training, teachers \textbf{will improve}.} (Si l'école investit...) \\
\hline
\textbf{2nd Conditional} & Hypothèse irréelle/peu probable au présent/futur (If + Prétérit, Would/Could/Might). & \eng{If we \textbf{had} better bandwidth, we \textbf{could} stream videos.} (Si nous avions...) \\
\hline
\textbf{3rd Conditional} & Regret / Hypothèse passée irréelle (If + Past Perfect, Would Have + PP). & \eng{If they \textbf{had backed up} data, they \textbf{would not have lost} it.} (S'ils avaient sauvegardé...) \\
\hline
\textbf{Mixed Conditional} & Cause passée, conséquence présente (If + Past Perfect, Would + Inf). & \eng{If the ministry \textbf{had invested} earlier, we \textbf{would not face} this crisis today.} \\
\hline
\textbf{Modals (Obligation/Rec)} & \eng{Must} (obligation forte), \eng{Should / Ought to} (conseil), \eng{Have to} (obligation externe). & \eng{Authorities \textbf{must allocate} funds.} / \eng{Students \textbf{should learn} cybersecurity.} \\
\hline
\textbf{Subjunctive (Mandative)} & Après \eng{It is essential/vital/crucial/imperative that} + Sujet + \textbf{Base Verb} (ou \eng{should} + BV). & \eng{It is essential that every teacher \textbf{receive} training.} (Subjonctif) \\
\hline
\end{tabularx}
\end{center}
\end{propriete}

\begin{attention}[Erreurs grammaticales types des francophones (Top 10)]
\begin{enumerate}
    \item \textbf{I am agree} $\rightarrow$ \eng{I agree} (verbe, pas adjectif).
    \item \textbf{Despite of / Despite... but} $\rightarrow$ \eng{Despite the cost} / \eng{Although it is expensive}.
    \item \textbf{Equipments / Informations / Advices} $\rightarrow$ \eng{Equipment / Information / Advice} (inombrables, pas de \eng{s}).
    \item \textbf{Suggest me to do} $\rightarrow$ \eng{Suggest doing} / \eng{Suggest that I do} (subjonctif: \eng{suggest I do}).
    \item \textbf{Responsible to} $\rightarrow$ \eng{Responsible for} + V-ing / Noun.
    \item \textbf{Enough powerful / Good enough} $\rightarrow$ Adj + \eng{enough} (\eng{Powerful enough}).
    \item \textbf{If I will have time} $\rightarrow$ \eng{If I have time} (Present après \eng{If} pour futur).
    \item \textbf{It depends of} $\rightarrow$ \eng{It depends on}.
    \item \textbf{He has been wrote / The report has been wrote} $\rightarrow$ \eng{Has been written} (Participe Passé irrégulier).
    \item \textbf{Faux-amis :} \eng{Actually} = en fait (pas actuellement $\rightarrow$ \eng{Currently}). \eng{Eventually} = finalement (pas éventuellement $\rightarrow$ \eng{Possibly}). \eng{Sensible} = raisonnable (pas sensible $\rightarrow$ \eng{Sensitive}).
\end{enumerate}
\end{attention}

\courssub{A Commented Model Essay}

\begin{exemplebox}[Modèle commenté : « The difficulties of using ICTs in Cameroonian schools » ($\approx$ 260 mots)]
\textbf{Title: Bridging the Digital Divide in Cameroonian Schools}

\textbf{Introduction:} In today's digital age, Information and Communication Technologies have become indispensable tools for quality education. \textbf{However}, in many Cameroonian schools, the integration of ICTs remains a daily struggle rather than an opportunity. \textbf{This essay examines the major infrastructure and human resource challenges hindering ICT use and proposes investment and training as sustainable solutions.}

\textbf{Body 1: Infrastructure Deficits.} \textbf{First}, poor connectivity is the most visible obstacle. Many schools, especially in rural areas, rely on weak mobile data or have no internet access at all. \textbf{For instance}, a teacher in the Far North region may wait twenty minutes to load a single educational video. \textbf{Consequently}, online research and interactive platforms are inaccessible during class hours. \textbf{Moreover}, the hardware is often outdated. Computer labs frequently contain decade-old machines that cannot run modern software, and the student-computer ratio can reach 10:1, making individual practice impossible. \textbf{Additionally}, frequent power outages damage equipment and interrupt lessons, as few schools possess automatic generators or UPS systems.

\textbf{Body 2: The Human Factor.} \textbf{Furthermore}, technical gaps are worsened by a lack of human capacity. Most teachers have not received adequate pre-service or in-service training on digital pedagogy. \textbf{As a result}, they lack the confidence to integrate ICTs into their lesson plans. Cybersecurity awareness is also low; weak passwords and phishing vulnerability put school data at risk. The absence of on-site IT technicians means that minor glitches become permanent breakdowns.

\textbf{Body 3: Towards Solutions.} \textbf{To address these issues}, a multi-pronged approach is required. The government must prioritise investment in fibre-optic infrastructure and renewable energy for schools. \textbf{Specifically}, public-private partnerships could help renew the computer fleet at lower costs. \textbf{Equally important}, the Ministry of Secondary Education should institutionalise mandatory, continuous digital training for all teachers and deploy at least one IT maintenance officer per school cluster.

\textbf{Conclusion:} \textbf{To sum up}, while ICTs hold transformative potential for Cameroonian education, the current digital divide is deepened by poor infrastructure, inadequate equipment, and insufficient training. \textbf{Ultimately}, bridging this gap demands not only financial resources but also strong political will and strategic planning. Only then can Cameroonian students fully benefit from the digital revolution.
\end{exemplebox}

\begin{aretenir}[Analyse du modèle (Points clés à reproduire)]
\begin{itemize}
    \item \textbf{Structure claire :} Thesis statement en fin d'intro (souligné dans le texte ci-dessus). Chaque paragraphe de développement commence par un \textbf{Topic Sentence} (gras dans le modèle) et un connecteur de séquence (\eng{First}, \eng{Furthermore}, \eng{To address}).
    \item \textbf{Connecteurs variés :} \eng{However}, \eng{For instance}, \eng{Consequently}, \eng{Moreover}, \eng{Additionally}, \eng{Furthermore}, \eng{As a result}, \eng{Specifically}, \eng{Equally important}, \eng{To sum up}, \eng{Ultimately}.
    \item \textbf{Grammaire ciblée :} Passif (\eng{are inaccessible}, \eng{are required}, \eng{be deployed}), Conditionnelles (\eng{If... were... would}, implicite dans \eng{To address}), Subjonctif (\eng{should institutionalise}, \eng{should receive}), Modaux (\eng{must prioritise}, \eng{could help}).
    \item \textbf{Vocabulaire précis :} \eng{digital divide}, \eng{outdated}, \eng{student-computer ratio}, \eng{power outages}, \eng{digital pedagogy}, \eng{cybersecurity awareness}, \eng{public-private partnerships}, \eng{multi-pronged approach}.
    \item \textbf{Style formel :} Pas de contractions (\eng{cannot}, \eng{do not}, \eng{it is}), pas de \eng{I think / In my opinion} (sauf conclusion), phrases complexes.
\end{itemize}
\end{aretenir}

\courssub{Je vérifie mes connaissances}

\begin{exercice}[QCM : Structure et connecteurs]
Choisissez la bonne réponse parmi les trois propositions.
\begin{enumerate}
    \item Where does the \eng{thesis statement} usually appear in a composition introduction?
    \begin{enumerate}
        \item At the very beginning (hook).
        \item At the end of the introduction.
        \item In the first body paragraph.
    \end{enumerate}
    \item Which connector introduces a \eng{contrast}?
    \begin{enumerate}
        \item Furthermore
        \item However
        \item Consequently
    \end{enumerate}
    \item What tense is most appropriate to describe \eng{general facts} about ICT difficulties (e.g. \eng{Slow internet affects productivity})?
    \begin{enumerate}
        \item Present Simple
        \item Present Continuous
        \item Past Simple
    \end{enumerate}
    \item Which of the following is a \eng{concluding} connector?
    \begin{enumerate}
        \item For instance
        \item In addition
        \item To sum up
    \end{enumerate}
    \item In a formal composition, contractions (e.g. \eng{don't, can't}) should be:
    \begin{enumerate}
        \item Used freely.
        \item Avoided (write \eng{do not, cannot}).
        \item Used only in the conclusion.
    \end{enumerate}
\end{enumerate}
\end{exercice}

\begin{exercice}[Vrai ou Faux ? Justifiez]
Lisez les affirmations suivantes. Écrivez \textbf{True} ou \textbf{False} et justifiez votre réponse en anglais (une phrase).
\begin{enumerate}
    \item \eng{Bandwidth} refers to the maximum data transfer rate of a network.
    \item A \eng{glitch} is a major, permanent system failure requiring hardware replacement.
    \item \eng{Cybersecurity} concerns only the protection of hardware against theft.
    \item \eng{Outdated software} often causes compatibility issues with new applications.
    \item \eng{Connectivity issues} can be caused by power cuts (\eng{power outages}) in Cameroon.
\end{enumerate}
\end{exercice}

\begin{exercice}[Vocabulaire : Complétez avec les connecteurs logiques]
Complétez le paragraphe suivant avec les connecteurs de la liste. Chaque connecteur s'utilise une seule fois.
\begin{center}
\textbf{moreover \quad however \quad as a result \quad for instance \quad to sum up}
\end{center}
\begin{textebox}[Gap-fill exercise]
Many employees struggle with digital tools at work. \_\_\_\_\_\_\_\_\_\_, lack of training is a major cause. \_\_\_\_\_\_\_\_\_\_, some companies refuse to invest in workshops. \_\_\_\_\_\_\_\_\_\_, a teacher in Buea cannot use the new platform because she never received instructions. \_\_\_\_\_\_\_\_\_\_, productivity decreases. \_\_\_\_\_\_\_\_\_\_, solving ICT difficulties requires investment and training.
\end{textebox}
\end{exercice}

\begin{exercice}[Conjugaison : Mettez le verbe à la forme convenable]
Conjuguez les verbes entre parenthèses au temps approprié.
\begin{enumerate}
    \item Look! The screen \_\_\_\_\_\_\_\_\_\_ (freeze) again. We \_\_\_\_\_\_\_\_\_\_ (lose) our data if we don't save quickly.
    \item Many schools \_\_\_\_\_\_\_\_\_\_ (not / equip) with fibre optic yet. The government \_\_\_\_\_\_\_\_\_\_ (promise) to connect them by 2025.
    \item If the internet \_\_\_\_\_\_\_\_\_\_ (be) faster, we \_\_\_\_\_\_\_\_\_\_ (can / download) the resources instantly.
    \item A new antivirus \_\_\_\_\_\_\_\_\_\_ (install) on all computers last week by the technician.
    \item Students \_\_\_\_\_\_\_\_\_\_ (must / respect) the cybersecurity rules to avoid hacking.
\end{enumerate}
\end{exercice}

\courssub{Je m'entraîne}

\begin{exercice}[Traduction : Français $\rightarrow$ Anglais]
Traduisez les phrases suivantes en anglais formel (style composition). Attention aux pièges (faux amis, ordre des mots, inombrables).
\begin{enumerate}
    \item Je ne suis pas d'accord avec l'idée que les TIC résolvent tous les problèmes.
    \item Malgré les pannes de courant fréquentes, nous devons terminer le travail.
    \item Les informations sur ce serveur ne sont pas sécurisées.
    \item Il est essentiel que les enseignants reçoivent une formation continue.
    \item À cause de la lenteur de la connexion, le téléchargement a échoué.
\end{enumerate}
\end{exercice}

\begin{exercice}[Correction d'erreurs typiques]
Les phrases suivantes contiennent chacune \textbf{une erreur} typique de francophone (grammaire, vocabulaire, orthographe, syntaxe). Soulignez l'erreur et réécrivez la phrase correcte.
\begin{enumerate}
    \item I am agree that ICTs are useful in education.
    \item The internet connection is very slow, so that we cannot work.
    \item We need more equipments to improve the network.
    \item He suggested me to buy a new router.
    \item The report has been wrote by the manager yesterday.
    \item Despite the high cost, but the software is necessary.
    \item There is a lot of informations to process.
    \item If I will have time, I will fix the bug.
    \item She is responsible to update the database.
    \item The computers are not enough powerful for this software.
\end{enumerate}
\end{exercice}

\begin{exercice}[Remise en ordre : Structure d'une composition]
Les paragraphes suivants d'une composition sur \eng{The difficulties of using ICTs in Cameroonian schools} sont mélangés. Remettez-les dans l'ordre logique (Introduction, Body 1, Body 2, Body 3, Conclusion) en numérotant de 1 à 5. Justifiez brièvement l'ordre choisi (connecteurs, progression d'idées).
\begin{enumerate}
    \item \textbf{Paragraph A:} To sum up, while ICTs offer great opportunities for learning, the challenges of connectivity, equipment and training are real. The government and school authorities must invest in infrastructure and teacher training to bridge the digital divide.
    \item \textbf{Paragraph B:} Furthermore, the lack of maintenance and technical support worsens the situation. When a computer breaks down, it often stays broken for months because there is no technician on site. Consequently, the few available devices are overused and deteriorate faster.
    \item \textbf{Paragraph C:} In today's digital age, Information and Communication Technologies have become essential tools in education. However, in many Cameroonian schools, using ICTs remains a daily struggle. This essay examines the main difficulties faced by students and teachers and suggests possible solutions.
    \item \textbf{Paragraph D:} First, poor internet connectivity is a major obstacle. Many schools, especially in rural areas, rely on weak mobile data or have no connection at all. As a result, online research and e-learning platforms are inaccessible during class hours.
    \item \textbf{Paragraph E:} Another critical issue is the insufficient and outdated equipment. Most computer labs have old machines that cannot run modern educational software. Moreover, the student-computer ratio is often 10:1, making individual practice impossible.
\end{enumerate}
\end{exercice}

\begin{exercice}[Transformation : Voix active $\rightarrow$ Voix passive]
Transformez les phrases à la voix passive. Conservez le temps verbal. N'incluez l'agent (\eng{by...}) que si nécessaire.
\begin{enumerate}
    \item Hackers steal sensitive data every day.
    \item The technician is repairing the server now.
    \item The school has bought new laptops for the lab.
    \item Power cuts often damage computer hardware.
    \item Students must submit their assignments online.
    \item The virus corrupted the database yesterday.
    \item The ministry will launch a new e-learning platform next year.
    \item We cannot solve this problem without better funding.
\end{enumerate}
\end{exercice}

\courssub{Je me prépare au Bac}

\begin{exercice}[Compréhension et langue : Texte original]
Lisez attentivement le texte ci-dessous, puis répondez aux questions.
\begin{textebox}[ICT Challenges in a Yaoundé Workplace]
At “CamTech Solutions”, a medium-sized enterprise in Yaoundé, the morning routine often starts with frustration. Employees turn on their computers, only to stare at loading screens for twenty minutes. The company relies on a single ADSL line that struggles to support fifty users. “Bandwidth is our biggest enemy,” says Marie, an accountant. “When the internet is down, I cannot access the cloud accounting software. I have to wait or use offline Excel sheets, which doubles my workload.”

Frequent power outages worsen the situation. Although the office has a generator, the switch is not automatic. Each cut causes an abrupt shutdown, risking data loss and hardware damage. Last month, a power surge destroyed the main server. The company lost three days of transactions because the last backup was a week old.

Cybersecurity is another concern. With no dedicated IT security officer, staff use weak passwords like “123456” or “CamTech2024”. Phishing emails regularly reach inboxes. Last week, an intern clicked a malicious link, installing ransomware that encrypted the HR folder. The company paid a ransom of 500,000 FCFA to recover the files.

The manager, Mr. Fono, acknowledges the problems but cites budget constraints. “We cannot afford fibre optic or a full-time IT team,” he explains. “We survive with what we have.” Yet, the hidden costs of lost time, ransoms and hardware replacement far exceed the price of a proper infrastructure. Until investment becomes a priority, CamTech Solutions will remain stuck in “survival mode”.
\end{textebox}

\textbf{Comprehension questions (answer in English, full sentences):}
\begin{enumerate}
    \item What are the three main ICT difficulties described at CamTech Solutions?
    \item Why does Marie have to use offline Excel sheets?
    \item What happened when the intern clicked the phishing link?
    \item Why does Mr. Fono hesitate to upgrade the infrastructure?
    \item According to the text, are the current “savings” actually cheap? Explain.
\end{enumerate}

\textbf{Language work:}
\begin{enumerate}
    \setcounter{enumi}{5}
    \item Find in the text a synonym for \eng{sudden and unexpected} (paragraph 2).
    \item Find in the text the opposite of \eng{strong passwords} (paragraph 3).
    \item Rewrite the sentence in the passive: \eng{A power surge destroyed the main server.}
    \item Put the verb in the correct form: \eng{If the company \_\_\_\_\_\_\_\_\_\_ (invest) in fibre optic last year, they \_\_\_\_\_\_\_\_\_\_ (not / lose) so much time now.} (Mixed conditional).
    \item Turn into reported speech: Mr. Fono said: \eng{We cannot afford a full-time IT team.}
\end{enumerate}
\end{exercice}

\begin{exercice}[Traduction : Français $\rightarrow$ Anglais]
Traduisez le paragraphe suivant en anglais. Soignez la cohérence, les connecteurs, les temps et le vocabulaire spécifique aux TIC. (Environ 80-100 mots).
\begin{quote}
Dans de nombreuses administrations camerounaises, l'utilisation des outils numériques reste un parcours du combattant. Tout d'abord, la connexion internet est souvent instable, ce qui rend l'accès aux serveurs distants quasi impossible. De plus, les ordinateurs sont vétustes et tombent fréquemment en panne. Par conséquent, les agents perdent un temps précieux à attendre les réparations. En outre, le manque de formation empêche le personnel d'exploiter les logiciels disponibles. Pour remédier à cela, l'État devrait investir dans la fibre optique, renouveler le parc informatique et organiser des sessions de formation continue. En définitive, moderniser l'administration passe par une maîtrise effective des TIC.
\end{quote}
\end{exercice}

\begin{exercice}[Sujet de composition type Baccalauréat]
\textbf{Consigne :} Traitez le sujet suivant en \textbf{250 à 300 mots} environ. Respectez la structure \textit{Introduction / Development (2-3 paragraphs) / Conclusion}. Utilisez des connecteurs logiques variés. Relisez-vous pour corriger les fautes (accords, temps, orthographe, inombrables).

\begin{center}
\textbf{Write a composition on the following topic:}
\end{center}
\begin{quote}
\eng{Write a composition of 250-300 words on the difficulties of using ICTs in your school and suggest solutions.}
\end{quote}

\textbf{Indications pour la rédaction :}
\begin{itemize}
    \item \textbf{Introduction :} Hook (actualité numérique), reformulation du sujet, \textbf{thesis statement} (annonce du plan : difficultés + solutions).
    \item \textbf{Développement :} 2 ou 3 paragraphes distincts. Ex: 1. Infrastructure (connexion, matériel, électricité). 2. Facteurs humains (formation, maintenance, cybersécurité). 3. Solutions concrètes (investissement, partenariats, formation, maintenance préventive). Utilisez \eng{First, Furthermore, Moreover, However, Consequently, For instance}.
    \item \textbf{Conclusion :} Résumé des points principaux, ouverture/recommandation finale (\eng{To sum up, In conclusion, Ultimately}).
\end{itemize}

\textbf{Auto-évaluation / Évaluation par les pairs (Peer Assessment) :}
Après rédaction, échangez votre copie avec un camarade. Notez sur 20 en utilisant la grille ci-dessous (4 points par critère). Rédigez un commentaire constructif en anglais (2-3 phrases).

\begin{center}
\begin{tabularx}{\linewidth}{|X|c|c|c|c|c|}
\hline
\rowcolor{popblueL}
\textbf{Critère} & \textbf{4} & \textbf{3} & \textbf{2} & \textbf{1} & \textbf{0} \\
\hline
\textbf{Task Achievement} (Réponse au sujet, nombre de mots, solutions proposées) & & & & & \\
\hline
\textbf{Organisation} (Paragraphes, connecteurs, progression logique) & & & & & \\
\hline
\textbf{Grammar \& Syntax} (Temps, voix passive, modaux, phrases complexes) & & & & & \\
\hline
\textbf{Vocabulary} (Lexique TIC, variété, précision, pas de franglais) & & & & & \\
\hline
\textbf{Mechanics} (Orthographe, ponctuation, majuscules, pas de contractions) & & & & & \\
\hline
\end{tabularx}
\end{center}

\textbf{Commentaire du pair (en anglais) :} \eng{Your essay correctly identifies... However, you should work on... Good use of...}
\end{exercice}

\newpage

\coursec{Corrigés des exercices}

\begin{corrige}[Exercice 1 — QCM : Structure et connecteurs]
\begin{enumerate}
    \item \textbf{b)} At the end of the introduction. La thèse (\eng{thesis statement}) clôt l'introduction après l'accroche et le contexte ; elle annonce le plan.
    \item \textbf{b)} \eng{However} exprime le contraste. \eng{Furthermore} = addition ; \eng{Consequently} = conséquence.
    \item \textbf{a)} \eng{Present Simple} : il exprime les faits généraux et les vérités permanentes (\eng{Slow internet affects productivity}).
    \item \textbf{c)} \eng{To sum up} introduit la conclusion. \eng{For instance} = exemple ; \eng{In addition} = addition.
    \item \textbf{b)} Dans une composition formelle, on évite les contractions : on écrit \eng{do not}, \eng{cannot}, \eng{it is}.
\end{enumerate}
\end{corrige}

\begin{corrige}[Exercice 2 — Vrai ou Faux ? Justifiez]
\begin{enumerate}
    \item \textbf{True.} \eng{Bandwidth refers to the maximum amount of data that can be transferred over a network per second.} (La bande passante désigne le débit maximal d'un réseau.)
    \item \textbf{False.} \eng{A glitch is only a minor, temporary malfunction, not a permanent failure.} (Un \eng{glitch} est un dysfonctionnement mineur et temporaire.)
    \item \textbf{False.} \eng{Cybersecurity protects data, software and networks against attacks, not only hardware against theft.} (La cybersécurité protège aussi les données et les logiciels.)
    \item \textbf{True.} \eng{Outdated software often cannot run or communicate with new applications, which causes compatibility issues.}
    \item \textbf{True.} \eng{Power outages shut down routers and servers, so they directly cause connectivity issues in Cameroon.}
\end{enumerate}
\end{corrige}

\begin{corrige}[Exercice 3 — Vocabulaire : Complétez avec les connecteurs logiques]
\eng{Many employees struggle with digital tools at work. \textbf{Moreover}, lack of training is a major cause. \textbf{However}, some companies refuse to invest in workshops. \textbf{For instance}, a teacher in Buea cannot use the new platform because she never received instructions. \textbf{As a result}, productivity decreases. \textbf{To sum up}, solving ICT difficulties requires investment and training.}

\textbf{Justifications :}
\begin{itemize}
    \item \eng{Moreover} : addition d'une cause (le manque de formation s'ajoute à la difficulté).
    \item \eng{However} : contraste (les entreprises refusent d'investir malgré le besoin).
    \item \eng{For instance} : illustration par un exemple concret.
    \item \eng{As a result} : conséquence (la productivité diminue).
    \item \eng{To sum up} : conclusion du paragraphe.
\end{itemize}
\end{corrige}

\begin{corrige}[Exercice 4 — Conjugaison : Mettez le verbe à la forme convenable]
\begin{enumerate}
    \item \eng{Look! The screen \textbf{is freezing} again. We \textbf{will lose} our data if we do not save quickly.} — \eng{Look!} appelle le Present Continuous (action en cours) ; 1\textsuperscript{er} conditionnel : \eng{if} + Present Simple $\rightarrow$ \eng{will} + base verbale. (Registre formel : \eng{do not} plutôt que \eng{don't}.)
    \item \eng{Many schools \textbf{have not been equipped} with fibre optic yet. The government \textbf{has promised} to connect them by 2025.} — Present Perfect passif (\eng{yet} = marqueur ; focus sur le résultat, agent secondaire) ; Present Perfect pour une promesse faite avec conséquence présente.
    \item \eng{If the internet \textbf{were} faster, we \textbf{could download} the resources instantly.} — 2\textsuperscript{e} conditionnel (hypothèse irréelle au présent) : \eng{if} + Prétérit (\eng{were} pour toutes les personnes au style soutenu) $\rightarrow$ \eng{could} + base verbale.
    \item \eng{A new antivirus \textbf{was installed} on all computers last week by the technician.} — Prétérit passif (\eng{last week} ; sujet singulier $\rightarrow$ \eng{was} + participe passé).
    \item \eng{Students \textbf{must respect} the cybersecurity rules to avoid hacking.} — Modal \eng{must} + base verbale (obligation forte).
\end{enumerate}
\end{corrige}

\begin{corrige}[Exercice 5 — Traduction : Français $\rightarrow$ Anglais]
\begin{enumerate}
    \item \eng{I do not agree with the idea that ICTs solve all problems.} — Piège : jamais \eng{I am agree} ; \eng{agree with}.
    \item \eng{Despite frequent power cuts, we must finish the work.} — \eng{Despite} + nom ; on peut aussi écrire \eng{Despite the fact that power cuts are frequent, ...}.
    \item \eng{The information on this server is not secure.} — \eng{Information} est inombrable (pas de \eng{s}) et le verbe reste au singulier ; \eng{secure} (adjectif), pas \eng{secured} ici.
    \item \eng{It is essential that teachers receive continuous training.} — Subjonctif mandatif : base verbale \eng{receive} (pas de \eng{s} à la 3\textsuperscript{e} personne) ; variante : \eng{should receive}.
    \item \eng{Because of the slow connection, the download failed.} — \eng{Because of} + nom ; variante : \eng{Due to the slow connection, ...}.
\end{enumerate}
\end{corrige}

\begin{corrige}[Exercice 6 — Correction d'erreurs typiques]
\begin{enumerate}
    \item Erreur : \eng{am agree}. \eng{I \textbf{agree} that ICTs are useful in education.} (\eng{agree} est un verbe.)
    \item Erreur : \eng{so that}. \eng{The internet connection is very slow, \textbf{so} we cannot work.} (\eng{so} = conséquence ; \eng{so that} = but.)
    \item Erreur : \eng{equipments}. \eng{We need more \textbf{equipment} to improve the network.} (inombrable.)
    \item Erreur : \eng{suggested me to buy}. \eng{He suggested \textbf{buying} a new router.} ou \eng{He suggested \textbf{that I buy} a new router.}
    \item Erreur : \eng{wrote}. \eng{The report has been \textbf{written} by the manager.} (participe passé irrégulier \eng{write–wrote–written} ; on supprime aussi \eng{yesterday}, incompatible avec le Present Perfect, ou on passe au Prétérit : \eng{was written yesterday}.)
    \item Erreur : \eng{Despite ... but}. \eng{\textbf{Despite} the high cost, the software is necessary.} ou \eng{\textbf{Although} the cost is high, the software is necessary.} (jamais deux marqueurs de concession ensemble.)
    \item Erreur : \eng{informations}. \eng{There is a lot of \textbf{information} to process.} (inombrable.)
    \item Erreur : \eng{will have}. \eng{If I \textbf{have} time, I will fix the bug.} (Present Simple après \eng{if}.)
    \item Erreur : \eng{responsible to update}. \eng{She is responsible \textbf{for updating} the database.} (\eng{responsible for} + V-ing.)
    \item Erreur : \eng{enough powerful}. \eng{The computers are not \textbf{powerful enough} for this software.} (\eng{enough} se place après l'adjectif.)
\end{enumerate}
\end{corrige}

\begin{corrige}[Exercice 7 — Remise en ordre : Structure d'une composition]
Ordre correct : \textbf{1 = C, 2 = D, 3 = E, 4 = B, 5 = A}.

\textbf{Justification :}
\begin{itemize}
    \item \textbf{C (Introduction)} : accroche (\eng{In today's digital age}), contraste (\eng{However}) et thesis statement (\eng{This essay examines... and suggests...}).
    \item \textbf{D (Body 1)} : connecteur de séquence initial \eng{First} ; premier problème : la connectivité.
    \item \textbf{E (Body 2)} : \eng{Another critical issue} enchaîne sur un deuxième problème : le matériel obsolète.
    \item \textbf{B (Body 3)} : \eng{Furthermore} ajoute un troisième problème (maintenance) ; \eng{Consequently} en tire la conséquence.
    \item \textbf{A (Conclusion)} : \eng{To sum up} résume les trois difficultés et formule la recommandation finale.
\end{itemize}
La progression est logique : problèmes (connexion $\rightarrow$ équipement $\rightarrow$ maintenance) puis synthèse et solutions en conclusion.
\end{corrige}

\begin{corrige}[Exercice 8 — Transformation : Voix active $\rightarrow$ Voix passive]
\begin{enumerate}
    \item \eng{Sensitive data \textbf{is stolen} every day.} (Present Simple passif ; agent \eng{by hackers} omis car évident.)
    \item \eng{The server \textbf{is being repaired} now.} (Present Continuous passif : \eng{is being} + participe passé.)
    \item \eng{New laptops \textbf{have been bought} for the lab.} (Present Perfect passif.)
    \item \eng{Computer hardware \textbf{is often damaged} by power cuts.} (agent utile : c'est la cause.)
    \item \eng{Assignments \textbf{must be submitted} online.} (modal + \eng{be} + participe passé.)
    \item \eng{The database \textbf{was corrupted} yesterday.} (Prétérit passif.)
    \item \eng{A new e-learning platform \textbf{will be launched} next year.} (futur passif : \eng{will be} + participe passé.)
    \item \eng{This problem \textbf{cannot be solved} without better funding.} (modal négatif + \eng{be} + participe passé.)
\end{enumerate}
\textbf{Rappel de méthode :} l'objet de la phrase active devient le sujet ; on conjugue \eng{be} au temps du verbe actif, puis on ajoute le participe passé.
\end{corrige}

\begin{corrige}[Exercice 9 — Compréhension et langue : Texte original]
\textbf{Comprehension questions :}
\begin{enumerate}
    \item \eng{The three main difficulties are poor bandwidth on a single ADSL line, frequent power outages that damage hardware and cause data loss, and weak cybersecurity (weak passwords, phishing and ransomware).}
    \item \eng{Marie has to use offline Excel sheets because when the internet is down, she cannot access the cloud accounting software, which doubles her workload.}
    \item \eng{When the intern clicked the phishing link, ransomware was installed and encrypted the HR folder, and the company had to pay a ransom of 500,000 FCFA to recover the files.}
    \item \eng{Mr. Fono hesitates to upgrade the infrastructure because of budget constraints; he claims the company cannot afford fibre optic or a full-time IT team.}
    \item \eng{No, the current “savings” are not actually cheap, because the hidden costs of lost time, ransoms and hardware replacement far exceed the price of a proper infrastructure.}
\end{enumerate}

\textbf{Language work :}
\begin{enumerate}
    \item \eng{abrupt} (dans \eng{an abrupt shutdown}, paragraphe 2).
    \item \eng{weak passwords} (paragraphe 3).
    \item \eng{The main server was destroyed by a power surge.} (Prétérit passif : \eng{was} + participe passé \eng{destroyed}.)
    \item \eng{If the company \textbf{had invested} in fibre optic last year, they \textbf{would not lose} (ou \textbf{would not be losing}) so much time now.} — Conditionnel mixte : \eng{if} + Past Perfect (cause passée) $\rightarrow$ \eng{would} + base verbale (conséquence présente).
    \item \eng{Mr. Fono said (that) they could not afford a full-time IT team.} — Discours rapporté au passé : \eng{cannot} $\rightarrow$ \eng{could not} ; \eng{we} $\rightarrow$ \eng{they}.
\end{enumerate}
\end{corrige}

\begin{corrige}[Exercice 10 — Traduction : Français $\rightarrow$ Anglais]
\textbf{Proposition de traduction (environ 95 mots) :}

\eng{In many Cameroonian administrations, using digital tools remains an uphill struggle. First, the internet connection is often unstable, which makes access to remote servers almost impossible. Moreover, the computers are outdated and frequently break down. As a result, civil servants waste precious time waiting for repairs. In addition, the lack of training prevents staff from making full use of the available software. To address these problems, the State should invest in fibre optic, renew the computer fleet and organise continuous training sessions. Ultimately, modernising the administration requires a genuine command of ICTs.}

\textbf{Points de vigilance :}
\begin{itemize}
    \item \eng{an uphill struggle} / \eng{an obstacle course} pour « un parcours du combattant » (pas de traduction littérale).
    \item « vétustes » = \eng{outdated} (pas \eng{old-fashioned} qui connote la mode).
    \item « les agents » = \eng{civil servants} / \eng{employees} (pas \eng{agents}, faux ami : \eng{agent} = agent secret/commercial).
    \item « renouveler le parc informatique » = \eng{renew the computer fleet} / \eng{replace the computer stock}.
    \item Connecteurs : \eng{First, Moreover, As a result, In addition, Ultimately} — un par fonction.
\end{itemize}
\end{corrige}

\begin{corrige}[Exercice 11 — Sujet de composition type Baccalauréat]
\textbf{Barème indicatif (sur 20) :} Task Achievement /4 ; Organisation et connecteurs /4 ; Grammaire et syntaxe /4 ; Vocabulaire /4 ; Mécanique (orthographe, ponctuation) /4. Pénaliser : hors-sujet, moins de 200 mots, contractions systématiques, franglais.

\textbf{Proposition de rédaction modèle (environ 280 mots) :}

\eng{\textbf{Overcoming ICT Difficulties in My School}}

\eng{In today's digital age, Information and Communication Technologies have become essential for quality education. However, in my school, using ICTs remains a daily challenge rather than an opportunity. This essay examines the main infrastructure and human difficulties we face and suggests realistic solutions.}

\eng{First, poor connectivity is the most visible obstacle. Our school relies on a weak mobile data connection, and the network collapses whenever more than twenty students log in. For instance, during computer science practicals, a single webpage can take five minutes to load. Consequently, online research and e-learning platforms are almost unusable. Moreover, the equipment is outdated: most machines in the computer lab are over ten years old, and the student-computer ratio is nearly 8:1, which makes individual practice impossible. Frequent power cuts make matters worse, as the school has no generator or UPS system.}

\eng{Furthermore, the human factor deepens these technical problems. Many teachers have never received proper training in digital pedagogy, so they lack the confidence to use the available tools. In addition, there is no IT technician on site; when a computer breaks down, it may stay unrepaired for a whole term.}

\eng{To address these issues, several measures should be taken. The government must invest in fibre-optic connections and solar power for schools. Public-private partnerships could also help renew the computer fleet at lower cost. Finally, it is essential that every teacher receive continuous training and that a maintenance officer be appointed for each school cluster.}

\eng{To sum up, although ICTs offer great opportunities, poor infrastructure, outdated equipment and insufficient training hold my school back. Ultimately, only strong political will and sustained investment can bridge the digital divide and allow students to benefit fully from the digital revolution.}

\textbf{Points de vigilance pour la correction des copies :}
\begin{itemize}
    \item Vérifier la présence d'un \textbf{thesis statement} en fin d'introduction et d'un \textbf{topic sentence} par paragraphe.
    \item Exiger au moins 5 connecteurs variés et correctement ponctués (\eng{However,} avec virgule ; pas de \eng{Despite... but}).
    \item Sanctionner les inombrables (\eng{equipments, informations}), les contractions et les faux amis (\eng{actually, eventually, sensible}).
    \item Valoriser le passif, les conditionnels et le subjonctif (\eng{It is essential that every teacher receive...}).
\end{itemize}
\end{corrige}

\newpage

\coursec{Fiche bilan}

\courssub{Fiche Bilan -- Lesson 60 : Writing Compositions on ICT Difficulties}

\begin{popfigure}
\begin{tikzpicture}[
  mindmap,
  grow cyclic,
  every node/.style={concept, execute at begin node=\hskip0pt},
  concept color=popblue,
  level 1/.append style={level distance=4.5cm, sibling angle=360/6},
  level 2/.append style={level distance=3cm, sibling angle=360/4},
  branch/.style={concept color=#1, edge from parent/.style={draw=#1, thick}},
  node font=\small\bfseries,
  text width=2.8cm,
  align=center
]
\node[concept, concept color=popdark, align=center]{ICT Composition\\Difficulties at Work/\& Learning}
  child[branch=popblue] {node[concept, align=center]{1. Task Analysis\\ \& Brainstorming}}
  child[branch=poporange] {node[concept, align=center]{2. Structure\\Intro / Body / Conclusion}}
  child[branch=popgreen] {node[concept, align=center]{3. Vocabulary\\ \& Connectors}}
  child[branch=poppurple] {node[concept, align=center]{4. Grammar\\ \& Verb Tenses}}
  child[branch=poppink] {node[concept, align=center]{5. Model Essay\\ Commented}}
  child[branch=popgold] {node[concept, align=center]{6. Checklist\\ \& Assessment}};
\end{tikzpicture}
\legende{Carte mentale de la leçon : les 6 étapes clés pour réussir une composition sur les difficultés des TIC.}
\end{popfigure}

\begin{bilan}

\textbf{1. Lexique clé (Vocabulary) -- Thème : ICT Difficulties}

\begin{center}
\begin{tabularx}{\linewidth}{|X|X|X|}
\hline
\rowcolor{popblueL}
\textbf{English} & \textbf{Français} & \textbf{Contexte / Collocation} \\
\hline
\cle{connectivity issues} & problèmes de connexion & \eng{Poor connectivity hinders remote work.} \\
\cle{power outage} / \cle{blackout} & coupure d'électricité & \eng{Frequent power outages damage hardware.} \\
\cle{digital divide} & fracture numérique & \eng{The digital divide affects rural students.} \\
\cle{cybersecurity threat} & menace de cybersécurité & \eng{Phishing is a major cybersecurity threat.} \\
\cle{data breach} & violation de données & \eng{A data breach compromises privacy.} \\
\cle{technical glitch} & bug / incident technique & \eng{A glitch froze the video conference.} \\
\cle{digital literacy} & culture / compétences numériques & \cle{Lack of digital literacy} slows adoption. \\
\cle{bandwidth} & bande passante & \eng{Low bandwidth causes lag.} \\
\cle{hardware failure} & panne matérielle & \eng{Old computers suffer hardware failure.} \\
\cle{software compatibility} & compatibilité logicielle & \eng{Compatibility issues waste time.} \\
\cle{distraction} / \cle{procrastination} & distraction / procrastination & \eng{Social media leads to procrastination.} \\
\cle{eye strain} / \cle{burnout} & fatigue visuelle / épuisement & \eng{Screen time causes eye strain.} \\
\hline
\end{tabularx}
\end{center}

\vspace{0.5em}
\textbf{2. Connecteurs logiques (Linkers) -- Pour structurer la composition}

\begin{center}
\begin{tabularx}{\linewidth}{|X|X|X|}
\hline
\rowcolor{poporangeL}
\textbf{Fonction} & \textbf{Connecteurs (Formal)} & \textbf{Exemple court} \\
\hline
Addition & \eng{Furthermore, Moreover, In addition, Besides} & \eng{Moreover, costs are high.} \\
Cause & \eng{Due to, Owing to, Because of, Since} & \eng{Due to outages, work stops.} \\
Consequence & \eng{Consequently, Therefore, As a result, Thus} & \eng{Thus, productivity drops.} \\
Contrast & \eng{However, Nevertheless, On the one hand... on the other hand, Whereas} & \eng{However, ICTs offer opportunities.} \\
Example & \eng{For instance, Such as, Namely} & \eng{For instance, in Buea...} \\
Conclusion & \eng{To conclude, In conclusion, Ultimately, All things considered} & \eng{To conclude, training is vital.} \\
\hline
\end{tabularx}
\end{center}

\vspace{0.5em}
\textbf{3. Règles de grammaire essentielles (Grammar Toolbox)}

\begin{center}
\begin{tabularx}{\linewidth}{|X|X|X|}
\hline
\rowcolor{popgreenL}
\textbf{Point} & \textbf{Règle / Forme} & \textbf{Exemple traduit} \\
\hline
\cle{Present Perfect} & \eng{have/has + V-ed/Participe} -- Expérience, conséquence présente & \eng{Many workers \textbf{have faced} connectivity issues.} (Ont rencontré) \\
\cle{Present Simple} & Faits généraux, vérités permanentes & \eng{Power cuts \textbf{occur} daily in Douala.} (Surviennent) \\
\cle{Modals (Should, Must, Could)} & Recommandation, obligation, possibilité & \eng{Employers \textbf{should provide} training.} (Devraient fournir) \\
\cle{Passive Voice} & \eng{be + V-ed} -- Focus sur l'objet/action & \eng{Data \textbf{is stolen} by hackers.} (Est volée) \\
\cle{Relative Clauses} & \eng{who/which/where/whose} -- Préciser le nom & \eng{Students \textbf{who lack} laptops struggle.} (Qui n'ont pas) \\
\cle{Conditionnel (Type 1/2)} & \eng{If + Present, Will} / \eng{If + Preterite, Would} & \eng{If we \textbf{invested}, results \textbf{would improve}}.\ (Si nous investissions...) \\
\cle{Gerund vs Infinitive} & Après préposition $\to$ \cle{-ing} ; Après certain verbes $\to$ \cle{to + V} & \eng{Good at \textbf{solving} problems.} / \eng{Need \textbf{to upgrade} systems.} \\
\hline
\end{tabularx}
\end{center}

\vspace{0.5em}
\textbf{4. Structure type d'une composition (Essay Structure)}

\end{bilan}

\begin{methode}[Plan en 4 paragraphes]
\begin{enumerate}
\item \textbf{Introduction (1 paragraphe)} : Accroche (Hook) $\to$ Définition du contexte (ICT au travail/école au Cameroun) $\to$ \cle{Thesis statement} (annonce du plan : difficultés + solutions).
\item \textbf{Development -- Paragraph 1 (Difficulties)} : Phrase sujet + 2-3 arguments (ex: infrastructure, compétences, sécurité) + connecteurs (\eng{Firstly, Furthermore}) + exemples locaux (Eneo, MTN, Orange, zones rurales).
\item \textbf{Development -- Paragraph 2 (Solutions / Counter-arguments)} : Phrase sujet + mesures (formation, investissement, politique publique) + connecteurs (\eng{However, On the other hand, To remedy this}).
\item \textbf{Conclusion (1 paragraphe)} : Reformulation de la thèse + Synthèse des points forts + Ouverture / Recommandation finale (\eng{All things considered, the government must...}).
\end{enumerate}
\end{methode}

\vspace{0.5em}
\textbf{5. Méthode express : Rédaction en 4 étapes (Exam conditions)}

\begin{enumerate}
\item \textbf{Analyse (5 min)} : Souligner mots-clés (\eng{difficulties, work, learning, ICTs}). Cerveau en ébullition (mindmap) : 3 difficultés + 2 solutions.
\item \textbf{Plan détaillé (5 min)} : Rédiger la \cle{Thesis statement} + mots-clés par paragraphe + connecteurs choisis.
\item \textbf{Rédaction (35 min)} : Écrire lisiblement, 250-350 mots. Utiliser l'anglais standard (pas de SMS, pas de pidgin). Varier les structures (passif, modaux, relatives).
\item \textbf{Relecture (10 min)} : Checklist (voir boîte ci-dessous). Vérifier : \cle{Subject-Verb Agreement}, \cle{Tenses}, \cle{Prepositions}, \cle{Spelling}, \cle{Punctuation}, \cle{Connectors}.
\end{enumerate}



\begin{aretenir}[Checklist avant le Bac -- Composition ICT Difficulties]
$\square$ \textbf{Task Achievement} : Le sujet est-il entièrement traité ? (Difficultés travail \textbf{ET} apprentissage).
$\square$ \textbf{Structure} : Introduction (Hook + Thesis) -- Body 1 (Difficultés) -- Body 2 (Solutions/Nuances) -- Conclusion.
$\square$ \textbf{Coherence/Cohesion} : Paragraphes distincts ? Connecteurs variés (\eng{Moreover, However, Consequently}) utilisés correctement ?
$\square$ \textbf{Vocabulary Range} : Lexique spécifique TIC utilisé (\eng{connectivity, bandwidth, digital divide, cybersecurity, literacy}) ? Pas de répétitions excessives.
$\square$ \textbf{Grammar Range} : Présent Perfect, Passif, Conditionnels, Relatives, Modaux présents sans erreur ?
$\square$ \textbf{Subject-Verb Agreement} : \eng{The government \textbf{is}...}, \eng{Workers \textbf{face}...}, \eng{Data \textbf{is}...} (pas \eng{are}).
$\square$ \textbf{Prepositions} : \eng{Good \textbf{at} using}, \eng{Depend \textbf{on} electricity}, \eng{Access \textbf{to} internet}, \eng{Interested \textbf{in} learning}.
$\square$ \textbf{Spelling/Capitalization} : \eng{ICTs}, \eng{Internet}, \eng{Cameroon}, \eng{Yaoundé}, \eng{Buea} (majuscules). Mots difficiles : \eng{necessary, occurrence, maintenance, environment}.
$\square$ \textbf{Punctuation} : Virgules après connecteurs initiaux (\eng{However,}), points finaux, pas de phrases "run-on".
$\square$ \textbf{Word Count} : Entre 250 et 350 mots (idéalement $\approx$ 300).
$\square$ \textbf{Tone/Register} : Formel, objectif, pas de contractions (\eng{do not} au lieu de \eng{don't}), pas de langage familier.
$\square$ \textbf{Local Context} : Au moins un exemple camerounais concret cité (Eneo, zones rurales, bendskin, marché Mokolo, université de Yaoundé I).
\end{aretenir}

\findecours

\end{document}
